
%%%%%%%%%%%%%%%%%%%%%%%%%%%%%%%% SCRIPT FOR E-RARA AND MDZ FILES     %%%%%%%%%%%%%%%%%%%%%%%%%%%%%%%%%%%%%%%%%%%%%%%%
%%%%%%%%%%%%%%%%%%%%%%%%%%% fini le 30.04.2025 par F. GOY            %%%%%%%%%%%%%%%%%%%%%%%%%%%%%%%%%%%%%%%%%%%%%%%%
% !TeX TS-program = lualatex
\documentclass{article}
\usepackage[T1]{fontenc}
\usepackage{microtype}
\usepackage[pdfusetitle,hidelinks]{hyperref}

\usepackage{polyglossia}
\setmainlanguage{english}
\setotherlanguages{latin,greek}
\usepackage[series={},nocritical,noend,noeledsec,nofamiliar,noledgroup]{reledmac}
\usepackage{reledpar}

\usepackage{fontspec}
\setmainfont{TeX Gyre Termes}

\usepackage{sectsty}
\usepackage{xcolor}

\usepackage{fancyhdr}
\pagestyle{fancy}
\fancyhf{}
\fancyhead[LE,RO]{\nouppercase{\leftmark}}  
\cfoot{\thepage}
\renewcommand{\headrulewidth}{0.4pt}

% Redefine \section to remove numbering
\usepackage{titlesec}
\titleformat{\section}[block]{\normalfont\scshape\color{gray}}{}{0pt}{} % no number in heading
\titleformat{\subsection}[hang]{\normalfont}{}{0pt}{} % also remove subsection number
\titleformat{\subsubsection}[hang]{\normalfont\footnotesize\color{black}}{}{0pt}{}

% Modify how section marks are stored to exclude numbers
\makeatletter
\renewcommand{\sectionmark}[1]{%
	\markboth{#1}{}} % Only store the section title, without number
\renewcommand{\subsectionmark}[1]{%
	\markright{#1}} % Only store the subsection title, without number
\renewcommand{\numberline}[1]{} % Hide the section number in TOC
\makeatother

\begin{document}

\date{}
        \title{Commentarium in epistolam D Pauli Apostoli ad Timotheum primam : [Pellicanus Conrad], [1539]}
\maketitle
\tableofcontents
\clearpage
\begin{pages} 
\beginnumbering
        
\section*{AD THESSAL. CAP, III. }
\marginpar{[ p.463 ]}\pstart Salutatio mea manu Pauli, quod est signum in omni epistola. lta scribo. Gratia domini nostri lesu Christi cum omnibus uobis. Amen.  \pend\pstart Hoc symbolum obseruabitis in omnibus meis epistolis, siue ad uos, siue ad alios scriptis. Sic enim scribo, ne per subdititias epistolas quisquam uobis imponat. Haec autem ultima salutatio manu Apostoli conseripta est, qua gratiam domini nostri lesu Christi optatesse cum eis. Amen.  \pend
\phantomsection
\addcontentsline{toc}{subsection}{\textit{COMMENTARIVM IN EPISTOLAM D. PAVLI APOSTOLI AD TIMOTHEVM PRIMAM. }}
\subsection*{\textit{COMMENTARIVM IN EPISTOLAM D. PAVLI APOSTOLI AD TIMOTHEVM PRIMAM. }}\pstart FINIS.  \pend
\phantomsection
\addcontentsline{toc}{subsection}{\textit{ARGVMENTVM PER ERAS. ROTEROD. }}
\subsection*{\textit{ARGVMENTVM PER ERAS. ROTEROD. }}\pstart \huge\textbf{T}\normalsize IMOTHEVM matre ludaea, sed Christiana, patre Graeco, Paulus in ministerium adoptarat, probae indolis iuuenem, et sacris literis eruditum, quem tamen ob ludaeos coactus est circuncidere. Quoniam autem huic ecclesiarum curam delegarat sicut et Tito, quas ipse non poterat adire, instituit eum in functione episcopali, et in disciplina ecclesiae: non sic ad monens ut discipulum, sed ut filium et collegam. Id quo cum maiore faciat authori tate, subinde apostolicam authoritatem sibi uendicat. Admonet autem, ut reiectis qui ludaicas fabulas inueherent, ea doceret quae essent fidei et charitatis. Deinde principum, licet ethnicorum, authoritatem, quomam hinc pendebat ordo ciuita tis, et tranquillitas reip. adeo non spernendam Christianis, ut pro his orari iubeat. Quid uiros deceat in coetu ecclesiastico, quid mulieres, praescribit. Episcopum de pingit cum sua familia. Atque haec ferme tribus primis capitibus. Deinde admonet ne recipiat Iudaicas fabulas, de ciborum delectu, de uitando coniugio. Mox docet qualem se praestare debeat erga grandes natu uiros, erga iuuenes, erga anus, erga puellas, erga uiduas diuites ac pauperes, et ecclesiae stipe alendas, erga adoIescentiores et suspectae etiamnum aetatis. Ad haec praescribit quid praecipere debeat dominis, quid seruis, quid diuitibus. Admonens modis omnibus reiiciendas contentiosas quaestiones sophistarum, inani doctrinae specie sese uenditantium. Scribite Laodicea per Tychicum diaconum.  \pend
\phantomsection
\addcontentsline{toc}{subsection}{\textit{AD LECTOREM C.P. }}
\subsection*{\textit{AD LECTOREM C.P. }}\pstart COMMENTARIA in hanc epistolam primam, lucubrationum mearum primitiae sunt, quas tentaui uerius olim quam scripsi, Anno D. M. X XII. Basileae, postea autem impressae ibidem sunt me ignorante. Nunc uero reuisae et correctae, suum quoque in ordinem repositae sunt. In quibus tempori tum seruiens, phra si quoque alia me usum Lector non miraberis. Sed saluo per omnia iudicio tuo, bo ni consulas, quae conformia scripturae, uel paululum infirmiora, offenderis.  \pend
\section*{COMMENT. IN I. EPIST. }
\marginpar{[ p.464 ]}
\section{CAPVT PRIMVM.}
\phantomsection
\addcontentsline{toc}{subsection}{\textit{\huge\textbf{P}\normalsize AVLVS apostolus lesu Christi, secundum imperium dei saluatoris nostri et Christi lesu spei nostrae Timotheo dilecto filio in fide. Gratia et misericordia et pax a deo patre et Christo lesu domino nostro. }}
\subsection*{\textit{\huge\textbf{P}\normalsize AVLVS apostolus lesu Christi, secundum imperium dei saluatoris nostri et Christi lesu spei nostrae Timotheo dilecto filio in fide. Gratia et misericordia et pax a deo patre et Christo lesu domino nostro. }}\pstart Super mutatione nominis Pauli ex Saule, plures multa ac uaria tradiderunt. Simplicior ratio ea uidetur, quae et sacris probatur testimoniis. Quia intra tribum Beniamin, unde et natum se Paulus fatetur, plausibile nomen erat Sauli, siquidem regale et celeberrimum, ob regem primum electum a domino ex filiis Israel, quo non fuit secundum saeculum dexterior, sed et sanctior, praeter inuidiam, qua sancto Dauidi usque ad mortem aduersatus est, et inobedientiam inde profectam, con tra ordinationem diuinam, de regno Dauidis, cui totis uiribus reluctabatur. Cum itaque penes Iudaeos de Iudaismo conuersos ad Christum, inuidiam annexam haberet hoc Saulis nomen, ob similia studia eiusdem uiri contra Christum et suos discipulos: sic ut Christus quoque eundem insimulare uoluerit Saulini, ut ita dicatur, ingenii, dicens, Saule Saule, ad quid me, nempe Messiam Dauidis filium persequeris? Vnde et Lucas Saulum tunc tantum nominat, cum persecutiones Christo Dauidi il latas a Paulo commemorat. Ideo ne Christi persecutoris nomine, sicut re et opere, notorius habitus fuerat, perpetuo uocaretur, Sicut Dauidem regem sacrum et electum perpetuo persecutus est priscus Saulus, Saulus suum nomen mutauit, et humilius delegit gentile Romanumque Pauli nomen. Si Latine interpreteris, uel operatorium et sedulum in sancto opere: si Hebraice, mirabilis. Qui certe non ni si de suis consueuit laboribus gloriari. Nostri pontifices et electi apostolici illi, hu mile nomen plaerunque ponunt, et arrogant sibi imperialia nomina, ut ex loanne Leo, ex Theodorico Alexander, ex Iuliano lulius, uastor ex uastatore. Non spernenda tamen interim cura sit nominum et cognominum, ne uel ridicula uel obscoena infantibus obtrudantur, unde uiri aliquando erubescant, et posteri, maxi me doctores: honesta et ad pietatem ac uirtutes commonitoria nomina habere, non fuerit inutile. Quod et ethnicis olim curae fuisse comperimus, et sanctis patriarchis: imo priscis nostris Germanis, Graecis quoque et Latinis in prisca religione. Apostolus.) A Christo iam nomen acceperat uas electionis: erat re uera doctor gentium in fide et ueritate, erat ecclesiarum pastor, et plantator, sed se missum et officio delegatum fatetur, ne quid altius uideri uelituel nominari: ne honorem quaerere, sed legationis officium agere se testetur, plenum laboribus, curis, periculis. Non otiosum, gloriosum, saeculare, quale nunc sectantur multi, qui aposto lici uocari diligunt, re uera regibus mundaniores, et a Christo alienissima conuer satione, apostatici potius iudicandi quam̃ apostolici. Iesu Christi.) Cuius nunc gloriosum ac saluificum nomen subinde Paulus celebrat in omnibus suis epistolis, cuius nuper inuocatores, Saulus iam ante, ad mortem crudeliter fuerat persecutus. Nunc portat ac profert coram gentibus, regibus, et filiis Israel salutare nomen ipsum, in quo solo seruamur fideles, nomen nouum, quod os domini nominauit. Saluator est enim lesus, omnium sibi credentium. Qui ut priscus ille losue, filius Nun, inducit ad promissionis terram, quod non poterat Moses. Qui et de Babylonica liberat captiuitate, ut lesus filius Iosedech, summus pontifex semper oran do pro populo dei. Qui et unctus est, rex regum, et dominus dominantium. Sacerdos quoque magnus semper assistens pro nobis, qui plenus gratia et ueritate, filios habet omnes unctos, a quo nomen nostrum consequimur, ut Christiani no\pend
\section*{AD TIMOTH. CAP. I. }
\marginpar{[ p.465 ]}\pstart imnemur et iimus. luxta imperium dei saluatoris nostri, et lesu Christi spel nostrae.) Contra stimulum recalcitrantem me, iubet incedere, et misit legatum, ac or dinauit subditum, nihil tale cogitantem omnipotens deus, iam mundum in tempore constituto saluare uolens, iam conciliatus per Christi completum mysterium, qui nunc se exhibet piissimum seruatorem, olim uero seuerus, sed iustus iudex, et impietatum uindex durus, etiam suo peculiari populo. Is, inquam, me misit, cuius imperio subtrahere me non possum, eidem seruio, tibi Timotheo et perte Christianis omnibus scribens, non meo nomine, sed eius praecepto et uerbis, qui in me loquitur et scribit, ne humanum putetur ac mendax, aut fabula inanis, sed uerbum dei uiuum. Per excellentiam is noster dominus agnoscendus est a nobis, qui nos de seruitute liberauit, et iugo satanae, errorum atque peccatorum. Eidem domino serui dati sumus a deo patre, et per eundem filii dei facti per adoptionem, cohaeredes secum per suum sanguinem ac mortem. Non enim omnipotentis dei patris ac creatoris, sed et Christi dominium agnoscimus, qui cum in principio uerbum deus existens, perque ipsum facta sunt omnia, uerbum consubstantiale. Ne uel ludaeis uel paganis scrupulum ingeramus, si uel deum coeli, uel Christum filium, non confiteamur unum esse omnipotentem deum. Qui est spes nostra.) Vnica Christus est fiduciae nostrae anchora, fidelis in promissis, qui pro nobis satisfecit, ei cui omnia obligamur. Qui gratis nos saluos facit credentes, sine quo me ra desperatio, sicut et error, mors et infernus omnes manet. Ipse enim solus et ne mo alius est uia, ueritas et uita nostra, sicut et spes nostra. Maledictus omnis qui sperat in alium a Christo: ipse est spes nostra. Si quis spem suam iactat in Christi dignissimam matrem et dei genitricem Mariam, is uiderit quid intelligat, credat, uel confiteatur, ne Christi solius gloriam alteri tradens, matrem offendat, et rideat, titulum impingens, quem minime agnoscit, sed odit ex animo, quanto sanctior. Ne creatori hominem praeponat, ne matrem ex gratia Christi, Christo praeferat, qui natura est filius dei. Genuit quidem nobis Christum spem nostram, in quem et ipsa sperans et credens, saluari debuit. Cuius benignitati nihil addi potest, qui fons est bonitatis et misericordiae, plenus gratia, a qua et mater quoque saluata est, et omnes electi saluantur. Vel, quod bene intelligunt, bene quoque eloquantur et uelint, aut certe taceant a blasphemia in Christum, quam unicam mater beatissima maxime odit et auersatur semper. Christus autem solus clamat, Venite ad me qui laboratis et onerati estis, ego reficiam uos. Non hoc matri detulit, sicut nec minora. Quid, inquit, mihi et tibi est mulier? Ex quo tamen eidem sacratissimae crea turae nihil honoris derogatur uel felicitatis, cuius impensissimum semper uotum est, gloria filii sui, quae est eadem cum patris et spiritus sancti gloria. Timotheo dilecto filio in fide.) Cum Graecus fuerit ex patre, licet ludaea matre et Christiana, nomen habuit τιμόβεet, quod honorans deum, sua lingua interpretatur: ut uidere liceat honorifici uiri honorificum nomen. Habebat testimonium bonum a fratribus et familiaribus. Erat eruditus in lege, et disciplinis saecularibus, adhuc iuuenis, Paulo gratus et sodalis, ideo circuncidi illum permisit, propter Iudaeos superstitiosos, ut illis gratificaretur et scandalum in tempore amputaret. Factus est episcopus in Asia Ephesi. Tales autem decet esse episcopos, bono testimonio, continentes, doctos, fideles, fama et uita spectabiles, ne nomen domini blasphemetur. Dilecto dicit. Erat quippe uirtuosus et integer, amore tanti apostoli dignus. Minus autem exprimit dilectio quam amor. Tam amor enim quam charitas dilectionem superant. Filio.) Qui paternas dotes omnium charismatum ac uirtutum, ut germanus Pauli filius probe referebat. Ideo eleganter Erasmus retulit, Germa\pend
\section*{COMMENT. IN I. EPIST. }
\marginpar{[ p.466 ]}\pstart no filio, cuius nimirum Paulum non solum non puderet, sed insignem sectatorem, sodalem et coapostolum habuit. Vtgaudent patres probi de filiis sensatis, sic doctores de studiosis discipulis sibi gratulantur. In fide.) Ne carne genitum aestimares, nobiliorem indicat filiationem, sicut fides naturae praefertur, quo ad deum. Vt autem pater generando similem naturam communicat, sic Paulus huic filio fidem in Christum ueluti per manum contradidit, eamque nequaquam sterilem et inanem, sed foecundissimam, qua plurimos Christo populos genuit, et nutriuit uerbo uitae, doctrina euangelica. Fidem autem prae caeteris dotibus commendat ubique  in discipulo Apostolus, non circuncisionem, non scientiam legis, nec parentes ac stirpem generis, uel fortunas. Fides enim facit Christianum, et Christianorum antistitem, deique filium adoptiuum et gratum. Gratia, misericordia et pax.)Optat episcopo Ephesiorum non diuitias, non sanitatem, non honores, quae hodie magni habentur, et sectantur apostolici, sed dei optimi fauorem et gratiam, qua nobis beneficus est deus, intra nos bona operatur, amat ut filios dilectissimos, indignis et immeritis munera spiritualia largitur, et non destituit necessariis temporalibus. Gratia haec non est habitus aliquis qualitatis, in animam influxus, stertens et otiosus, donec aetatis progressu suscitetur a libero arbitrio ad bonum operandum, et ad meritorios actus, ut loquuntur scholastici, sed fauor dei est, et instinctus diuinus, desuper allatus a patre luminum, ad opera dei in nobis, quem ut summum bonum in hac uita nobis optare debemus. Quem et Paulus imprecatur et optat amico in primis. Si em pro nobis deus fuerit, fauoreque prosequatur, quis contra nos? Misericordia.) Clementissimus deus et pius dominus, qui nostris quidem offenditur malis, tolerat tamen propitiaturque ́, blasphemantibus etiam, ut Paulo, benefacit, ignorantibus parcit, errantes reducit, hanc nobis innotescere cupit semper, et hanc sperare nos docet. Pax.) quae est conscientiae confidentis tranquillitas. Gratiam et misericordiam dei, animum exhilarantem, similiter et externam pacem pro dei misericordia optat, ut libere a fidelibus deo uacari possit. Haec primum impensissime optabatur a Christi discipulis, inter lupos periclitantibus continuo de uita, quae pax est praedicationi ueritatis opportuna. Quae et in fine synaxeos sancto semper osculo imprecabatur. A deo patre nostro, et domino lesu Christo domino nostro.) A patre luminum desuper ueniunt nobis optima quaeque , gratia scilicet, pax et misericordia, a quo omne bonum, quae precum instantia consequi Christus nos docuit, et communi fratrum studiosaque pre catione obtinere. Proinde in nouo iam testamento, deus Pater noster uult nominari, in ueteri autem Dominus exercituum, zelotes scribitur. Fratres enim sumus Iesu Christi filii sui. Est autem pater uniuersae creaturae: specialiter autem fidelium, recipientium filium suum ac audientium. Quis igitur filiorum cunctatur adire talem patrem misericordiarum, et totius consolationis, nisi degener et nothus, impie et maligne sentiens de indulgentissimo patre pientissimoque ad suos filios? Et domino Iesu Christo.) Christus quoque secundum humanitatem, in cruce oblatam per obedientiam, dominus factus est non solum fidelium, sed uniuersae creaturae. Cuius honorem agnoscunt coelestia, terrestria et inferna. Data est ei omnis potestas in coelo et in terra. Omnes angeli et dominationes seruiunt Christo domino dominantium. Subiugatos habet reges et principarus, regna et imperia, uelint nolint tribunali suo iudicandi ab illo, comparebunt. Dominus est uniuersorum, in cuius ueste scriptum, Rex regum, et dominus dominantium. Domino nostro.) Secundo adiicitur, quod fidelium et redemptorum proprie dominus sit, qui side recognoscunt et deuotione illius dominium, mandatis acquiescunt, doctrinam  \pend
\section*{AD TIMOTH. CAP. I. }
\marginpar{[ p.467 ]}\pstart capiunt, promissis credunt, uirtutum uestigiis imitantur. Dominum eum non agnoscunt, qui euangelium suum contemnunt, uerbo non credunt, prohibent au tem, uel non pro uiribus promouent.  \pend
\phantomsection
\addcontentsline{toc}{subsection}{\textit{Sicut rogaui te ut remaneres Ephesi, cum irem in Macedoniam, ut de nunciares quibusdam, ne aliter docerent, neque intenderent fabulis et genealogiis interminatis, quae quaestiones praestant magis quam aedificationem dei quae est in fide. }}
\subsection*{\textit{Sicut rogaui te ut remaneres Ephesi, cum irem in Macedoniam, ut de nunciares quibusdam, ne aliter docerent, neque intenderent fabulis et genealogiis interminatis, quae quaestiones praestant magis quam aedificationem dei quae est in fide. }}\pstart Vide modestiam tanti apostoli ad discipulum, non mandat, non urget, non compellit, sed syncere, mansuete, paterne, Christiano pectore, ut legitimus Christi discipulus, apostolicomore rogat, quae a Christo didicerunt. Amore duci uolunt Christiani, et spiritu filiorum affici, non minis, non fulminibus, non anathae matis compelli adea quae non conueniunt. Vtremaneres, inquit, Ephesi.) Aegre diuelli poterat fidelis et docilis discipulus, a uerbis uitae aeternae, quae ex ore Pau Ii iugiter spiritus domini loquebatur. Vt Petrus quoque testatur, et populus Christum sequens, imo et Herodiani impii milites. Quem et uix quandoque nauigio poterat declinare Christus, et locis desertis. Sic nimirum Paulum aemulabatur pius Timotheus, cupiens per labores, pericula et itinera, magistrum sequi, in patriam alienam, Macedoniam, sed rogatur a Paulo, ut Ephesi maneat. Ephesus autem Asiae minoris urbs est, quae tunc cultrix superstitiosissima erat simulachri Dianae, ab Ioue delapsi de coelo, quo currebatur non solum ex Asia uniuersa, sed et abomni orbe. Vt nunc stulti Christiani Loretum uisitant et Aquisgranum, caeteraque superstitioni consecrata. Quaestum ciuitati multum importabat, ubi et Paulus biens nio praedicans, plurimum profecit in Christianis urbis illius, multa quoque passus a gentilibus. Lege caput 19. Actuum. Cui doctorem urbi Paulus Timotheum praefecit, quem antea ad Macedoniam miserat, et cum quo Ephesum redierat. Interim tamen obtestatur, ne sensu proprio nitatur discipulus, rogat ut faciat ac loquatur iuxta doctrinam et institutum Pauli, ne nouis doctrinis abducatur: nec sinat sub ingredi pseudapostolos. Vt denuncies quibusdam, ne diuersam sequantur doctrinam.) Hoc uolui, inquit, faceres Ephesi, denunciares scilicet quibusdam Iudaizantibus, et aliam ab euangelica doctrina subintroducentibus. Officium prae lati ecclesiastici primum sit et doctoris, ostendere et docere euangelio solo necessario credendum, et sacris literis, eidem soli insistendum, eo quod insertum sit ue teri testamento, ueluti figuratum in figura, et ueritas per umbram declarata. Quorum testamentorum, quae unum uere dicuntur euangelium, claritate tantum differentium fides abunde sufficit ad omnium fidelium salutem, et ad deo complacendum fide et opere. Quicquid uero propheticae et apostolicae doctrinae aduer satur, anathema sine dubio iudicetur. Aduersatur autem philosophia gentilium in quibusdam, ideo cum delectu diligenti suscipienda, non nisi profectioribus in fide, ne cui noceat mel felle uenenóue gustatum. Cui tamen philosophiae, quia ha ctenus toti instetimus, ideo hodie euangelium praedicantibus nouitatem improperamus, cum nihil antiquius sit uere Christiano. Aduersantur quoque traditiones humanae, quibus conscientiae illaqueentur, cultus dei, praeter uerbum dei, instituitur, quarum obseruantia mereantur homines uitam aeternam, transgressores post hanc uitam purgatorium, expiandum poenis uel donariis. Haec Apostolus denunciauerat minime praedicanda in Epheso apostolica ecclesia, a loanne apostolo et Paulo instituta. Subintrauerant autem Iudaei, curiosorum dogmatum obscuri et superstitiosi doctores, quorum auditores fuisse uidentur Apelles, Ebion,  \pend
\section*{COMMENT. IN I. EPIST. }
\marginpar{[ p.468 ]}\pstart Valentinus, et similes, qui de Eonum suarum propagationibus, quasi nostra aeta te theologi de Scotistarum subtilissimis formalitatibus plurima garriebant, et solum ingeniorum suorum, et reuelationum gloriam et admirationem uenabantur ob quaestum. De quibus Tertullianus et lrenaeus, in primo contra Haereticos. Hii quoque docebant genealogias infinitas, et nunquam finiendas, miras quoque fa bulas de deo et angelis, etiam apostolorum temporibus, ueluti reliquias haeresum Iudaicarum. Continuo quoque unus addebat super alterum, multa quoque diligen tia per registra sua describebant maiorum nomina, uolebatque quilibet de nobilio rifamilia natus prae aliis uideri. Fabulas denique Paulus appellat, quicquid extra dei uerbum et sacras docetur literas. Vt est totum Iudaicum illud Talmud, et ob scurissima dubia, inuolutaque cabala. De quibus omibus subinde sequitur. Quae quaestiones praestant magis quam̃ aedificationem dei, quae est in fide.) Cauendum sum, mopere Paulus docet ab eis dogmatibus, quae ansam praebent quaestionum: non autem dicit quarum et qualium quaestionum. Certe uanas, curiosas, gentilitias, philosophicas, metaphisicas quaestiones, omnium nostrorum temporum scholasticorum et summariorum digne tractando damnasset et prohibuisset Paulus. Quibus, ut experientia infelici didicimus, nihil discitur, nisi de omnibus dubitare, contradictoria propugnare, fallendo uincere, uincendo uideri, cathedras quoque  et loca inter illos digniora sortiri. Quis hodie numeret quaestionum libros, et quaestionariorum illorum nomina, qui scripserunt in magistrum illorum Lombardum Petrum? etiamsi quis Tritemii scriptores ecclesiasticos, non omnes enumeratos, recenseat. Interim tamen uix unus e centum habeatur, qui sacram tractauerit scripturam totam. Quin et qui tractauerunt sacra, ut D. Carensis, et alii quidam, eo modo tractarunt, ut praestet non tractasse. Quasi uel Petrus uel Paulus Christi no mine iusserint huiusmodi quaestiones moueri tractarique ́, sic studuerunt iam a trecentis annis quaestionibus. Ideo euanuerunt cogitationibus, et spiritui ecclesiastico utilia neglexerunt, imo contemptui habuerunt, ut protrita, et minime erudita. Quaestiones enim praestant, non soluunt. Semper enim uani isti habent quod quaerant, et stulti quaerere malunt, q̃ doceri. Cum Christus pietatem docuerit, simplicitatem, integritatem et fidem. Aedificatio autem dei scripturis fit sacris, in quibus fides, spes et charitas ueraque docetur pietas, quam isti philosophastri et quaestionarii ignorant, nunquamque assequuntur. Quae est per fidem.) Catholica do cumenta non habentur nisi per fidem: nec aedificatur ecclesia, nisi sacris literis et fide, quae sola deus credi mandauit atque doceri. Estque ueri Christiani hominis profectus tantum per fidem. Vt quoquis fide in deum et uerbum suum magis profi ciat, eo sit sanctior, et pie eruditior. Opera denique fide iustificant, et domino placent, sine fide nulla aedificatio, sed subuersio est.  \pend
\phantomsection
\addcontentsline{toc}{subsection}{\textit{Finis autem praecepti est charitas, de corde puro, et conscientia bona, et fide non ficta. }}
\subsection*{\textit{Finis autem praecepti est charitas, de corde puro, et conscientia bona, et fide non ficta. }}\pstart Praeceptum deique uoluntas, charitate perficitur, ex fide in Christum saluatorem nostrum, et reconciliatorem, qui a deo patre missus est nobis saluator. Fides autem per charitatem fructificat, cum opera charitatis sint fidei sequelae et fructus, ut sine eis fides non sit nec appareat. Charitas autem duplex, Dei, qua beneficia dei agnoscimus ex illius bonitate et gratia, qui fidem exigit et praestat, fideque placatur et colitur, ea qua toti a sua uoluntate pendemus, et nosmetipsos eius bonitati committimus, iugiter suspensi ad beneplacitum illius, in omnibus operibus nostris perficiendum: hoc non fitsine fide uerborum dei, et facrarum literarum  \pend
\section*{AD TIMOTH. CAP. I. }
\marginpar{[ p.469 ]}\pstart praeambulo. Denique charitas proximi necessaria fideli est, et ex flde operatur. Si non operatur, non est charitas, nec fidem ostendit. Noli iactare fidem et Christianismum, sine operibus charitatis. Non satis facies deo uerbo tantum et confessio ne fidei: uiuax est, actione augescit et proditur, sine exercitio languescit, imo expirat et moritur, uel mortua manifestatur. Nec ficta charitatis opera iuuant, sed e puro corde profecta, fide uera purificato, ut quae bona fiunt pro dei tantummodo gloria fiant, sine respectu mercedis uel praemii. Ne summo bono iuxta Augustinum, uti uideamur, et hypocritas agere, duploque domino displicere, dum mali tia nostra illius bonitatem et omniscientiam fallere desideramus et tentamus. Conscientia bona, ait, et fide non ficta.) Explicandi gratia haec adduntur, cum idem sit puro corde, bona conscientia, et germana fide agere domini uoluntatem, et benefacere proximo.  \pend
\phantomsection
\addcontentsline{toc}{subsection}{\textit{A quibus quidam aberrantes conuersi sunt in uaniloquium, uolentes esse legis doctores, non intelligentes neque quae loquuntur, neque de qui bus affirmant. }}
\subsection*{\textit{A quibus quidam aberrantes conuersi sunt in uaniloquium, uolentes esse legis doctores, non intelligentes neque quae loquuntur, neque de qui bus affirmant. }}\pstart Errare necesse est omnes, qui fabulis et quaestionibus insistentes, dei uerbum hominibus saluandis propositum negligunt, uiam uidelicet ueritatis, scripturam et Christum, quem omis docet scriptura noua et uetus, cuius scopus est unicus, et finis totius legis, ipsa fraterna charitas. Otiosi enim a cura proximorum, et Chri sti mysterium ignorantes, et minime iuxta Christum uiuere cupientes, quid aliud agerent, q̃ sensu humano curiosa, uana, superflua, nihil ad pietatem facientia, cogitantes, inuoluerent sese animi gratia talibus quaestionibus, tantummodo seipsos quaerentes. Doctores legis se uocant.) Soli autem sunt qui uerba dei reuerentur, quod credentes humiliat, et sensu deprimit, nec sinit intumescere. Solum enim Christum et spiritum Christi doctorem agnoscunt uere pii: qui enim uere sapiunt, humiliantur, et ignorantiam plurium fatentur: quanto magis proficiunt, tanto plura se ignorare deprehendunt, inscitiam suam agnoscunt ingenue. Nec doctio res uideri uolunt quam sint, curtam suppelectilem suam intuentur, non ambiunt ti. tulos, non iactant magnalia, sed laudati erubescunt. Non intelligentes etc.) Liquet legentibus Valentinianorum et Apellis, aliorumque illius temporis uanorum scripta doctorum, quam plura scire se uideri ambierint, uerbisque inusitatis sua commenta inuoluerint, praeter scripturam plurima somniantes. Sic autem loqueban tur, ut nec ipsi intelligerent, nec hii quibus promercede et quaestu uerba facie. bant. Quanta autem sit in Scholasticis nostris quaestionum obscuritas, uel inde li quet, quod perpetuo consequenterque et semper discipuli incertiores a magistris discedunt, quaestionibus sese fatigant, argumentis inuoluunt, quibus se nequeunt explicare. Nimirum regula ueritatis amissa, canonica scriptura, consequi aliud non potuit, que ut a Christiana synceritate defecerimus, gentilium et ludaeorum sa pientia et superstitione degenerauerimus, lumen uero animarum perdiderimus, interim intellectum diuinae legis, fidem quoque certam, et amorem dei ac proximi amiserimus.  \pend
\phantomsection
\addcontentsline{toc}{subsection}{\textit{Scimus autem quia bona est lex, si quis ea legitime utatur. }}
\subsection*{\textit{Scimus autem quia bona est lex, si quis ea legitime utatur. }}\pstart Dubitare nemo potest, legem et quicquid a deo est, bonam esse. Cumque lex doctrinam sonet, uitae honestae institutionem, a malo quocunque dehortationem, ad finem utique bonum ordinatur, ad charitatem scilicet dei et proximi, ut necessario bona uideatur et sit. Verum nihil tam bonum, quo non possumus abuti ad malum: quinimo nulla legis quoque obseruatio ad literam bona est, nisi seruiat chari.  \pend
\section*{COMMENT. IN I. EPIST. }
\marginpar{[ p.470 ]}\pstart tatt: finis enim praecepti charitas est. Quam si quis negligat, nihil profuerit mora dicus obseruata litera legis: abusus est potius, si militet aduersus charitatem, quae uera sit charitas, quam solam scopum habere debet quicquid constituitur inter homines. Quin et deus quoque suam uult legem charitati seruire, ut patet de sabbatho in Euangeliis, et de Dauide, panes lege prohibitos, cum suis in necessitate fugiens, comedente.  \pend
\phantomsection
\addcontentsline{toc}{subsection}{\textit{Scientes hoc, quia lex iusto non est posita. }}
\subsection*{\textit{Scientes hoc, quia lex iusto non est posita. }}\pstart Non abutitur lege, qui nouit finem legis esse charitatem, quam qui cum fide habet, sine qua haberi non potest, iustus est. Iustitia enim ex fide est, quae per charitatem operatur. Eidem lex non imponitur ut onus et iugum, sed ut regula actionum, directio sancta et utilis, et diuinae uoluntatis insinuatio, quam pius non ob seruatut legem, et quasi ex necessitate, sed in spiritus libertate, ex amore recti, et gaudio gratificandi deo summo bono, cui haec certo placere non ambigit, lege dei admonitus. Non minus facturus, etiam sine constrictione legis coactus. Non enim in sic iusto spiritum libertatis impedit lex imposita, sed dirigit et certum facit, deo placere quod agitur, a quo legem accepit, ut studiossus et alacrius operetur, domino obtemperaturus pro uirili.  \pend
\phantomsection
\addcontentsline{toc}{subsection}{\textit{Sed iniustis et non subditis, impiis et peccatoribus sceleratis et contaminatis, patricidis, et matricidis, homicidis, et fornicariis, masculorum concubitoribus, plagiariis, mendacibus, periuris. }}
\subsection*{\textit{Sed iniustis et non subditis, impiis et peccatoribus sceleratis et contaminatis, patricidis, et matricidis, homicidis, et fornicariis, masculorum concubitoribus, plagiariis, mendacibus, periuris. }}\pstart Iniustis, inquit. Qui enim sine fide uiuentes iniusti sunt, dominum non timent, iudicium non uerentur, uitam futuri saeculi non attendunt, uel non credunt, legem dei non aduertunt. Hos necesse est a malis studiis suis, non tantum dei legibus manutenendis, sed et ab hominibus statuendis, legitimis cancellis retineri, coërceri et puniri: non tam in ipsorum bonum, quam eorum quibus cum conuersan tur. Non subditis, inquit, uel inobsequentibus.) Vt qui sine lege quacunque sua potius sponte sanctis monitis a malis retrahi non uolunt, legis timore, et poenarum interminatione in suum bonum a uetitis arceantur et malis. Impiis.) addit. Qui deum non credunt uel contemnunt, aut non curare nostra suspicantur, ideoque si ne fide quicquid libet agere non uerentur, horum malignitati lege obsistendum est, ut si non dei iram, saltem hominum poenas, legibus ferendas timeant, et iuxta eas uiuere compellantur. Pii autem homines non expectant legum mandata, sed suapte synceritate et probitate sancte conuersantur et agunt, dulcique charitatis tractu uinciuntur, potiusque praecurruntquam ducantur: tantum abest, ut cogi uelint uel debeant. Peccatoribus.) Cum simus peccatores omnes, quantumlibet religiosi et pii, ut nec deo amatissimus discipulus Ioannes et apostolus, se peccatorem audeat diffiteri, nec quisquam sanctorum, sequitur, omnibus quidem esse legem positam hominibus, quia peccatores. Nihilominus tamen iusti sunt et pro tanto exleges, quia in fide et charitate uiuentes, sine lege, hoc est legis coactione, bona libere, sponte et alacriter agunt. Qui uero ad legis necessitatem coacti uiuunt, et, ut fit, legem ipsam ostendentem peccatum odiunt, hi uere peccatores et iniusti dicuntur, poenitentia et conuersione indigent, donec tales fuerint, non saluantur. Sceleratis.) Erasmus Irreuerentibus. Qui non sua sponte maioribus reuerentiam obediendo exhibent: hi, ne rerum ordo turbetur, legibus disciplinandi sunt, ut cui honorem, honorem: obedientiam, obedientiam, et c. Honorem enim maioribus per contemptum negantes, pacem rumpunt et concordiam, lege ergo coërcendi sunt, et per magistratum humiliandi. Contaminatis,) Impuris uel pro\pend
\section*{AD TIMOTH. CAP. I. }
\marginpar{[ p.471 ]}\pstart phanis. Soli ergo puri, mundi et sacri, quorum mentem dei spiritus inhabitat, sine legis necessitate, non tamen sine lege uiuunt, omnia enim dei mandata studiose conseruare satagunt: et si id minus succedat, indolescere non tardant, nec emen dare. Habent enim dei legem, et beneplacitum in uotis semper, cui sponte obediunt. Contaminati autem dicuntur, qui omnium spurcitiarum genere se foedos ostentant. Patricidis, matricidis, homicidis.) Hii legibus opus habent, et ut formidine poenae malum oderint, qui carent amore uirtutis, qui naturam malitia su perant, contumeliam quoque ingerunt creatori, et bestias superant crudelitate. Ideo tam graues contra hos leges etiam homines statuerunt, quibus tam incon, dita malignitas et atrocitas compescatur: interim tamen uidemus tam efferas quas dam gentes, apud quas crudelissima poenarum genera uix proficiunt, quo minus grassentur homines in seipsos. Scortatoribus.) Hic uidemus fornicationem et scortationem inter crimina a Paulo computari, secundum etiam dei leges, ne quis licitam mentiatur, hominum generationi et educationi publicoque conuictui ma xime contrariam. Masculorum concubitores.) Crimen infandum secundum spe cies suas non paucas, deo, naturae et hominibus abominabile, ac coelesti ultione dignum, acerbis quoque et duris legibus uindicandum. Infidelium hominum peculiare crimen, et dei contemptorum, qui deum non credentes, uel omnino postponentes, delicias tantum suum deum statuunt, et humanam denique sapientiam di uinae praeferentium. Vtraque uitia, legibus matrimonii dominus reprimere statuit, quibus contra, coelibatus ex idololatrarum consuetudine maxime suffragatur. Plagiariis.) Qui homines pro mancipiis furto subducunt, rapiunt, uendunt, abu tuntur ad crudelitatis obsequia, omnis humanitatis interim obliti, ut pyratae, et tyranni infidelissimarum gentium. Tales autem plagiarii uel fures manciparii di cuntur. Quemadmodum qui sacra furant, sacrilegi: qui de publico aut principis fisco, depeculatores: qui iumenta aliena abducunt, abigei: ita qui liberos aut seruos alienos, plagiarii dicuntur apud Iureconsultos. Mendacibus.) Legibus cer te coercendi suntomnes, qui uerbis circumueniunt, decipiendo, nocendo, irridendo, mentiendo, ueritatem negando, impugnando, conculcando, inuoluendo mundum erroribus, falsis opinionibus imponendo simplicibus, et cauponando ueritatem, ut quaestui satisfiat. Periuris.) His qui in dei contumeliam, ac diuini nominis dedecus, blasphemiam, domino impingunt, aut falsitatem, mendacium, infidelitatem, iniustitiam uel ignorantiam, uel qui se uerbis huiusmodi testantur esse homines sine deo. Sed et hi qui quae dedecori sunt maiestati diuinae, quolibet modo impingere conantur, tentant uel imprecantur. Omnibus his leges dentur, et iuxta leges rigide sunt plectendi.  \pend
\phantomsection
\addcontentsline{toc}{subsection}{\textit{Et si quid aliud sanae doctrinae aduersatur, quae est secundum euangelium gloriae beati dei, quod creditum est mihi. }}
\subsection*{\textit{Et si quid aliud sanae doctrinae aduersatur, quae est secundum euangelium gloriae beati dei, quod creditum est mihi. }}\pstart Seruantibus sane euangelicam doctrinam, a Christo et apostolis et ab ipso de nique Paulo praedicatam, quae sola sana est, et credentibus electisque necessaria, utilis et certa, non est in foro animae et salutis necessaria lex quaecumque humana, quan doquidem nec Mosaica, quae dei lex est, Iusti enim sunt qui euangelio credunt, simpliciter uerbo dei acquiescunt, proximo faciunt quae sibi fieri cupiunt, Christum amplectuntur ut suum saluatorem et fratrem. Quicquid autem huic aduersatur doctrinae, hoc lege quoque prohibitum est, et legis uirtute opus habet, quae reuincat, deterreat puniatque peccatores, transgressores et iniustos. Vt qui charitatis uin culo dulci deo connecti nolunt, legum laqueis innectantur, alioqui iustis non est  \pend
\section*{COMMENT. IN I. EPIST. }
\marginpar{[ p.472 ]}\pstart lex posita. Euangelium.) non tam praemium fausti nuncii nobis est, quam ipsa laeta denunciatio, maximorum ac spiritualium bonorum promissorum, ad laetificandas conscientias piorum hominum, quae et in primis ac insigniter in usum susce pta est praedicari, de mysterio dei ad homines per Christi domini promissionem, incarnationem, praedicationem, beneficia, passionem, mortem, et resurrectionem. Quae omnia nunc instant, et usui solatioque sunt credentibus, quia scilicet deus per ditum mundum in Christo sibi perfecte reconciliauerit, et omnium peccatorum satisfactionem per Christum, et in Christo iam receperit. Vt non tam de peccato rum nostrorum condigna satisfactione, et ad deum nostrum conciliatione ac remissione culpae simus soliciti, quam de poenitentia et emendatione, ut posthac pientissimi patris beneplacitis toto affectu et desiderio insistamus, utque grati filii in fide et charitate, toti illi summo bono nostroinhaereamus. Euangelium autem glo riae nominauit. Quid enim deo gloriosius, q̃; sua bonitas, per Christum et euangelium nobis reuelata ? Docet enim nos, ut non tam gloriosum sit illi, quod in prin cipio per suam omnipotentiam creatus est mundus, quam quod in fine temporum pro nobis immolatus est Christus. Glorificatur euangelio filii sui deus toto terrarumorbe, deprimit sibi, et subiicit omnem gloriam huius mundi, omnem princi. patum, omnem uanam sapientiam, glorificat obseruatores, ac credentes facit uere sanctos et sapientes: cui quicquid contrariatur uel repugnat, imo non concor dat, uanum est. Beati dei.) Deficiunt Apostolo uerba, dominum deum honoris digno epitheto adnominare cupienti. Quicquid enim de deo dicitur, minus est a dignitate dei. Beatum simpliciter dicit, quia omnis beatitudinis fons, author et sco pus est, quem scire, amare et tenere consummata felicitas, et absoluta est perfectio. Cui nihil bonorum abest, a quo solo omne bonum, qui denique ad capessendam beatitudinem electos omnes creauit. Denique omnes homines saluare cupiens, pro quibus omnibus omnia fecit sancta ac iusta uoluntate sua. Creditum.) sibi euangelium Paulus affirmat et praedicat, de quo et gloriatur, ut in eo deo suo seruiat, fideliterque incumbat, quamlibet non fuerit de numero duodecim apostolorum, sed persecutor ecclesiae. Ex gratia tamen uocatus, et missus ad gloriam euangelii disse minandam, et gratiam dei per Christum inter gentes, tanquam fidelis, non adulterator quaestuosus et uanus, frigidus uel deses. Creditum fideliter praecspuo legato, ut thesaurum inexhauribilem in orbem uniuersum spargendum: semen quoque fructum centesimum prolaturum.  \pend
\phantomsection
\addcontentsline{toc}{subsection}{\textit{Gratias ago ei qui me confortauit in Christo lesu domino nostro, quia fidelem me existimauit, ponens in ministerio. }}
\subsection*{\textit{Gratias ago ei qui me confortauit in Christo lesu domino nostro, quia fidelem me existimauit, ponens in ministerio. }}\pstart Quanquam fuerim magister legis, exercitatus quoq in uaniloquio scholae Pha risaeorum, magister noster, et humanis literis probe eruditus, uel non indoctus, gratias tamen ago deo meo, qui per infusam mihi fidem in dominum lesum Chri stum, me licet indignum, sola sua gratia dignum fecit, quem ministerio suo subde ret, ut praedicem dominum lesum Christum, sine omni respectu priuati uel pubus  lici commodi mei. Sed in gloriam dei patris, et in laudem gratiae bonitatis suae, an nunciandae per me in uniuersum orbem. Eiusmodi autem uocatio ad spiritualia bona, ac uerbi dei praedicationem, soli deo accepta referatur: non enim talemmereri nostrae est facultatis uel arbitrii, praeuenimur ad id dei bonitate, et gratia gra tis data, ut nihil de nobis gloriemur. Sed dominum honorificemus, qui uocat quos elegit in opus suum, et uocationi adaptat:et quicquid per nos promouetur boni, ipse agit, promouet ac absoluit, ut ipsum in suis beneficiis agnoscamus, et per fi\pend
\section*{AD TIMOTH. CAP. I. }
\marginpar{[ p.473 ]}\pstart dem et gratiae exercitium recipiamus maiora. Qui me confortauit.) Cum nihil propriis possem uiribus, fortissimum tamen euangelii sui propugnatorem effecit super coaetaneos meos, dudum iustitia legali inflatum humiliauit, et exinaniuit, ut nihil tribuerem naturae uiribus, perque suam gratiam confortauit in omni bono. Nouo spiritus Christi feruore imbuit et inflammauit, ut non solum arserim in amo re dei per Christum, sed etiam uehementer multos uerbi praedicatione alios passim inflammauerim. In Christo Iesu domino nostro.) Omnia autem haec acta, accepta dictaque sunt in gloriam et laudem Iesu Christi, quem pater omnipotens et dominus, mundo lumen, recreatorem, redemptorem ac doctorem instituit, cu ius uerbi euangelici me praedicatorem dedit, cuius quoque gratia et meritis confor tatus, deo gratias ago. Quia fidelem me iudicauit.) Non quidemre uera fidelem inuenit, qui fastuosus Pharisaeus eram, sed per meram suam gratiam sum id quod sum, qua et decreuit me facere fidelem sibi, quem praesumptuose fidelem reperit propugnatorem traditionum paternarum, scilicet ludaicarum. Vt tanto clarius elucesceret eminentia bonitatis diuinae, quanto indigniorem deligeret in tantum opus, ebuccinandi gloriam dei in omnes gentes per Christum filium suum. Po tens in ministerio.) Vt sim non propriae quidem doctrinae inuentor, sed uerbi sui humilis minister: idoneus tamen et non piger, nec propriae utilitatis uel gloriae ca ptator, sed domini dei minister fidelis et sedulus in opus sanctum, non in gloriam praelationis, et regimen mundi in temporalibus quibuscunque , sed in ministerium subiectionis, ut eidem ministrare nunquam cessem, sed indefesso conatu impleam cursum ministerii quod suscepi. Cuncta enim non solum apostolica, sed Christia na in genere uocatio, alia haberi non debet, q̃ ministerium quoddam charitatis, a Christo nobis commendatae, qua mutuum gratificando famulemur, propter quod et uoluitin ecclesia multos relinqui semper imbecilles, infirmos ac paruulos in Christo nutriendos, uerbo et exemplo, alterum indigere alterius opera, do naque ministeriorum singula singulis impertiri pro singulorum necessitatibus.  \pend
\phantomsection
\addcontentsline{toc}{subsection}{\textit{Qui prius eram blasphemus, et persecutor, et contumeliosus, sed misericordiam dei consequutus, quia ignorans feci in incredulitate. }}
\subsection*{\textit{Qui prius eram blasphemus, et persecutor, et contumeliosus, sed misericordiam dei consequutus, quia ignorans feci in incredulitate. }}\pstart Vel hinc liquet immeritum me sola gratia uocatum a domino, in tam sacrosanctum ministerium, quod iam ante Christum eundem dominum meum, quem nunc praedico dei filium, Christum saluatorem mundi, blasphemo ore criminatus sum coelestem eius doctrinam, ac dei patris, mundo per eum exhibitam charitatem, stupenda quoque eius audita non uisa miracula credere nolui: dixi, ut caeteri pharisaicae scholae sectatores, daemoniacum fuisse eum, haereticum, seductorem, plus quam peccatorem, et amicum fautoremque peccantium, contemptorem rerum diuinarum pariter et humanarum. Docui haec eadem alios, seduxi quos potui a do ctrina Christi: quos uerbis non potui, factis compuli. Persecutor.) Non sufficie bat oris blasphemia, et clamor sacrilegus, quin et ubique persequebar discipulos euangelii et doctores. Primum quidem Hierosolymis discipulos, Stephanum, sed et epistolas accepi ad omnes synagogas, ut si quos inuenissem huius uiae uiros et mulieres, uinctos perducerem in Hierosolymam. Non nego impietatis meae infelicitatem, sed confiteor me peccasse supra modum in Christum et Christi eccle siam. Nec solum persecutus sum, sed usque ad mortem inquisiui, ubique protraxi, tra didi, accusaui, incitaui crudeles, damnaui, lapidaui, occidi. Sed misericordiam adeptus sum.) Vbi superabundauit crimen, ibi et exuberauit gratia: praeuenit me non solum sine meritis, sed pessime incedentem suscepit clementer. Impediuitin  \pend
\section*{COMMENT. IN I. EPIST. }
\marginpar{[ p.474 ]}\pstart malo, corripust clamore e coelis, indicauit errorem, suppeditauit doctorem, non solum de ecclesia Damascena, sed etiam de coelo, et de tertio coelo, et ad supereffluentiam, ut plura quae didici ibi, eloqui non expediat. Addidit animum, publice confitendi Christi gloriam, et euangelicam ueritatem, utque subinde toto pecto ris ausu defenderem, quos prius uastator acerrime fueram persequutus ubique . Quae haec mutatio fieri poterat, nisi a dextra excelsi? Quomodo conuerti potuissem ab inolitis doctrinis, nisi gratia dei benigne praeuentus? Tanto celebriorem in me suam gloriam auxit, quanto blasphemias meas ac crudelitates patentiore grassatione euulgari constituit. Quod ignorans feci per incredulitatem.) Repletus eram scientia traditionum humanarum, quam a magistris nostris didiceram: promissiones diuinas, de modomittendi Messiam et Christum suum, nemo illorum docuerat. Legis iustitiam sufficere mihi abunde putabam, praeter quam nullam sciebam. Leuitici sacerdotii ceremonias maximi ducebam, templum crebro, utreligiosior, uisitabam. Videns itaque per Christi doctrinam et apostolorum, alienari fratres meos ab humanis ac Iudaicis illis traditionibus, quas paruifaciebant, Christum dei filium asserebant uerbum dei, patri coaeternum, et coomnipotens, et coaequale, sabbathum infringentes, manus non lauantes, paucas uel nul las templo hostias offerentes, orationes publicas et solennes priuatis et clancula riis ad deum desideriis postponentes, sacerdotum et scribarum authoritatem ele uantes, et caetera plura. Existimaui, falsa scientia excaecatus et ignorans, omnes tales cum Christo illo suo iustissime eliminandos, occidendos, memoriam nominum ac personarum delendam, hoc ipso uoto arbitratus obsequium me praestare deo, zelosus sine uera scientia. Facile dei misericordiam consequimur, si ignoranter peccamus, nisi ignorantiam maligne affectemus peccato in spiritum sanctum, qui inspirat, cuius inspirationem ad bonum detestari, ex certa scientia, est in spiri tum sanctum peccare. Excusatio autem manet ac relinquitur ei, qui per incredu litatem peccat, studiosusque uoluntatis diuinae existens, et alioqui pius, non intelli git aliquid secundum scripturas esse necessario credendum, ut hic Paulus, idque ob defectum doctorum. Non uiderat Saulus Christum, miracula ex auditu tantum didicerat, utriusque rumoris authores audierat. Hinc pontifices et scribas ac pha risaeos, Christianorum et Christi aduersarios iuratos et infestissimos, quos ipse suspiciebat tanquam primores, et ueritatis indagatores credebat. Haud aliter ac hodie tota fere Gallia suos moratur doctissimos Parisieñ. theologos Sorbonicos, magistros nostros. Hinc simplices idiotas, humiles apostolos, et doctrinam crucis despiciebat ut scandalosam, parum credebat miraculis, homo natura suspicax pro doctrina suorum magistrorum, et tenax traditionum paternarum. Ideo incredu litas fere eidem inuincibilis uideri poterat, ut misericordiam dei assequeretur ob ignorantiam, alioqui zelosus scripturae discipulus et aemulator existens. Non sic de pharisaeis aliis ac pontificibus, qui Christum audierant, miracula uiderant, sciebant quia a deo uenisset magister, et ꝙ nemo haec signa facere poterat, nisi deus fuisset cum eo. Peccabat tamen interim Saulus, non curans audire apostolos, pro phetas, et Christi uerba ac opera non conferebat, nec scripturas scrutabatur, implicatus doctrinis pharisaicis et traditionibus, magnatorum magis iudicio, quam propriae conscientiae recollectioni intentus, synceritati simplicium ac bonorum minus credebat, tumidus corde, scientia animali, et consuetudine maiorum irretitus. Quid hic fusius dici possit, satis norunt, qui spiritu sunt pietatis et sacrae eru ditionis imbuti. Impii suo malo pergunt detestari scripturas sacras et fidem, mor dicus autem et contumaciter defendere usum, quaestum, luxum et fastum.  \pend
\section*{AD TIMOTH. CAP. I. }
\marginpar{[ p.475 ]}
\phantomsection
\addcontentsline{toc}{subsection}{\textit{Superabundauit autem gratia domini nostri cum fide et dilectione, quae est per Christum lesum. }}
\subsection*{\textit{Superabundauit autem gratia domini nostri cum fide et dilectione, quae est per Christum lesum. }}\pstart Fauor et clementia dei, qua me praeordinauit ad suum obsequium, eatenus copiose me prosequuta sunt, ut quia peccator eram semper, promereri illam non poteram, alioqui gratia dici minime potuisset. Sed hostem gratiae in Christo, qui electos, licet ignorans persequutus sum, elegit me ad portandum nomen suum co ram gentibus, et sibiipsi in seruum, conferendo fidem et agnitionem mysterii sui per Christum, dispensati ad omnes gentes, eamque dilectionem saluatoris lesu Chri sti, et ecclesiae electae ex gentibus, quam nullis meis studiis obtinere, sed per gratiam Christi Iesu suscipere poteram clementer indultam. Eandem denique tanto gradu donauit, qui merito superabundans diceretur, quia multum peccaueram nimio et irrationabili zelo, cum ignorantia tamen excusabili penes Christum, quam electo nollet imputare. Supra modum autem gratia collata dicitur, quia sine prae dicatione, spiranti minas et caedes in Christum, hoc est in Christi discipulos, reue latio tertii coeli delata est, Christoque doctore imbutus est de altissimis mysteriis dei, Idque subito spiritu sancto replentem et consummante, ut Christi domini magis appareret maiorque benignitas. Fide insuper uera implebatur per gratiam, quae sine dilectione suscipi non poterat, nec subsistere: ideoque magna fides in Paulo, amplam quoque charitatem ostendit, quae eum plus omnibus laborare faciebat, et solicitum reddebat de singularum ecclesiarum profectibus, ut prius omnes uoluerat profligasse. Gratia igitur Christi donata cum fide et dilectione, opera fidelium commendat, non legis Mosaicae opera, nec alia quaecunque : non enim iustificant quen quam, sed sola fides per Christum mediatorem, qui est nostra iustitia, obseruantia digna legis, et remissio peceatorum.  \pend
\phantomsection
\addcontentsline{toc}{subsection}{\textit{Fidelis sermo et omni acceptione dignus, Quia Christus lesus uenit in mundum, ut peccatores saluos faceret, quorum primus ego sum. }}
\subsection*{\textit{Fidelis sermo et omni acceptione dignus, Quia Christus lesus uenit in mundum, ut peccatores saluos faceret, quorum primus ego sum. }}\pstart Certissima traditio fidei nostrae est, et indubitabilis promissio dei, articuluś apostolicae praedicationis, cuius fides Christianum facit, et fidelem. Sermo, quo nihil certius dici potest, et credi utilius, amplecti quoq suauius et felicius conser uari. Omni acceptione dignus.) Nihil mundo optabilius potuit reuelari, plausi bilius euangelium annunciari nequit. Cuius praedicatio omnium antiquissima, statim primis parentibus, et omnibus patriarchis promissa credita et sperata fuit. Cuius fide saluati sunt omnes, quam fidelissime crediderunt, et desiderantissime praestolabantur implendam, quam nunc nos dei misericordia gaudemus impletam, et gratias agimus. Argumentum est autem et scopus non solum totius huius epistolae, sed euangelicae et apostolicae praedicationis, et summa pietatis fideique ́ Christianae. Christus Iesus uenit in hunc mundum peccatores saluos facere). Qui dudum in patriarchis et prophetis loquebatur, nunc dixit, Assum, de patris sinu sponte obediens et hunc mundum perditum cernens, misericordia motus, redenptor aduenit. In hunc tam miserabilem mundum errore, ignorantia, infirmitate malitiaque repletum per omnia. Fuit quidem semper omnia in omnibus, nulli unquam abfuit creaturae. Gubernauit et impleuit mundum diuinitate sua, sed in hominis forma ac ueritate nunquam in mundo prius hoc modo apparuerat. Nunc utero uirginis egressus in hunc mundum, uerus homo et deus, pro homine certauit cum daemone, quem deus uicit per hominem. Omnes quidem peccatores reperit, ideo omnes saluare uenit, non iustos, qui uel sint uel esse se putent: nemo enim saluatur nisi humiliatus, et peccati sui conscius, ut humiliter redemptorem  \pend
\section*{COMMENT. IN I. EPIST. }
\marginpar{[ p.476 ]}\pstart desideret, de seipso desperans, iuslitiam suam nihili faciens, misericordiam petens, et indulgentiam promittenti credens. Saluos autem facere, est de peccato liberare: non quidem ne insit donec uixerimus in hoc corpore, sensus deprimente, sed per fidem et charitatem deo unire, amicum deo facere, peccata omnia non ad da mnationem imputari, conscientiae certam consolationem imprimere, deque dei mi sericordia et gratia certum reddere, pro totius mundi delictis deo patri satisfacere, infernum claudere credentibus, deo autem semper uiuere. Haec breuissima est euangelii summa, alta mente reponenda fidelibus, et in spiritus dulcedine semper ruminanda, simulque magnitudo peccati, quod tanti redemptoris necessitatem in duxit, ut dei gratiam magnificemus, et culpam nostram grauissimam agnoscamus, ut humiliemur. Quorum primus ego sum.) Ego, inquit, ut Christum dominum primus sum persecutus, et in eius membris occidi uoto et facto, ecclesiamque acerrime afflixi, et primum martyrem coronaui. Sic primo omnium, tantus inimicus ego, ut alter Semei Beniamita, Christi domini Dauidis filii clementiam et gratiam sum expertus. Ego Christum blasphemaui, ipse me uas electionis gentibus praedicatorem transmisit: ego ipsum e mundo profligare cupiebam, ipse me sectatorem fecit: occidere uolebam et mactare, ipse me eorundem propugna torem fidelium instituit: nomen Christi exosum habui, quod nunc praedico et magnifico: Iudaeos tuebar, quos nunc impugno atque confundo. Omnia secus atque uo lebam a domino miserante cesserunt, ut miraculum conuersionis meae clarius in dei gloriam elucesceret.  \pend
\phantomsection
\addcontentsline{toc}{subsection}{\textit{Sedideo misericordiam sum consequutus, ut in me primo ostenderet Christus lesus omnem patientiam, ad informationem eorum, qui credituri sunt illi in uitam aeternam. }}
\subsection*{\textit{Sedideo misericordiam sum consequutus, ut in me primo ostenderet Christus lesus omnem patientiam, ad informationem eorum, qui credituri sunt illi in uitam aeternam. }}\pstart Sanctos peccare nonnunquam sinit enormiter, ut sanctiores, feruentiores, hir miliores, fideque et fiducia firmiores resurgant, de seipsis nihil boni uel non nimium praesumant. Omnia tribuant clementiae dei, in quam solam sperent, ardentius pro peccato orent, sibiipsis ex animo displiceant. Petrum et Paulum apostolos summos, ac Christianorum doctores, ob id peccare permisit, ut clementiae suae speculum et exemplar intuerentur omnes, non quidem ad peccandum, ad quod natu ranimium prona suopte genio impellit, sed ad non desperandum, humiliter fidendum, ad informationem fidelium qui credituri erant. Si enim tanto persecutori tanta deus benignitate indulsit tanta crimina, quid nos non miseri et humiles peccatores speremus ?Interim quoque discimus, neminem iudicare, de nullo cri minoso desperare, neminem pro deplorato abiicere uel damnare, neque in sanctis quoque personis perpetuam innocentiam exigere. Quandoquidem nonnunquam humilis peccator gratior deo sit, quam sanctulus superbus de sua sanctimonia praesu mens. Consequimur autem misericordiam, non nostris uiribus, operibus, poenitentiis, sed gratuito dei dono sua nos gratia praeueniente. Omnem autem patientiam.) Hoc est absolutam ac perfectam ostendit dominus in sancto suo apostolo Paulo, ut maiorem exigere nullus possit. Non graues iniungit poenitentias, non dure uindicat admissa, sed gratis absoluit abomnibus: prioris tantum uitae ueram conuersionem emendationemque iniungit, ut culpam agnoscamus, pro ea humiliter ueniam inquiramus, firmiterque obtinendam speremus. Ad informationem eorum, etc.) In Paulo siquidem deus uelut exemplar proposuit mundo, ut discant quid sibi sperare debeant, quicunque relictis moribus pristinis, crediderunt in Chri stum. In uitam aeternam, ait: fortassis ad dei misericordiam referendum est, qua  \pend
\section*{AD TIMOTH. CAP. I. }
\marginpar{[ p.477 ]}\pstart tredentes, quantumlibet peccatores, per clementiam informat, et adducit in ui, tam aeternam, quam nullis nostris meritis promeremur, sed dei misericordia credentes assequimur per Christum dominum.  \pend
\phantomsection
\addcontentsline{toc}{subsection}{\textit{Regi autem saeculorum immortali, inuisibili soli deo honor et gloria in saecula saeculorum. Amen. }}
\subsection*{\textit{Regi autem saeculorum immortali, inuisibili soli deo honor et gloria in saecula saeculorum. Amen. }}\pstart Continuo et maxime Paulus memoriam facit gratiae deiad se persecutorem ueritatis, quam gratiam Christus dominus illi largitus est, qui post tam horrendam piis omnibus persecutionem, Christo in membris suis illatam, tantam rependit misericordiam, longanimitatem, patientiam, et affluentiam donorum spiritualium, sibi Paulo cui inimicus esset delatam, et mirabilem uocationem ad apostolatum, ut exemplar fieret peccatoribus, fide in Christum saluandis. Qua permotus memoria beneficii tanti, cum gratitudine et admiratione tantae gratiae susceptae exclamat, beneficium agnoscens, et gratiam reddens, dicit, Regi saeculorum etc. Rex autem saeculorum, omnipotens deus dicitur, qui coelum condidit, et omnsa coelo conclusa, sapientia sua gubernat omnia, et regit clarissime agnita, qui in omnia saecula semper immutabilis perseuerat. Ne Solem, stellas et coelos reges nostros credamus uel suspicemur, sed qui haec omnia fecit in gloriam suam, quae et suo nutui subduntur et obediunt: in cuius manu sunt non solum corda, sed et uoluntates, animae, et uita regum omnium, qui sine eo nec momentum uiuunt, nec uel minimum possunt. Est autem Saeculum in scripturis, spatium diuturnum, uel uitae ut plurimum humanae, uel centum annorum, saepe etiam quinquaginta: item, anni restantes ad quinquagesimum iubilaei secundum legem, etiamsi pauci sint. Hic autem saeculorum nomen accipiendum longaeuitas omnium uiuentium. Rex enim estcreator omnium, semper uiuens, quantumlibet sequantur uicissim hominum generationes et centenarii, uel aerae quaecunque . Nos autem ut aqua dilabimur, ipse uero fluuius est perennis, iugiter perseuerans, uiuificans creaturas, infinitus et immobilis existens, dat cuncta moueri. Immortali) ait, Qui non so lum immutabilis est, et nihil perdere potens, semper perfectus et felix, sed omnibus quoque uitam donans, in quo uiuimus, mouemur et sumus. Non ut dii falsi, qui olim homines fuerunt mortales, sed summum bonum, ideoque ; sine summo consistit malo, quod habetur mors, et nihilum uel corruptio. Neque etiam in Christo mortua est, nec mori potuit diuinitas, sed tantum caro destituta ab anima triduo, ambo diuinitati inseparabiliter iuncta. Hic autem Paulus ipsissimam intendit diuinitatem, quae aliud non est quam pater, filius et spiritus sanctus, unus deus, et una deitas, infinita perfectione et duratione. Inuisibili.) Patrem enim et diuinitatem nemo uidit unquam mortalium, nedum oculis, sed nec perfecte intellectu quocunque creato, etiam angelico. Filius autem qui est in sinu patris, ipse perfectissime naturaliter nouit, comprehendit, ac eandem reuelat, cui et quantum uult. Est enim deitatis natura speculum, non naturale, sed uoluntarium, si uultuidetur, et quantum uult, a nemine tamen potest perfecte uideri, quod finita capacitas, infinitum bonum perfecte complecti non possit uel capere. Sed nec unquam alicui uisus est pater. Omnia autem agit per filium, per quem fecit et saecula, quo uerbo suo omnia administrat. Et potestas omnium in coelo et in terra data est Christo dei filio. Soli deo honor et gloria.) A quo omne bonum, ideo ei debetur honor, qui est bonitatis testimonium. Sicut itaque regi mortali honor dependendus, non regi, sed deo referendus est, cuius rex est minister, a quo et regia dona suscepit, et quem praecipit honorari. Sic pro sanctitate, uirtute, sapientia in hominibus quibuscunque   \pend
\section*{COMMENT. IN I. EPIST. }
\marginpar{[ p.478 ]}\pstart honor omnis deo debetur, qui solus largitur haecomnia, ut non tam illi uel illa, quam deus ipse optimus et maximus in illis honoretur. Alioqui creaturae tribuentes, quod solius est dei, idololatriam committimus, et utendis fruimur. Vthodie ui, demus honorari sanctos, a quibus petimus quae solus deus donare potest, et qui a se talia docuit postulare. Et sic deo propria, illis contra illorum uota et studia com munscamus. Soli deoomnis sit honor, a quo solo omne bonñ. Dum uero dicitur, Soli deo, non excluditur filius et spiritus sanctus, quin et illi sunt unus cum patre, immortalis et inuisibilis dominus atque deus. Pater autem omnium perfectionum diuinarum fons est et abyssus, emanationumque diuinarum origo. Quae licet deitas eadem sit in tribus, non tamen eodem modohabendi: omnia enim a patre sunt, et ipse a nullo, filius a patre, spiritus sanctus a patre et filio, pater autem a nullo. Est autem gloria dei, recognitio et attributio omnis boni, quia dominus uniuersorum sit ipse, cuius bonitas relucet in rebus omnibus, qui ab omni creatura colendus cognoscitur, et omnia sapientissime moderatur, cui omnia seruiunt. Haec omnia gloria dei sunt, et clara infinitae et immensae bonitatis agnitio ac professio. Quod additur, In saecula saeculorum, infinitatem durationis ac aeternitatem signi ficat carentis principio et fine, cuius omnes temporis differentiae simul sunt, apud quem non est uicissitudinis cuiuscunque obumbratio. Amen.) Attestationis no, ta est, et fidei omnium quae dicta sunt proxime de diuinitatis excellentia et gloria,  \pend
\phantomsection
\addcontentsline{toc}{subsection}{\textit{Hocpraeceptum commendo tibi fili Timothee iuxta prophetias quae de te praecesserunt, ut milites in eis bonam militiam, habens fidem et bonam conscientiam. }}
\subsection*{\textit{Hocpraeceptum commendo tibi fili Timothee iuxta prophetias quae de te praecesserunt, ut milites in eis bonam militiam, habens fidem et bonam conscientiam. }}\pstart Vt intelligas fili Timothee, quid maxime te agere uellm, quidque te deceat, et exigat a te dominus, qui et prophetiis quibusdam de te reuelatis, praemonstrauit te dignum futurum episcopum, et sancti spiritus admonitione, id officii tibi iniunctum est, quod non ego impono tibi, sed quod a spiritu Christi et uerbo dei tibi iniungitur, hoc tibi commendo, commonefacioque istius praecepti, ut memine ris solicite, militare bonam militiam contra infideles pseudapostolos, criminosos Christianos et falsos fratres, eos ut magnifice superes et uincas uera fide Christi, bonaque conscientia. Ex quibus intelligimus fuisse Timotheum optimae famae, et sanctae conuersationis uirum, eruditum in sacris literis, de quo lam ante praedixerant quidam prophetae, quorum in primitiua ecclesia multi erant, futurum aliquando Ti motheum idoneum episcopum, et orthodoxae fidei catholicaeque doctrinae propugnatorem egregium. Quod et ipsum fortassis Paulus cognouerat in spiritu, del reuelatione commonitus, ob idque syncerius dilectum episcopum ordinauerit in urbe celeberrima. Ephesi autem tunc erant aduersarii multi euangelicae doctrinae, uariae conditionis, gentiles in primis, qui pro Diana sua, ob quaestum zelabant pro sua idololatria defensanda. Iudaei quoque Paulo eiusque sociis et dogmatibus exitium iugiter meditabantur. Pseudapostoli legem Mosis cum euangelio seruandum do centes. Falsi denique fratres, qui infirma fidelium traducebant, persecutiones susci tabant, inordinateque uiuentes, scandalo erant sanctae ecclesiae. Hos omnes non po terat Timotheus aliis armis debellare quam fide et conscientia, fortiter et constanter adhaerendo sacris literis, et uerbis sacrae doctrinae, ex lege, prophetis et hagiogra phis (quae tunc fere solae extabant) quibus contra aduersarios omnes illos, uerbo scilicet ueritatis, et animi fortitudine, firma quoque fide, constanter defenduntur, et sine trepidatione docentur, catholica et apostolica primitiuae ecclesiae dogmata. Interim sibi testimoniumprae bente et audaciam conscientia sua, cum immacu.  \pend
\section*{AD TIMOTH. CAP. I. }
\marginpar{[ p.479 ]}\pstart lata ac religiosa conuersatione, coram omni ecclesia, quae eum audacem reddere poterat, ad animose uiriliterque occurrendum, etiam apostolatus sui authoritate, et spiritus sancti uirtute, qua perfusus erat, omnibus illis aduersariis potestatibus, doctrina et moribus, uolentibus corrumpere commissam sibi ab Apostolo ecclesiam Ephesiorum, exemplum futurus episcopis successoribus.  \pend
\phantomsection
\addcontentsline{toc}{subsection}{\textit{Quam quidam repellentes circa fidem naufragauerunt, ex quibus est Hymenaeus et Alexander, quos tradidi fatanae, ut discant non blasphemare. }}
\subsection*{\textit{Quam quidam repellentes circa fidem naufragauerunt, ex quibus est Hymenaeus et Alexander, quos tradidi fatanae, ut discant non blasphemare. }}\pstart Erant Ephesi, qui contra conscientiam et agnitam ueritatem Apostolum traducebant, confingebant in eum mendacia, et obloquebantur: qui et semel susceptam euangelii fidem abiiciebant, Hymenaeus et Alexander. Hii uerbis Apostoli pertinaciter restiterant, seditiones mouerant contra Paulum, Iudaei ex genere, et a fide apostatae, Pauli persecutores subdoli, et non semper in manifesto aduersantes. Vnde et Paulus Timotheum admonet de Hymenaeo, ut caueat, nihil illi fidat, in secunda epistola. Satanae tradere, est excommunicare, et eiicere a coetu fidelium, ut a morbidis agnis sibi sani et boni Christiani caueant et non inficiantur. Prodidit ergo illorum malitias, et ecclesiae patefecit, ut et illi resipiscerent, et reliqui sibi de periculo prospicerent. Talis olim in ecclesia excommunicatio agebatur, sic publicorum criminum propalatio, et exomologesis in aedificationem ecclesiae fiebat. Quas utinam consuetudines probatissimas, iusta forma, secundum apostolicas traditiones ecclesia recuperet, ut reformentur mores, et disciplina non cesset, ut discant non blasphemare. Blasphemia nimirum est, peruersa docere, secreta diuina temeraria curiositate disquirere. Nihil enim commune habet humana ratio collata diuinis. Compescendi diligenter in ecrlesia sunt, qui dominum blasphemant, uerbi illius authoritati detrahunt, qui gentilium errores postliminio reducunt, qui Iudaicis fabulis et haereticorum uenenis inficiunt ecclesiam, qui religionis praetextu gentilium ritus commutant in speciem religionis, quos satius fuisset in totum reiici, ne gentilium ceremoniis ecclesia plus quam Iudaica grauetur, puritasque doctrinae et operum commaculetur, et tandem opprimatueritatem.  \pend
\endnumbering\beginnumbering\section{CAPVT II.}
\phantomsection
\addcontentsline{toc}{subsection}{\textit{\huge\textbf{O}\normalsize Bsecro igitur primum omnium fieri obsecrationes, orationes, postulationes, gratiarumactiones pro omnibus hominibus. }}
\subsection*{\textit{\huge\textbf{O}\normalsize Bsecro igitur primum omnium fieri obsecrationes, orationes, postulationes, gratiarumactiones pro omnibus hominibus. }}\pstart Cum igitur necesse habeas militare periculosam militiam, cumque fide tibi et conscientia opus sit, ut imminentibus undique et semper malis occurrere possis, et officio satisfacere tuo, spiritu sancto te uocante commisso, quae sine dei gratia ac uirtute, nec a te praelato haberi possunt, neque ab alio. Ideo adhortor, non impero, non importunus urgeo, ut dominus, sed ut fidelis tibi pater filium obsecro, quem pro syncero affectu et filiali amore, scio non minus complecti hortamentum, quam seruus imperium. Non enim patris praeceptum filius expectare debet et exigere, sed solo nutu insinuasse filio beneplacitum patris, abunde suffecerit, ut protinus perficiat sedulo, quod patrem optare crediderit. Sic non mandatis, imperiis et coactionibus in dei negotiis est agendum. Nec Christus nec apostoli id tentarunt, sed usi sunt admonitionibus, consiliis, orationibus, et adhortationibus, quibus in ecclesia proficere studuerunt. Vt nihil mandauerint, nisi quae pro charitate seruanda seriosius proponebant. Vt ante omnia fiant orationes.) Cum sanctum ac deo acceptum, uel probandum, nihil prorsus ualeamus, nostris tantum freti uiribus, sine dei auxilio, nec uerba, nec opera, nec exempla, nec aliqua solicitudo nostra aliquid prosit uel promoueat, necesse est ante omnia  \pend
\section*{COMMENT. IN I. EPIST. }
\marginpar{[ p.480 ]}\pstart nostra conamina, dominum implorare, auxilium illius petere ac assistentiam, ut ipse bonum inspiret, desiderari faciat ac perfici: nihil enim aliunde obtinebimus.  \pend\pstart Deprecationes.) Dicuntur heae orationes ad deum, quibus aduersa quae merui mus peccatis nostris, per dei misericordiam amoueri precamur, iram scilicet dei, et poenam condignam, conscientiae angustias, incommoda corporis et spiritus, tentationes perniciosas, et proximas periculo, ac casus peccatorum. Obsecrationes.) Sunt, quibus per omnia sacra et deo grata, consequi desideramus dona dei nobis necessaria, spiritualia in primis, deinde et corporalia. Vt quod nostris meritis obtinere non possumus, eius gratia, et Christi domini meritis et intercessione a domino postulemus nobis concedi in gloriam dei. Interpellationes.) Sunt cum unus pro alterius necessitate ex charitate solicitus intercedit, et profratre orat, ut uel aduersa amoliantur, uel optabilia consequantur, cum scilicet pro inuicemoramus, ut pariter saluemur. Gratiarumactiones.) Sunt, eum pro acceptis beneficiis dominum laudamus, et omnia eidem accepta reddimus, et a nobisipsis nihil nos habuisse uel potuisse fatemur. Vicissim nos offerentes uoluntati suae, cuius beneplacito per omnia subiiei desideramus. Hoc est enim eucharistiae uerum ac Christianissimum sacrificium, laudis scilicet et confessionis, ac gratiarumactionis. Pro omnibus hominibus.) Non solum orandum pro amicis, fidelibus, iustis, ac familiaribus, ac benefactoribus, uerum omnium maxime proeis qui magis in digent, dei maiori opus habent misericordia. Vt sunt inimici nostri, ut charitate imbuantur Christiana, infideles, Iudaei, haeretici, ut respectu gratiae diuinae, lumine uerae fidei et ueritatis illustrentur, et domino credant, iniusti ac impii ut resipiscant: ignoti nobis, ut domino miserante a peccatis liberentur. Nihil unquam inter orandum occurrat discriminis, quo minus affectuose pro uno hominum genere oremus, quam pro alio, uel excludere quoscunque praesumamus, uel intendamus frigidius. Nec tantum oremus obbeneficia, ut hypocritae solent, et incommodum nostrum, sed sola fraterna charitate ad orandum permoueamur pro egentibus: et quia id domino placere credimus, similia abaliis orationum suffragia desiderantes. Non iubet Apostolus orari pro defunctis, sed pro hominibus, non pro animabus fidelium mortuorum, quae non iam nobis proximae sunt, sed deo. Viui no bis ubique commendantur, mortui deo:nos contra egimus, totam fere nostram solicitudinem impendentes in eos, demortuos scilicet, de quibus nobis dominus nihil praecepit. Pro his non piget nimium multa dare, agere, pati, interim tamen fratres indigentes, maximopere a deo nobis commissos, dissimulamus, et nihil iuuamus. Sed haec excogitauit auaritia cleri, ut morituri multa legarent pro se mortuis de diuitiis suis, quas posthac, ut non suas, relinquere cogebantur. Haeredibus debita iure naturae ea intercessoribus pro animabus facile largiri poterant, maxime quia purgatorii poenas his mediis euadere posse instruebantur praeter uerbum dei, et sine sanctorum exemplo.  \pend
\phantomsection
\addcontentsline{toc}{subsection}{\textit{Pro regibus, et omnibus qui in sublimitate constituti sunt, ut quietam et tranquillam uitam agamus, in omni pietate et castitate. }}
\subsection*{\textit{Pro regibus, et omnibus qui in sublimitate constituti sunt, ut quietam et tranquillam uitam agamus, in omni pietate et castitate. }}\pstart Apostoli temporibus reges mundi infideles erant et ethnici, aduersarii fidelibus, sed tamen pro eis orari uoluit, imo tanto solicitius instetit, quanto magis clementiae diuinae influxus illis pro illorum regimine necessarius erat, ut boni fierent, ut tractabiles, mansuete praesint, benigne dominentur, quo ualeant sustineri, ut cor illorum regatur a domino, quo subditi tranquillius regantur: ut illis desuper praestetur sapientia et timor dei, ut ament iustitiam, et tueantur innocen\pend
\section*{AD TIMOTH. CAP. II. }
\marginpar{[ p.481 ]}\pstart tiam. Et omnibus in sublimitate constitutis.) Addita coniunctio expositiue accipitur. Regibus, id est sublimibus. Sunt em reges qui in magistratu aliquo constituti, regunt, moderantur et praesunt. Constituri uero non ab alio que ̃ a deo. Omnis enim potestas a deo est, etiam mala, propter peccata populi, et imminens subditis dei iudicium. Per eos enim deus non raro sua exequitur iudicia. Dabo, inquit, tibiregem in furore meo. Effoeminati erunt principes eorum, etc. Maximam itaque  nobis poenam deprecabimur, si bonum et religiosum principem orationibus a deo promerebimur. Propter peccata enim populi dominus regnare facit hypocritam, id est impium. Sublimitas autem est, regere subiectos, iuxta stilum spiritus sancti, et a deo ordinata: non tantum ut honorentur praefecti, sed multo magis ut onerentur, solicitudine maiori prodesse cupientes eis, quibus ut per charitatem et sapientiam seruiant, a deo dati sunt. Reges gentium dominantur, et benefici dicuntur, non sic reges fidelium. Hi enim sciunt sese tanto modestiores gerere debere, quanto sunt plurium commodis et profectibus destinati, et obsequio charitatis ac fidei deuincti. Debetur ipsis honor, reuerentia, timor, obedientia, ac sustentatio a subditis: ipsi uicissim obligantur amorem, fidem, curam, tutelam, solicitudinem, instructionem, consolationem, regimen, coërcionem a malis, et discolis poenam. Nec obmittendum negligentius, quod non prece, pretio et instantia, sublimitas sit ambienda, uel regimen aliorum, sed constitui oportet ab aliis, nihil nisi subditorum bonum respicientibus, neque sine consensu eorum, si tamen rationabiliter sibi consultum cupiant, et non omnem malint subterfugere disciplinam, ut non tam seruitus habeatur, que ̃ administratio disciplinae. Vt placidam et quietam uitam agamus.) Annitendum omni studio regibus, ut populus beneuolenter et cum gaudio, non ex tristitia obediat, sed placide, syncero affectu, sine murmure et displicentia, si fieri possit, et sine uiolentia:ut sint omnia sine tumultu, odio, seditione et rebellione. Sed Christiana pace, tranquille conuersentur, ut fratres ac domestici discipuli Christi, pacifici, modesti, mansueti et humiles. Haec est enim uita felicissima, planeque beata, et Christiana charitate iucunda. Cum omni pietate et honestate.) Faciunt haec duo perfectam quietem, si scilicet deum nostrum colamus fide, spe et charitate, ut totos nos committamus ei ut patri: et omnia quae occurrunt, patris prouidentia et pietate nobis euenire certo credamus. Tunc enim murmurabimus nunquam, temere non iudicabimus illius iudicia et dispensarionem, sedut sunt omnia optima cordis pace, ut patris pientissimi beneficia amplectemur. Castitas hic pro omni morum grauitate accipienda, ut honestas. In conuersatione enim ad proximum, morum honestas et comitas admodum nos commendat fratribus. Cauebit quisquis huiusmodi est, offensam et odia, perpetuam autem tranquillitatem fouebit et amicitiam, quae inter inhonestos ac morosos aut leues personas diu constare non potest.  \pend
\phantomsection
\addcontentsline{toc}{subsection}{\textit{Hoc enim bonum est, et acceptum coram saluatore nostro deo, qui omnes homines uult saluos fieri, et ad agnitionem ueritatis uenire. }}
\subsection*{\textit{Hoc enim bonum est, et acceptum coram saluatore nostro deo, qui omnes homines uult saluos fieri, et ad agnitionem ueritatis uenire. }}\pstart Quieta conuersatio nostra, et tranquillitas huius uitae in simplici obedientia, etiam ethnicis regibus exhibenda docetur, ut ecclesia careat dominantium mole stia, neque grassentur, ut magis opprimant, ne irritentur contra nos et euangelicam doctrinam, non scandalizentur nostris moribus, non ansam arripiant mouendi persequutiones, et blasphemandi nomen Christi. Cedit enim in gloriam religionis Christianae, imo domini dei nostri, qui per filium suum saluauit nos, si pacifice, mansuete ac sedulo dominis obsequamur, subiecti omnibus in timore del, osten\pend
\section*{COMMENT. IN I. EPIST. }
\marginpar{[ p.482 ]}\pstart dentes facto, optimi dei optimas leges scriptas esse in cordibus nostris. Vt ex bonis operilbus nos considerantes, glorificent deum, si quando respectu misericordiae a deo fuerint uisitati, et corde pio religiosoque donati, dominumque deum nostrum ex fidelium religione cognouerint, nobiscumque ; uenerari inceperint, orationibus nostris interim domino commendati, et exemplo acquisiti. Omnes, inquit, homines uult saluos fieri.) Est enim dominus deus summe bonus, nihil eorum odiens quae condidit, neminem uult perire. Adest sua gratia omnibus, ut ueniant ad agnitionem ueritatis, quam per propheticas euangelicas et apostolicas literas, communem fecit omnibus. Vbi denuo coniunctio. Et exposstiue sumenda uidetur, ut nihil aliud sit uelle saluos fieri, que ad agnitionem diuinae uoluntatis perducere, ut agnoscere ueritatem sit fieri saluum. Ipsa enim ueritas liberat a peccato, si illi credatur. Denique uocauitomnes, annunciatur omnibus, Christus pro omnibus mortuus est, omnes ad se trahit beneficiis, proomnibus tradidit semetipsum. Contemptores autem dei et negatores non computandi ueniunt inter omnes, sed redigentur in nihilum coram deo. Si itaque sermoincidat, Quare non omnes beatitudinem assequantur, leum omnes homines uelit saluos fieri ?Responderi potest, Quare non omnes capiant gratiam, quam sibi porrigi nemo ambigit? Quare non credant uerbo dei omnes? Quare nemo etiam cogi uelit ad bonum suum, nec debeat ? Nec se quispiam excusare potest de ignorantia, deque tractu instinctuum bonorum in conscientia, qui non cessant. Quod si gratiam sequendi quis abesse causetur, non tamen eam optare ac petere se non posse uere dicet. Domino autem dicente, Petite et accipietis. Is re uera nonnihil recipiet, qui sedulo dominum fuerit deprecatus. Vltima tamen responsionis anchora tandem erit, uoluntas dei, quae regula est iustitiae, cuius non sit inquirenda ratio, qui et nobis nihil obligatur, et bonorum nostrorum non eget. Voluntate signi et antecedente, uult omnes saluari, dicunt Theologi post Augustinum, et distributione accommoda euadunt. Denique agnitio ueritatis, sensus est Christi, qui ueritas est, qui euangelium praedicauit mundo, qui ueram fidem ac uiuam, portum salutis statuit, et iustitiae summam. Lux mundi est, qui illuminat omnem hominem. Humanae scientiae, et ueritas intellectus naturalis, ad salutem inutiles sunt, et mendacium dici possunt. Christum autem scire, sanctorum est, ideo stultus in hoc mundofiat, qui uoluerit fieri sapiens in deo et saluus.  \pend
\phantomsection
\addcontentsline{toc}{subsection}{\textit{Vnus enim deus, unus et mediator dei et hominum, homo Christus lesus, qui dedit redemptionem semetipsum pro omnibus. }}
\subsection*{\textit{Vnus enim deus, unus et mediator dei et hominum, homo Christus lesus, qui dedit redemptionem semetipsum pro omnibus. }}\pstart Breuis hic expressa est summa fidei Christianae, et ueritatis irrefragabilis, cuius solius cognitio firmiter persuasa, salutem operatur hominibus. Si, inquam, non simplici id notitia sciant, sed potius uera sapientia: qui scilicet sic norunt, ut sapiant q̃ suauis sit dominus. Hunc ergo tenorem fidei, et breuissimam regulam, simpliciter et breuibus discutiamus. Vnum esse deum religiose cognoscitur, quum praeter eum qui omnia condidit, nihil aliud, ut alicubi summum bonum desideratur, colitur, amatur, antefertur, uel confertur nihil. Non scilicet uenter, non caro et sanguis, aurum, honor, laus, et totus hic mundus. Omnia haec postponuntur, non tantum uerbo et lingua, ut facile possunt hypocritae, sed corde et sensu. Dum illud unum semper intendimus et attendimus, propter illud omnia agimus, non circa plurima occupamur et distrahimur. Hoc est uere agnoscere unum deum, praeter eum nihil desiderare, quod non propter illum uelimus. Vnus mediator.) uel conciliator. Secundus fidei huius breuissimae articulus est, credere et scire  \pend
\section*{AD TIMOTH. CAP. II. }
\marginpar{[ p.483 ]}\pstart Christum esse unicum mediatorem, quem pater coelestis statuit inter se et nos,ut cum essemus peccatores, et a deo proiecti et inimici, per Christum reconciliemur. Hic omnes homines peccatores esse, et reconciliatione egere ad deum, manifeste docetur, qui ut aduersarii et hostes, deum offensum sibi patiuntur, et abiectos esse damnatosque ex primi hominis transgressione significantur, sola Christi gratia conciliandos: et non id solum, sed et filiationis et haereditatis ius accepisse per Christi mysterium, et in Christo omnia bona, ut nostra sint, quae sunt fratris nostri. Vnus autem, non multi mediatores uel conciliatores docentur: unus enim Christus, unus saluator. Non ergo praeter eum interponendi sunt alii, quantumlibet sancti et pii homines, omnes enim peccatores sunt, et gratia dei indigent. Solus Christus suapte innocentia, dignitate et meritis proprie dictis, non ex gratia, nos peccatores conciliat patri, pro officio suo quod a patre accepit. Omnia denique pater dedit in manus filio unigenito, ad illum nobis currendum, quoties quomodo cunque urgemur, qui fons est pietatis, et fidelissimus frater, qui sese pro nobis totum obtulit, in acerbissimam simul et probrosissimam mortem crucis, non nisi ex maxima charitate et fide ad genus humanum. Non est ergo ut eum ueluti morosum nobis promereamur, sanctorum aliorum quorumcunque intercessionibus, benignior nobis est, quam omnes sancti esse possint, quantumlibet charitate fulgentes. Solus pro nobis mortuus est, et nos toties inuitat, pulsat, expectat, hortatur, allicit, rogat, et cogit aliquando. Venite, ait, ad me omnes etc. Venite edite absque argento etc. Audit nos semper per suam sibi soli propriam omniscientiam, adest per immensitatem, qua patri consubstantialis est: intendit singulis perfectissima charitate, quandoquidem sui sunt omnes: non distrahitur multitudine curandorum uel saluandorum, omnia nostra nouit melius quam nosipsi, quibus egeamus considerat, nobis dormientibus. Praecepit seriose, ut in solo suo nomine precaremur patrem. Honorem hunc gloriosissimum suum alteri prorsus nulli cedere potest: nec patitur communem ab aliis usurpari unigenitus filius. Homo, inquit, Christus Iesus.) Mirum quomodo dubitauerunt hominem non fuisse quidam haereticorum, qui sacras literas non negant, et hanc epistolam, qui hic simpliciter homo dicitur. Caro enim et frater noster est, tentatus per omnia, absque peccato, et pro nobis hominum nouissimus sese exhibuit, ac uermem nominauit, ut hominem perditum redimeret. Et primogenitus quoque dicitur ex multis fratribus, ut chirographum peccati deleret per crucem, uere pro nobis mortuus est, utque homo esset, qui pro humano satisfaceret genere. Christus Iesus, unctione sua consecrans omnes gentes, et Christianos faciens sua gratia, et spiritus sui unctione: passione quoque sua saluans gentiles et Iudaeos, faciens utraque unum. Qui dedit semetipsum redenptionem.) Obtulit seipsum in remissionem omnium peccatorum, sacrificium seipsum perfectum et maximum, omnium quoque praecedentium sacrificiorum finem et scopum, proque nobis captiuis commutationem sufficientissimam. Non raptus est ad iudicium mortis inuitus, sed sua sponte obediuit, et seipsum obtulit pretium deo patri, pro salute humani generis, et reconciliatione perfecta. Cum essemus inimici, et tales inimici, ut etiam eum nos ipsum crucifixerimus, homines scilicet, pro quibus se pretium et satisfactionem mediator obtulit. Nec dicamus, non esse nos posteros Iudaeorum, qui hodie quoque in glorificatum dominum nostrum et regem simile facinus non refugeremus, simili occasione permoti, qua Iudaeorum primores, quibus meliores inter nos non ita multos inuenire licebit, quin et acerbius ipsum iterum nostris uitiis crucifigimus. Pro omnibus.) Exposuit ac tradidit se in mortem pro omibus, nemo nisi nolens saluari, relinquitur. Sufficiens  \pend
\section*{COMMENT. IN I. EPIST. }
\marginpar{[ p.484 ]}\pstart fuit, et satisfecit deo patri pro omibus, nihil obmisit de suo, quo minus omnes redimi potuerint, et hodie possint. Excusatio nulla nobis relinquitur, qui plane scimus quid uel facere, uel desiderare, uel quae obmissa dolere, uel non indolentiam deprecari, et humiliter agnoscere debeamus, proque fide huiusmodi et charitate deum orare. Quis haec non se posse nisi contumax uere fatebitur ? Ad hanc uerita tis agnitionem, si quis illuminante deo peruenerit, ut agnoscat ac credat certa fide, unum suum deum, et conciliatorem unicum, et sibi deuotissimum dominum Iesum Christum, quodque seipsum dederit ex maxima gratia redemptionis pretium pro omnibus credentibus, sine dubio saluus erit, perfectus fidelis habebitur, et Christianus coram deo et Christo suo. Quae illi fides reputabitur ad iustitiam, sicut Abrahae. Hanc breuissimam summam totius noui testamenti et euangelicae doctrinae, dum olim apostoli praedicarent facile orbis attrahi ad credendum poterat, idolisque abiectis, Christiana religione imbui. Cum nondum tam incredibilia obtruderentur multa, qualia interim praeter sacram scripturam fidelibus, de plus que ̃ bis duodecim articulis neotericae fidei imponuntur necessario credenda, de Papatu multa, de indulgentiis, purgatorio, diuis, et diuorum intercessionibus, miraculis, delectu ciborum, excommunicationibus, confessionibus, casibus episcopalibus, uotis, peregrinationibus, Conciliis, sacramentis, impedimentis, Missis, horis canonicis, ceremoniis, festis, ordinibus, imaginibus, monasticis regulis, dispensationibus, fraternitatibus, consecrationibus. Et quis ea numerare possit, quibus hodie grauantur conscientiae piorum, et impiorum conuersio non iuuatur? Quorum nihil continetur in Symbolo apostolorum uel canonicis scripturis, imo nec priscis sacrosanctis sanctorum patrum conciliis: de quibus multa scribi possent, sed sufficiunt quae nunc passim leguntur. Nobis abunde iamdudum suffecit uerissimum diui Hieronymi dictum, Quicquid (credendorum et necessario agendorum) de sacris scripturis non producitur, eadem facilitate contemnitur qua probatur.  \pend
\phantomsection
\addcontentsline{toc}{subsection}{\textit{Cuius testimonium temporibus suis confirmatum est, in quo positus sum ego praedicator et apostolus. Veritatem dico, non mentior, doctor gentium in fide et ueritate. }}
\subsection*{\textit{Cuius testimonium temporibus suis confirmatum est, in quo positus sum ego praedicator et apostolus. Veritatem dico, non mentior, doctor gentium in fide et ueritate. }}\pstart Ad hoc Christus ueniens sese patri pro nobis obtulit, ut suo illo tempore toti mundo testaretur diuinae charitatis erga genus humanum beneficium, dudum promissum patriarchis ac prophetis, in redemptionem omnium hominum, ut credentes per eum saluentur. Noluit enim occultum esse, quod sacrificiorum figuris et propheticis promissionibus in ueteri testamento uoluiresse tam multis testatum. Quin et ego Paulus ad hoc nostrae redemptionis euangelium praedicandum assumptus sum ab eodem Christo domino, segregatus ac electus post alios, ut toto orbe gentibus hoc praedicem, et hanc dei tantam gratiam adnunciem. Nuncius missus a Christo et apostolus, quanquam cum Christo non uixerim, ex misericordia tamen sua me conuertit, ex persecutore et blasphemo in apostolum ac ueritatis doctorem ad omnes gentes. Veritatem dico, non mentior.) Modestissime iurat, quod notorium esset, et negari non poterat. Celebris, inquit, fama iactatur insaniae meae, cum spirans minarum et caedis in discipulos, subito in multorum comitatu e coelo uocatus sum, et a Christo conuersus. Deinde cum Barnaba a spiritu sancto segregatus in gentes, attestatione non modica miraculorum, profectu praedicationis maximo, a Iudaea incipiens per Asiam, usque ad Illyricum omnia repleui euangelio Christi. Non potest esse dubia apostolatus mei uocatio, et quod me Christus  \pend
\section*{AD TIMOTH. CAP. II. }
\marginpar{[ p.485 ]}\pstart gentium doctorem dederit, testantur signa apostolatus mei: qualem denique me gesserim in gloriam Christi, non in meum compendium aliquod temporale, nihil miscendo doctrinae, quam in me Christus non fuerit loquutus. Omnia sine hypocrisi et figmento, dextre ac plane ueraciter tradidi, ut accepi a domino. Doctorem se nominat, non in Vniuersitate, ut loquuntur, promotum: non enim Aristotelem audierat antea, et parua logicalia, sed spiritu sancto promotus, docebat sine inter missione uerbo, exemplo et scripturis, non tam nomine doctor q̃officio, assiduoque ́ docendi studio. Tales etiam non sic promoti, doctores erant Origenes, Auguslinus, Hieronymus, et similes. Quanquam eruditionis et fidei testimonium habaere, et factis consequi non fuerit inutile, sed sine quaestu, fastu et hypocrisi: utque professionis suae specimen ostendant studio, fide et exemplis, uerbo et facto, non serico uel annulis, ne inter comessationes et pocula pro gratia nummorum, et personarum respectu, officium et authoritas docendi in ecclesia commendetur indignis, cum fide et ueritate doctores agant, ut scripturis canonicis edocti, fidem fideliter doceant, sine permixtione falsorum dogmatum, quae resipiant sensus car, nis, ambitionem honoris, et quaesitum.  \pend
\phantomsection
\addcontentsline{toc}{subsection}{\textit{Volo ergo uiros orare in omni loco, leuantes puras manus, sine ira et disceptatione. }}
\subsection*{\textit{Volo ergo uiros orare in omni loco, leuantes puras manus, sine ira et disceptatione. }}\pstart Docuerat orandum pro consequenda uitae tranquillitate in hoc mundo, et ut ad agnitionem ueritatis omnes peruenirent saluandi. Nunc uiros docet ac mulieres instruit, quomodo orandum, et de aliis necessariis, incidenter et obiter, foeminarum mores instituendo: primum autem uiros, Quos, inquit, uolo orare in omni loco. Nondum apostolorum temporibus Christianis erant amplissima templa, sufficiebat locus profamilia quacunque , longo tempore in criptis praedicatum: ubicunque  ecclesia erat, duo scilicet aut tres uel plures, ibi locus orationis et praedicationis, in cubiculis, in agro, in naui, in montibus ac syluis: ubique orationis locus, ubi spiritus sancti templum consecratum habetur, quod sunt fidelium animi. Idipsum docuit ueritas Christus, Cum, inquit, oraueris, ingredere cubiculum, et ora patrem in abscondito. Multi hodie asueuerunt, ut non nisi in ecclesia orare ualeant, et exaudiri credant: si tamen uere orant, et non potius multa sensuum et mentis distractione adnumerent potius uerba, que ̃ ut corde, affectu feruoreque spiritus subleuentur in deum, in considerationem fragilitatum, defectuum ac miseriarum, quarum consideratione legitima oratio sit promouenda. In omni, inquit, loco.) Consecratus est a deo mundus totus, et homo orationi. Vt non in templo Hierosolymis, nec Romae, sed in spiritu et ueritate sciamus orandum. Id quod Christus docuit dominus, quique se sic orantes exauditurum promisit. Caetera diuinae potius prouidentiae committantur. Leuantes puras manus.) Non lotas, Pharisaeorum more, sed puras ab omni uiolentia, a sanguine innocentum, a munerum acceptione, usura, furto, a foedis libidinosisque contactibus, et ab omi malo opere. Nec uult Paulus extendi manus pharisaico more, sed desiderantium ex petitionis ardore. Non obsunt, sed prosunt quaedam ceremoniae orantibus, sed non nimiae, et quae mentis distractionem uel hypocrisim mouere possint. Sine ira et disceptatione.) Vbi docet nihil a deo postulandum, nisi prius orans proximo conciliatus sit charitate Christiana. Solet uirorum animos occupare furor et ira, ut superbia mulieres, incontinentia ac liuor. Ira autem charitati contraria, orationi locum omnino praecludit. Non ignoscit deus nobis, nec dona pietatis largitur, nisi peccantibus in nos antea indulserimus, nisi fratri offensam remiserimus, et ex animo dilexe\pend
\section*{COMMENT. IN I. EPIST. }
\marginpar{[ p.486 ]}\pstart rimus propter deum.  \pend
\phantomsection
\addcontentsline{toc}{subsection}{\textit{Similiter et mulieres in habitu ornato, cum uerecundia et sobrietate ornantes se, et non in tortis crinibus, aut auro, aut margaritis, uel ueste pretiosa, sed quod decet mulieres, promittentes pietatem per opera bona. }}
\subsection*{\textit{Similiter et mulieres in habitu ornato, cum uerecundia et sobrietate ornantes se, et non in tortis crinibus, aut auro, aut margaritis, uel ueste pretiosa, sed quod decet mulieres, promittentes pietatem per opera bona. }}\pstart Mulieres Apostolus docet post uiros, de idoneis earundem moribus, ne de forma superbiant amictu uel ornatu, quod agnatum mulieres patiuntur et gaudent, alioqui non pessimae, imo et bonae, castae ac prudentes, quae de speciei ornatuumque ́ cura uix sibi praecipere possunt. Admittit quidem Paulus, ut habitu ornato sint pro ingenio muliebri, non denique sordido amictu, uel nimium neglecto, quo abominatio sint maritis, et occasio forsitan exteras adamandi, sed modum praescribit, ut decet castas, uerecundas, et minime irritatrices libidinis. Munde quidem, eleganter ac honeste uestiantur, hoc est, ne nuditate et ornatu superfluo, pretioso uel curioso meretricibus assimilentur, appareantque affectu placendi in publico, maxime uero in coetu uel ecclesia, quod Christianis maritis nec placere potest neque debet. Ornentur autem legitimis monilibus, uerecundia scilicet et sobrietate, ex uultu, oculis, incessu, uerbis, subindicando moribus ipsis, inuitas sese offerri uirorum aspectibus, nolle ab aliquo desiderari ad malum, malle semper contegi, et aegre uidere uiros, multo molestius uideri. Quin et propter maritum uerecundia non nudentur: gratiosissimum enim ornamentum uerecundiam habet, quisquis est maritus sapiens et probus. Sobrietate.) Nulla uinolentia aut crapula aliquatenus notetur mulier, quae lubricitatis et libidinis certi sunt testes, teste Hieronymo. Quae et impedimenta praebent foetibus optatissimis, insipientiae et garrulitatis uitia secum adferunt, mendaciorum et litium matres, nihilque secretum seruaturae, seminaria malorum. Nascitur e sobrietate castitas, ut sunt sorores, continent inhonesta desideria, et libidines impediunt. Ornantes se, non in tortis crinibus.) Non compositione capillorum, et implicatura argenti et auri, non adulterino colore, uel cerussa fucoque peregrino, quibus iniuriam inurunt formae naturali, magis autem uirtutum exercitiis ornentur. Non auro aut margaritis.) Non onerent impensis maritos, colligant potius, et in thesauros reponant pro haeredibus, uel certe sanctius subueniant indigentibus, in opera pietatis uerae conseruent. Non addant superbiae naturali fomenta luxuriae, non certare foeminae uicissim optent de sumptibus, sed de modestia, sapientia et maturitate: pretium margarirarum et auri felicius expendetur in pauperes. Veste (inquit) non uestibus: quandoquidem superfluae uestes Christianos dedecent, et excusari non possunt, etiam in diuitibus. Non ad luxum amiciantur, sed ad honestatem et necessitatem. Vnica quidem non sufficit, sed superflua non decet, pretiositas intolerabilis. Scandalizat, irritat, zelum suscitat et exemplum, extenuat rem familiarem, difficile defenditur a tineis, et non sine negotio curatur ac solicitudine. Sed quod decetmulieres profitentes pietatem, aemulentur diligentius. Propensiores ipsas foeminas natura constituit ad cultum dei, ad multas ceremonias ac superstitiosas. Has discant non esse pietatem, frequentare ecclesias, gestare oratoria uel corollas, preculas mussitare, admirari cantus, et organis delectari, aspicere picturas idolorum, inclinari, incendere cereos, aqua conspergere sepulchra et seipsas. Haec quidem pietatis signa habita sunt iamdudum etiam apud ethnicos, sed minime sunt pieras, prorsus autem uersa sunt in impios usus. Pietatem magis profiteantur fideles foeminae et uiri, recordatione suorum peccatorum, suspiriis pro gratia et indulgentia, emendatione morum, et sancto proposito, attentius auscultari uerbis di\pend
\section*{AD TIMOTH. CAP. II. }
\marginpar{[ p.487 ]}\pstart uinis, audita meditari et ruminare, ut deinde liberis ea domi recenseant ac inculcent, ne quid aduersi uerbo dei admitti sinant in pueris, nec exemplis propriis in citent ad irreligiosa uerba uel opera, idque a tenera aetate et primis annis, nondum deprauatis animis. Orationes breues, sed feruidas dicant ex consideratione necessitatum profusas, et ex animo humiliato: non multas, nec toties repetitas pro more et praefixo numero. Non enim quantitatem precum iuxta numerum et tempus, sed qualitatem deus exigit et deuotionem. Offerant monetam probatam charitate, non loquacitate: domi orent, familiae prospiciant, marito obsequantur, et auxilientur in cura familiae. Caueant ab aspectu inuerecundo et uago. Lumen non lapidibus uel mortuis accendant, sed operi necessario, utili ac tenebris. Fides sine aquarum aspersione suffecerit. Pro omnibus istiusmodi exercitiis arripiant bona opera. Haec autem sunt, ut ministrent liberis, marito, familiae, uicinis, pauperibus, praegnantibus, parientibus, puerperis, infirmis, domui incumbant ornandae, lauandae, purgandae, prospiciendae die noctuque , iuxta Salomonis canticum. Omnia autem haecin fide peragant, persuasae certo talia deo placere, esseque ipsissimae pietatis officia, et opera a deo imposita, obedientiaeque sanctae merita, et fructus optimi. Si quando uacat, maxime autem feriis, lectioni incumbant, pueros ididem doceant, moneant, corrigant, et a malo exemplo in uerbis et operibus solicite sibi temperent.  \pend
\phantomsection
\addcontentsline{toc}{subsection}{\textit{iulier in silentio discat cum omni subiectione. Docere autem mulieri non permitto, neque dominari in uirum, sed esse in silentio. }}
\subsection*{\textit{iulier in silentio discat cum omni subiectione. Docere autem mulieri non permitto, neque dominari in uirum, sed esse in silentio. }}\pstart Viris quidem permittitur, imo iniungitur ut in coetu, quae de auditis non intelligunt, inquirant a senioribus et presbyteris in ecclesia docentibus: et si quid eis sedentibus spiritus reuelat, loquantur cum modestia et sine disceptatione, ad profectum ecclesiae. Non prohibetur a uerbo ueritatis proponendo uir qualiscunque  scripturas quacunque in lingua, sane intelligens. Voluit em̃ Paulus, ut passim quini aut seni pariter prophetarent, etiam intra ecclesiam. Cantica autem et organa psallentiumque clamores, ac tumultus inconditos, et studium assiduum canendi non docuit. In cordibus psallere, et spiritualibus hymnis domino laudes dicere, et canere mente, non tantum spiritu, hoc est sono, Christianos omnes docet. Mulieres uero in ecclesia sileant, domi uiros interrogent, si quid discendum fuerit, et in silentio:ne sensu suo abundare contingat, et litigare incipiant, uerbis et sensu proprio sibi complacentes. Docere mulieres non permitto, Paulus inquit: non ego ausim huic sententiae uerboque adiicere, in coetu uirorum, eatenus ut liceat foeminis docere in ecclesia foeminarum: sensu enim infirmae sunt, et non satis stabiles. Sed uiros presbyteros senes annis et moribus habeant agmina foeminarum, a quibus publice huiusmodi congregatio instruatur. Nemo uero negauerit, aliquot nonnunquam foeminas conferre posse laudabiliter et debere quoque de auditis ac praelectis. Et legere sacra scripta piorum doctorum etiam poterunt in coetu foeminarum, sic tamen ut intelligant, et non ignota in lingua. Verum sine intellectu orare, legere, canere, psallere, opus est sine omni fructu ac laude, maxima et perniciosa temporis perditio, et operum uel piorum uel alias utilium grauis iactura, qualia introduxit in ecclesiam sola pigritia ignorantiaque et superstitio monachorum, Olim et in Italia et in Germania quoque , uel teste Hieronymo, doctae erant sanctimoniales, quibus lectiones et cantica Latina lingua seruire poterant, et deuotioni conferre. Hodie in Germania et Gallia nihil intelligunt eorum quae iactant, uel sine iudicio aliquantulum suspicantur: quas non tam obedientia, quam hominibus et stultis monachis promittunt, excusas, que dementia culpar. Quid enim  \pend
\section*{COMMENT. IN I. EPIST. }
\marginpar{[ p.488 ]}\pstart inutilius institui poterat? quid stultius excogitari, quam ut tantundem temporis die noctuque expendant infelicissimae muliereulae in recitationibus, imo laboriosis cantibus uerborum, nihil prorsus illis significantium? Reliquum autem temporis ab huiusmodi tam molesta operis fatigatione, otio et remissione resolui et deperdi, ut omne earum tempus longe sit illis infructuosius, que ̃ somnium iners: interim tamen sibi de cultu diuino uehementer complacent, iactant merita gloriosa, quae et datis literis et sigillis uenditant benefactoribus: nec se patiuntur aliquatenus ab huiusmodi institutis diuelli. Loquimur nunc de eis monialibus, quae optimae habentur et religiosiores, misere deceptis a stultis monachis. Qui curae suae demandatas mulierculas, fidelius regere, et secundum canonicas scripturas docere, de domini dei sui certissima uoluntate, et de mandatorum illius obseruantia, qua domino seruuli seruire deberent ac possent. Quis enim domini sui seruitium seruus iactare potest, dum obmissis eis negotiis quae dominus prosequenda serio iniunxit, illis potius incumbit quae minime iussit, et suapte natura charitati non prosunt, quae est finis legis. De eis monialibus, quae sine disciplina uiuunt genio, miror quid mundus cogitet quantumlibet malus, et quid cogitent parentes, dum filias sic prostituunt miserrime, idque titulo diuini cultus, ubi tantum carni seruitur, et stolidis monachis confessoribus. Haec et similia maiore querela opus habent, quam ego ualeam consequi uerbis quibuscunque . Nec dominari in uirum.) Id enim foedum marito est et ridiculum, ac fomentum superbiae uxoribus, uirique ́ contemptus. Monere quidem maritum potest et debet, rogare, blandiri, et uincere malum in bono: interim tamen esse in silentio, parce loqui et breuiter, cum reuerentia ac timore: illico a uiro monita tacere, taceat: irato marito non respondeat, tempus captet opportunius: semel non crebrorepetat quae dicenda uiderit. Lingua mulieribus tanto diligentius moderanda est, quanto natura propensiores sunt ad cauilla et multiloquium. Nec circuncursari eas delectet per alienas aedes, non com morari in plateis, non foris plurimos ac in triuiis miscere sermones, domi agant potius cum silentio.  \pend
\phantomsection
\addcontentsline{toc}{subsection}{\textit{Adam enim primus formatus est, deinde Eua: et Adam non est sedu ctus, mulier autem seducta in praeuaricatione fuit. }}
\subsection*{\textit{Adam enim primus formatus est, deinde Eua: et Adam non est sedu ctus, mulier autem seducta in praeuaricatione fuit. }}\pstart Ne quopacto superbiendi ansam mulier sibi uendicet dominandi uiro, admonet Apostolus ut meminerint suae originis, quam a uiro sumpta accepit: qui primo conditus a deo, natura quoque perfectior, absolutior, sapientia clarior, et ad seductiones intentatas cautior atque fortior erat. Quo sopore merso, ex costa illius Eua diuina uirtute formata est:ut unum corpus de uno corpore, et os ex ossibus prodiens, amoris uinculo arctius astringerentur. Adam quidem prior uirtute et tempore, Eua posterior utroque :ideoque duci ea, instrui, parere et sequi debet, agno. scereque se dei ordinatione subiectam. Adam non fuit deceptus, non eum serpens aggredi primum audebat, qui eum fortassis minime decepisset sapientiorem foemina: non fefellisset tam impudenti mendacio, non ambiuisset excellentiam, non colloquium tam friuolum tulisset, nisi in gratiam mulieris consortis, amicae. Cui cum haec placere uideret Adam, delicias suas contristare non poterat, sed acquie uit, roganti mulieri attendens, oblitus praecepti uel negligens. Mulier autem seducta mendaciis, inuidia, dolo, blasphemia diaboli, permissione ac dispensatione diuina, nedum iusta, sed iustitiae originem nacta, qua proficere homo debuisset et progredi gratia dei praesente, acceptoque praecepto, retrocessit ipsa Eua, transgressa mandatum, et uerbis dei non credens, sed infidelitatis primum peceato su  \pend
\section*{AD TIMOTH. CAP. II. }
\marginpar{[ p.489 ]}\pstart perata, superbiaque et auaritia, gula et luxuria decepta est, et uiro insuper insidiata, blanditiis et instantia uicit, et in praeuaricationem traxit, ut pariter cum ea coeperit dubitare, si uere haec praedixeric uxor, scilicet imminere mortem, quam post gustum uideret superstitem ac uiuentem, ut et ipse peccaret. Cui semel seductrix facta, posthac directricis et doctoris munus ac officium praesumere et praestare nun quam par fuerit. Semel tam perniciosae suasioni consentiens Adam et seductus, meritosuspectam habuit omnem mulieris imormanenem et igemuim; ex ipsa experimento sumpto fragilitatis et insipientiae, cum perfecta esset, nunc bis misera et sui admonita, praesumere doctrinam non audeat solidae uernalis et sapientiae.  \pend
\phantomsection
\addcontentsline{toc}{subsection}{\textit{Saluabitur autem per filiorum generationem, si permanserit in fide et dilectione, et sanctificatione, cum sobrietate. }}
\subsection*{\textit{Saluabitur autem per filiorum generationem, si permanserit in fide et dilectione, et sanctificatione, cum sobrietate. }}\pstart Subdita quidem mulier sit uiro, et se humiliter gerere studeat inter uiros ac pudice, in ecclesia quoque nihil praesumat, agnoscens mentis imbecillitatem et inconstantiam, ut sibi docendi munus in ecclesia non arroget. Verum neutiquam ob id contemnenda est uiro, quam per fillorum et liberorum generationem beatam fecit dei munificentia, et insigni munere illustrauit. Quid enim admirabilius generatur in rerumnatura, ubi omnia tam sunt miraculosa, quam in utero soeminae: non tantum humanum corpus, sed et anima immortalis et rationalis creando in funditur, et infundendo creatur ?Cuius generationis, nutritionis, educationis pariter et instructionis officium dominus in matrem contulit et mulierem, incunabulaque tradendae pietatis et profectus matribus commisit circa prolem:ut sicut lacte pascunt corpusculum, sic solido cibo nutriantur a uiris: in fide ac sacris dogma tibus uiros foeminae perfectos audiant ac sequantur. Curet igitur liberorum mater, ut in ea quam a teneris unguiculis a se fidem et dilectionem des imbiberunt filii, in eis permaneant, iugiter inculcando fidem et timorem dei, ut sancti sint moribus, ab omni criminum contagione perpetuo sibi caueant. Carni et sanguini non indulgeant, dei mandata quibus eos mater edocuit, semper memoriter cogitantes, ut sancti sint et permaneant corpore et spiritu. Sobrietatem in primis soli cite curent, ut cibo, potu, somno, otio, lusibus, indulgentia, sodalitate, non nimis nec periculose imcumbant. Potius autem matres studeant liberis persuadere, quos domino genuerunt et proximis, fidem, dilectionem, sanctimoniam, pudicitiam, sobrietatem in uerbis et moribus. Ad hoc necessarius est foeminis matribus rerum illarum usus et studium, ut puellae didicerint quod matres docere possint: perniciosissima itaque negligentia est, si uirgines insolentius educentur, si sine timore dei, uerecundia, modestia, pudicitia, sobrietate et religione succrescamnt. Quomodo enim docebunt, quod nunquam didicerunt ? Quid sperandum de uxoribus, quae insolentia puellari deprauatae sunt, ut nesciant quid sit uerecundia uel honestas?  \pend
\endnumbering\beginnumbering\section{CAPVT III.}
\phantomsection
\addcontentsline{toc}{subsection}{\textit{\huge\textbf{F}\normalsize Idelis sermo. }}
\subsection*{\textit{\huge\textbf{F}\normalsize Idelis sermo. }}\pstart Quae tibi dilecto filio meo ad populum in ecclesia docendum scribout Christi apostolus, et tibi in eodem pater filio meo scribo, certissima habeas: et ea quae tam nunc tibi dixi de uiris et mulieribus, q̃ nunc scribenda cogito, scias a domino reuelata atque fidelia. Christi domini doctrinam esse, qui in me loquitur, non hominum uerba, id est mendacia. Nihil ex me loqui in ecclesia audeo, non opiniones ac fabulas comminiscor, fidelis est omnis dei sermo et efficax, quem mihi disseminandum credidit in orbem fidelibus.  \pend
\phantomsection
\addcontentsline{toc}{subsection}{\textit{Si quis episcopatum desiderat, bonum opus desiderat. }}
\subsection*{\textit{Si quis episcopatum desiderat, bonum opus desiderat. }}
\section*{COMMENT. IN I. EPIST. }
\marginpar{[ p.490 ]}\pstart Graece umοκοπ et, superintendens, obleruator, speculator et cultos dicuur, cuius officium episcopale est, pascere gregem Christi uerbo salutis, et pascuis doctrinae apostolicae atque canonicae. Scire utrunque testamentum, uetus ac nouum, gestare solicitam curam animabus Christo lucrificandis, procurare totius ecclesiae bona in spiritualibus, et a malis obseruare, inuigilare contra errores doctrinae et morum, haereses uerbo dei extirpare, plantare pietatis doctrinam, uisitare episcopatum suum, ut pastor gregum caulas: infirmos consolidare, agnos aegrotos sanare, confractos alligare, errantes reducere, pereuntes quaerere: curare utpascantur in herbis uirentibus, pascuis pinguibus, super montes Israel:iudicare inter pecus et pecus, opponere se bestiis pro ouibus, inuigilare ut secure dormiant in saltibus, quibus Christus pastor summus circumposuit benedictionem. Qui ergo episcopatum desiderat, huiusmodi exercitia praestare posse se sciens, uult quoque conferre in ecclesiam, is certe opus bonum desiderat, ex charitate, propter deum, et ut benefaciat multis a domino sibi commendatis. Desiderare autem honorem, eminentiam, census, pompam, dominia, luxum, otium, castra, oppida, prouincias, regna, uenationes, et id rerum multa, non est hoc desiderare opus, multo minus bonum opus. Tales hodie sunt, qui ambiunt ac desiderant episcopatus, uenantur, emunt, rapiunt, etc. Tempore Pauli et apostolorum non sic. Nihil hii minus sunt quam episcopi, principes sunt et dicuntur, ac haberi uolunt: et huius mun di principes sunt, satrapae Imperii, quibus priscus adhuc placet status, cum essent flamines, archiflamines, primates, ac prouinciarum Romanarum praesides, et sa crorum pontifices apud ethnicos. Qui si animarum saluti, et doctrinae pietatis per occupationes saeculares incumbere uel non possent uel nollent: et, ut nunc fit, non tam episcopi nominari eligerent, que esse administratores ecclesiarum, saltem con tenti essent emolumentis prouincialibus, et turpem illum quaestum, honoribus suis iniquum, abhorrerent, qualem accipiunt de scortationibus sacerdotum, spu riorum redemptionibus, et casibus reseruatis: de quibus lucris nemo honestus se notari pateretur. Operam autem sedulam nauarent, ut sub administratione et tutela sua essent, sub praefectura sua uel sua diocoesi priuatarum ecclesiarum ueri epi scopi, quales hoc loco describit Apostolus, quales habuit primitiua ecclesia mul. tos, quorum catholicam doctrinam, et uitae innocentiam, et apostolicos mores il li ecclesiarum principes moderarentur atque defenderent: et honesta prouisione rerum temporalium, non ad ullam superfluitatem prospicerent, in maioribus eccle siis et minoribus iuxta exigentiam honestatis. Id certe quod facile fieri posset, et sine tumultu uel confusione prouinciarum et imperii, uel quorumcunque regnorum. Exonerarentur uel pro tanto, si non pro toto, illi nostri episcopi, grauissima sarcina conscientiarum suarum, quibus nunquam satisfacere possunt uel uolunt, et interim principes consulerent rebuspublicis, et principibus masoribus imperioque , ut hactenus. Et tamen fides, ac morum innocentia, et ecclesiarum pietas minime periclitaretur, ut nunc factum tanto tempore uidemus et dolemus. Quin et omnes nos Christiani homines nostro quoque modo episcopi simus necesse est, uel hac parte, ut officium nostrum sit uisitare, curare, iuuare, consolari proximos: unicuique enim de proximo suo mandauit dominus. Qui communitatem curat, custodit, nutrit, docet, fouet, fideliter regit, hic uere est episcopus: nomen est enim functio nis et operis, officii et sedulitatis, ac fidei, non honoris et eminentiae: sicut sunt antistes, praepositus, prior, abbas aut praeceptor.  \pend
\phantomsection
\addcontentsline{toc}{subsection}{\textit{Oportet enim episcopum irreprehensibilem esse. }}
\subsection*{\textit{Oportet enim episcopum irreprehensibilem esse. }}
\section*{AD TIMOTH. CAP. III. }
\marginpar{[ p.491 ]}\pstart Dicit non sine causa, Oportet igitur: est enim quod sequitur, totum de necessitate officii, et apostolici praecepti, imo diuini. Quid hodie facerent, qui corrupte larum suarum patrocinium quaerunt, tueriq uolunt praescriptione diu turni tem poris, idque etiam ex uerbis dei et scripturis canonicis? si hic uel alibi inuenirent hoc uerbum, Oportet episcopum esse coelibem, etiamsi continere non possit:opor tet eum habere omnes dioecesis sui decimas, oportet primos fructus illos de qualibet ecclesia percipere, etc. Profecto ob uerbum illud Oportet, dicerent esse mandatum diuinum, sub poena peccati mortalis et aeternae damnationis oseruandum, Nunc autem quid dicunt ad hoc, Oportet episcopum esse irreprehensibilem: non quidem sine peccato, pro quo nemo ab alio reprehensibilis sit, qui omnes peccatores sumus, et in multis offendimus omnes: sed ut nihil criminis impingi possit episcopo, nulla saecularis macula, non aliqua nota infamiae, quo minus aptus sit docere, obsecrare, increpare, arguere libera fronte criminosos quoscunque . Quis igitur nesciat, optimos quosque hactenus habitos, multosque passim parochos et paroecianos uiros doctos, benignos, solicitos, deum timentes, nullo crimine notabi les, sed hoc solo reprehensibiles uitio, quia continentiae gratiam non acceperant: quorum qui meliores erant, concubinas assequebantur, ne deteriora contingerent: reliqui uero continentiam, quam minime habebant, singebant, interim inardescentes ad aliorum coniuges et filias, propter quas et in foro confessionis abos minabilia parabantur. Quomodo talis, non dico hoc crimen adulterii, uel scortationis succrescens in ecclesia, acriter arguere et punire potuisset? sed nec alia quaecunque crimina, uel semel in ecclesia sic criminosus libere taxasset. Sunt autem epi scopi, et uere dicuntur, omnes ecclesiarum praelati, non solum magnarum, sed etiam minimarum. Habeantur ergo irreprehensibiles pastores ecclessarum, innoxii, integrae famae: hoc autem fiet, si qui nunc continuo sequuntur obseruata fuerint. Subinde enim multis docet, quomodo possit esse episcopus irreprehensibilis, quae et consequenter erunt diligentius annotanda.  \pend
\phantomsection
\addcontentsline{toc}{subsection}{\textit{Vnius uxoris uirum, sobrium, ornatum, prudentem, pudicum, hospitalem, doctorem, non uinolentum, non percussorem, sed modestum, non litigiosum, non cupidum, sed suae domui bene praepositum, filios habentem subditos cum omni castitate. }}
\subsection*{\textit{Vnius uxoris uirum, sobrium, ornatum, prudentem, pudicum, hospitalem, doctorem, non uinolentum, non percussorem, sed modestum, non litigiosum, non cupidum, sed suae domui bene praepositum, filios habentem subditos cum omni castitate. }}\pstart Sancti patres et patriarchae Hebraeorum plures simul habebant uxores, et a domino ueluti etiam amicissimi ferebantur, propter ingens desiderium multorum bonorum filiorum, quae ut habebatur, sic erat, et promissa fuerat sanctis benedictio domini felicissima, secundum hoc saeculum. Christus idipsum nusquam legitur prohibuisse, uel etiam improbasse. Diuinae literae pluralitatem uxorum nusquam da mnant, nec tamen consulunt, neque laudant. Hoc autem loco episcopis et diaconis prohibetur huiusmodi uxorum pluralitas. Et cum iam statim praemissum legamus uerbum Oportet, non sine causa Chrysostomus periclitatur, ne Paulus exigere uideatur matrimonium ab omni episcopo. Nec idipsum abre et leuibus rationibus. Certe primum ut non tententur a satana. Secundo, ne infamiam et suspicionem incontinentiae facile incidant. Tertio, ut habita occasione, discant familiam regere, sicque etiam apti reddantur ad gubernandas ecclesias. Quarto, ut hospitales esse possint, et parati ad susceptionem pauperum et peregrinorum, patriarcharum more. Quinto, ut a studio scripturae sacrae discendae, meditandae et docen dae domestica cura minimum distrahantur. Sexto, ut aliae cuiuscunque mulieris non sit necessaria suspecta quaecunque familiaritas. Septimo, quod sunt in domo, nego  \pend
\section*{COMMENT. IN I. EPIST. }
\marginpar{[ p.492 ]}\pstart tia non pauca, quae sine foemina expediri nequeant. Ocrauo, cohabitatiouiri et mulierss extra matrimonium, non potest carere incommodis plaerisque , animae, conscientiae, famae, et rei familiaris. Vt interim taceam sa crilegum illud et scandalosissimum scandalum concubinatus illorum episcoporum, hoc est plebanorum: quod ipsum liquer probe consideranti, fontem esse omnis Christianae iacturae, da mnationis, infelicitatis et omnium criminum colluuionis in orbe Christiano, in omnibus statibus. Quibus abunde consultum fieret, si tamen ecclesiae reformationem praelati ferre uolent uel possunt, si maritorum lege uiuerent sacerdotes. Inuenirent enim facile (cessantibus impiis legibus humanis, et contra uerbum dei fancientibus) boni illi paroeciani honestam proborum uel etiam diuitum (si uellent et expedire uiderent) puellam uirginem, in disciplina nutritam a primis annis, castam, fidelem marito, filiis, domui, familiae. Praeesset huiusmodi utiliter, et cum parsimonia domui sacerdotis, reponeret pro haeredibus filiis, quae de uictu, amictu et eleemosyna superessent: disceret a marito uerba dei, doceret paruulos, honoris aliquando consequendi intuitu, et spe uerisimili, prae filiis aliorum promouendorum ad honesta. Esset amabilis omnibus ac honorabilis in ecclesia. Sacerdos eam pudice diligeret, et prudenter aleret prolem ex ea, dei beneficio susceptum. Vacaret quoque cum tranquillitate sacris et bonis literis, sine trepidatione doceret et cum fiducia omnem ueritatem. Subditos diligeret, diligeretur a piis, defenderetur a fidelibus, corrigeret sine rubore et iusto necessarioque zelo de omnibus, haberet sibi fauorabiles, beneficos, promptos, et doctrinae obedientes primores in magistratu, unde mores facile restitui possent nimium corrupti. Quid mul tis ?Infinita bona consequerentur. Nunc autem, proh dolor et pudor, sacerdos parochus quotidie uel frequenter celebrans sacra sua, et sacramenta ministrans, publicum autem, ut fit, scortum nutriens de decimis, oblatlonibus, eleemosynis, et pauperum sudoribus, cum uel sine uerecundia, et cum grauissimoscandalo, nónne se damnationis mancipium nouit? et certissime secundum legem domini, et humanas constitutiones sese cogitur impium et sacrilegum iudicare ? Crimine graui et publico scandalosus, et transgressor legum diuinarum et humana rum. Non uult nec potest dimittere concubinam, nisi alia surroganda infelici mi seria, dum statim suecedunt ex noua noui filii, et perpetua calamitas, et scandalum intolerabile. Orationes suas, quas uocant Canonicas horas, (quia canonicis scripturis ediscendis ac intelligendis instare, multis horis indies noctuque deberent, iuxta antiquissimas probatasque sanctiones ecclesiasticas) has, inquam, si dicit, ibi continuo omia offendit contra suam conscientiam, infernum sibi interminantia: cum horrore ergo et nausea omnia legit. Odit potius sacrarum studia literarum, otio solo sese ac uitiis accommodat, ut fallat tempus, et taedium superet, infeliciterque ́ consolatur alea, perpotationibus, et reliquis id rerum. Vritur importunius, grauatur libidine insatiabili ex orio: et quia prohibita magis placent, et illicita, ten tatur ardentius. Id quoque temporibus facit etiam religioni suae consecratis. Sequi tur continuo, et nunquam cessat conscientiae uermis, uel desperatio infernalis. Fit ex hoc procliuior ad omnia uitia. Coram deo confunditur, coram conscientia, populo et familia. Sacramenta tractat ut impurissimus porcus, sine fide, reueren cia et honore, sed potius cum horrore et perfidia, ad commodum temporale tantum respiciens: tandemque fit totus sine deo, qui apud deum pro populo intercessor habetur, et ob cultum dei nutritur in blasphemiam dei. Quid autem interim huiusmodi sacrifsculi tam sordidi agit uel cogitat infelicissima meretrix et concubina ? Filios fortassis generat, si tamen abortum non procurat, ne uel episcopo pretium  \pend
\section*{AD TIMOTH. CAP. LII. }
\marginpar{[ p.493 ]}\pstart Ianguinis obligetur, uel a libidine arceatur continuanda. Deserere filios et scortatorem non potest, otio allecta, ac libidine uictus et amictus. Non fert redire ad parentes uel consanguineos: non facile procum inueniet, non maritum assequetur: nouit se despectam omnibus, etiam a filiis propriis: non ignorat facili occasio ne se eiiciendam a sacerdote tali, si libidinosior superuenerit meretrix: miserrimam post obitum sacerdotis, siquidem otio asueta, seruire nolit, et deliciis resoluta non possit: ideo interim sua maxime curat, furatur, exigit a scortatore quicquid ualet, praeparat uestes et multas, et meretriceas, ut post illum uel cum illo aliis placeat, quos et facile admittit, sciens se non propterea adulterarum lege abiiciendam, uel non esse adulteram, quae non est ligata marito. Filios non amat, negligit, non docet, sed ut non suos abiicit et propellit. De animae suae falute desperat, quam et lu dis publicis a rusticis daemonibus persequi, et uenationibus nocturnis patitur et audit iniuriosissime tractari, ueluti ab infernalibus spiritibus. Desperans ergo, re labitur ad turpitudines, et ad uitia fit nimium exposita ac procliuis. Vicinis uiris et mulieribus fit odibilis, rixosa, onerosa et intractabilis, dum se contemni non pa titur, dum speciosior et ornatior incedit, quam honestae matronae, dum irritat sacerdo tem contra uicinos, sese quoquo modo dehonestantes, uel etiam negligentes. Eilios et filias indulgentius nutrit et fouet: non cohibet insolentes, non uitia redarguit in eis, quibus se sentit refertam ac scatentem. Ingerere se temerarie annititur et effrons primam coetui foeminarum, iactat conditionem uictus sui, otii et libidi num: quae omnia probae matronae et puellae indignissime ferunt. Deinde filii sacerdotum talium inter tot scandala educantur, ut non sint illis amplius scandala, sed consuetudo prouinciae et gentium, imo ecclesiarum. A parentibus ut illegiti mi negliguntur, sine spe consequendi honoris, prorsus extorres ab omni honesta et graui conuersatione. Fiunt sic manumissi rebelles, indisciplinati, blasphemi, ri xosi, procaces, lusores, et parentum contemptores: non ferunt disciplinam etiam apud exteros, non uerbo, non facto dictis honestorum praeceptorum obtemperant: ab omnibus despiciuntur ut spurii: nihil facile eis boni et honesti alicubi con fiditur, peiores ut plurimum deprehenduntur omnibus coaetaneis: audaces ad omne flagitium, diuinorum omnium aduersatores: expelluntur ante tempus in exteros, libertatem captant ad omne nefas, cuius semina susceperunt ex educationae domoque paterna. Haec aetas et nobis posterior uix credet, quae hactenus infelici tem porum deprauatione uisa communiter sunt et perpessa in Germania nostra et alibi. Quanquam non defuerint semper, qui animo coniugali cohabitauerint, ho nestatis iacturam et ipsi deplorauerint, et prolem, quanquam raro, ad honores instituerint capescendos, maxime per Curtisanos et Romanas merces, tanquam eca clesiae filios, sic enim sunt appellati. Et quid interim populo accidit ? Intuetur miser dei populus et ecclesiarum uulgus, sacerdotem et pastorem suum speculum abominationis: uerbum dei, si audit ex illius ore, non credit, contemnit magis, de trahit, blasphemat, sectatur pastorem per abrupta gradientem in praecipitium grex infelix: correctiones non audit, uel non suscipit, ridet potius, cum culpa redarguit doctorem, incalescit magis et feruet ad proprias libidines, exemplo impii scortatoris. Dolum suspicatur in doctrina et uerbo dei, subuersionisque imposturam: sancta dogmata non internoscit, sed abiicitomnia: euitat ac declinat ecclesiam, iussis non obtemperat, fidem ueram ignorat, uitia sectatur, utiam asueta: non fert monitorem quempiam de fide uel moribus. Conscientia nulla fit criminum, uel si adsit, cauteriatur: sicque totus Christianismus perditur et euanescit. Peccatorum confessio quidmirum siobmittitur, tam notorie criminoso facienda Decime negantur ec\pend
\section*{COMMENT. IN I. EPIST. }
\marginpar{[ p.494 ]}\pstart clesiis. Ceremoniae ut sordidae propter fordidum ministrum irridentur. Interim male audit ipsum totum sacerdotium, et omnis denique Christiana religio. Verbum et fides contemnunt, non praedicat sacra scriptura, non discit, adulterinae doctrinae surrogantur, uitia publica palam defenduntur, excusatur error, scandalum iusti ficatur, piorum ora obturantur, exquiruntur ansae et rationes defensionum. Decanorum quoque fauores ambiuntur, redimuntur muneribus, excogitantur grauamina laicorum, exiguntur oblationes, aggrauatur stola, quam nominant, quae stus scilicet officiorum sacramentalium, mortuorum exequiae et peractiones commendantur et exiguntur: et pro sacrificiorum multiplicatione et anniuersariis statuendis, purgatorium Vergilianum incalescit. Quandoquidem modicis redditibus et accidentalibus quaestibus haec sacerdotis famula uel familia non ualeat enutriri. Turpe lucrum cupidus episcopus imperat, Diocoesanus uel eius uicarius, censumque auget quotannis. Cogitur talis sacrificulus maiorum de plebe fauores promereri per fas et nefas, quo subsistere in officio ualeat, uel inuitis subdi tis. Officio suo diuino et Missationi in publica ecclesia et circa aram, assistuntno uellae turpitudinis plantationes, spurii sui filii. Breuiter, exemplum sunt huiusmodi concubinarii sacerdotes flagitiorum omnium, et uelamen malitiae magistratibus et populi rectoribus. Haec, inquam, omnia et his longe plura abominabilia et scandalosa, subintroduxit nobis infelicissima necessitas coelibatus, et slulra prae tensa religio, contra stimulum calcitrandi, et angelicam puritatem imponendi eis, qui carne et sanguine uelint nolint opprimuntur et ardent: paucissimorumque ́ priuilegium perpetuae castitatis, imo uirginitatis, quae iactatur, tot millibus homi num utriusque sexus commune faciunt, reclamantibus conscientia, experientia, et dei uerborum obedientia. Quibus omnibus tot malis abunde consuleretur, si uerbo dei et sacrae scripturae institutis ac praeceptis obediretur. Honestissimum matrimonium fuit sacerdotibus Mosaicae legis, sicut et primitiuae ecclesiae Christianae episcopis, ut patet ex Ecclesiastica historia Eusebii et aliorum. Conditor quoque legis dominus omniporens, et amator hominum, prolem et matrimonium ubique et plurimum commendauit. Patriarcharum nota ex scripturis sunt sancta coniugia, et prophetarum, sicut et apostolorum. Christus toto suo euangelio, inter tot sacerdotum ac scribarum scelera uel uitia, ne semel quidem uxorum uel plurium uel unius habuit mentionem, praesciuitomnia, et matrimonium commune uoluit omnibus, interdictum autem meretricium et fornicationem. Primitiua ec clesia apostolica hunc coelibatum sub uoto omnino nesciuit. Orientalis sacrosan cta ecclesia hactenus ignorauit et ignorat hodie. Pontificorum indoctorum, superstitiosorum ac impiorum nuper orta tyrannis, haec animarum naufragia subornauit, quibus et ipsimet pereant, satis perspicue alicubi ostendentes fructus suae turpissimae puritatis, quo confundantur cum uniuersis suis impiis constitutionibus, quas spiritus sancti et uerbi dei flatu propediem exufflandas ueluti aranearum tela, et conscientiae laqueos minime possumus dubitare, nisi uerbum dei nullum mundus susceperit uel habeat. Cum pusillis clericis hactenus interdicta mul ta fuerint ob quaestum, quod pinguibus pontificibus perpetuum fuit ad luxum. Potuit quidem, in summa, uideri priscis quibusdam episcopis commodus admodum coelibatus clericorum. Primo, ut purius uiuerent. Secundo, dignius sacrosancta tractarent mysteria. Tertio, expeditius scripturae operam nauarent. Quar to, terrenis actibus minus inquietarentur. Quinto, ut minus grauaretur populus uel ecclesia sustentatione uxorum et prolis. Cum uero omnia iamdudum in diuer sa cessisse uideantur, eoque deuentum esse sciatur, ut sacrificuli isti passim im  \pend
\section*{AD TIMOTH. CAP. III. }
\marginpar{[ p.495 ]}\pstart purissime uiuant secundum praedicta, scientibus ac permittentibus plebe et primoribus. Secundo, ludibrio quoque et pro nausea habeant sacrosancta omnia homines ir religiosissimi. Tertio, ludo potius ac poculis incumbant iugiter quam scripturis canoni cis. Quarto, otio inertissimo resoluti prolabantur in publica crimina, quae etiam episcopi, si uelint, emendare nequeant, Romana authoritate praepediti uel exem ptionibus nonnunquam. Quinto, spuriis quoque et meretricibus non solum grauetur, sed et publice diffametur ecclesia. Verbum quoque domini ueniat in contenptum, et clerus non indigne in odium, ac in scandalum toti Christianae religioni, maxime apud infideles quoscunque . Ideo nulli sit dubium, obligari arctissime archiepiscopos et episcopos in solidum, sub aeternae maledictionis iustissima certissimaque uindicta, si grauissimam dei iram euadere uelint, ut ocyus ac sine dilatione hoc scandalum coelibatus, et uotorum perpetuae castitatis, ut ipsi intelligunt, et perniciem inexplebilem de medio tollant. Si non concilio generali, quod legitimum sustinere nolunt, qui lucem odiunt et ueritatem, saltem prouinciali Germanico, ubi concubinatus scandalosius relucet. Sin minus et hoc cesserit, quilibet episcopus, ut nune sunt, in sua dioecesi huic pestilentissimo succurrat morbo, communi consensu et synodo. Quod si neglexerint episcopi, debent ac possunt authoritate dei omnipotentis et id praecipientis per spiritum sanctum per seipsos qui continere non possunt, experientia teste, et naturall conditione uel complexione admoniti, propria authoritate sibi uxorem legitimam ducere, praemissa sufficiente instructione plebis, ut sciant praeceptum dei non esse clericalem coelibatum, Interim commode magistratuum fauore redempto, et uerbo dei institutis, si uelint consentire, alioqui uel officio resignandum magis, uel inuitis omnibus, domi no acquiescant, necessario et ineuitabili mandato omnipotentis dei. Sine quo ma trimonio saluari neutiquam possint, nec praeesse populo dei, nec satisfacere uera bo dei, ned apostolicae traditioni obedire, quam summus ille, ut uocant, pontifex sua authoritate in integrum restituere collapsam religionem obligaretur, sicut et dispensationem in cibis et uotis reliquis, quae minime obseruari contra conscien tiam omnium palam intuetur, non ignorans aedificare ad gehennam, quicquid committitur contra conscientiam. Quae sit autem apostolica uera traditio, hoc lo co et alibi legitur. Episcopum, id est pastorem in spiritualibus populi Christiani, oportere habere unam uxorem, nisi possit capere donum dei, quod a solo deo paucis et ex raragratia donatur, munus inclytum uirginitatis, et magno periculo seruatur. Cuius denique continentiae difficultatem, imo impossibilitatem, publice testatur abominabilis spurcitia totius ferme cleri:non solum plebanorum, de quibus dictum est, canonicorum et episcoporum, sed etiam monachorum et fratrum, monialium quoque et nonnarum patens impudicitia, nimium generalis ad irrisionem status clericalis iustissimam. Nihil contra hanc castitatem, ut nunc seruatur, satis dure seueriterque dici uel scribi potest, ab aliquo earum rerum non inexperto. Sunt enim omnia quae dicuntur, longe quoque atrociora, quam pro rei indignitate a nobis proponi et plangi possint. Excessum hunc nostrum pro temporis huius occasione et studiorum, alius quispiam fusius et ordinatius tractabit, ad ea me magis idoneus. Sobrium) addit. Non idcirco uxor adiuncta est ab Apostolo, ut episcopus uacet libidini cuicumque , sed ut illi medeatur. Minus enim a carne et satana infestan tur infra cancellos benedicti a deo matrimonii, quam promiscua illa ac libera libidine palantes, ut animalia, quibus non est disciplina uel intellectus. Quin et hic statim Apostolus aliud continentiae remedium adiicit episcopo coniugi, ut sobrius sit: cibo et potu utatur, ad corporis non uoluptatem, sed necessitatem: parsimoniam  \pend
\section*{COMMENT. IN I. EPIST. }
\marginpar{[ p.496 ]}\pstart seruet, ne et ecclesia grauetur expensis, uel corpus ac ingenium sneptumsit ad exercitla spiritualia, digna sacerdoti ac doctori ecclesiastico, scillcet ad studia sacrae scripturae, praedicationes, consultationes, consolationes, orationes, praelectiones, scriptitationes, etc. Exemplum sobrietatis praebeat, non pretiosa nec plurima uel delicata fercula, sed utatur semper cibo potuque pro corporis ratione, et officiorum executione opportunis. Non sic se die uno exinanire studeat irrationabili abstinentia, ut sequenti die denuo oneretur, ne utriusque diei sacrum opus intermitti con tingat. Vt nunc consueuerunt stulte ieiunantes superstitiosuli in uigiliis festorum, et ad honorem sanctorum abstinere, ut festis ipsis sequentibus diebus repleri, et neglecta resarcire possint, auidiusque pasci ac delicatius, quando toti sese spiritualibus dedere debere ratio dictat, et ecclesiastica mandat traditio. Tunc enim carnalissimi inueniuntur, et ad succumbendum tentationibus procliues magis, quam dum laborando cibum absumunt sobrie captum. Ornatum) dicit. Non sericis et aurophrygiatis uestibus, argento uel gemmis, siue in capite siue sandaliis, aut pedo, uel incuruo pastorali baculo. Haec enim omnino dedecori potius sunt sacerdoti, margaritis ornata, sed pudicis mansuetissimisque moribus redimitus sit, fa ctis ac habitu uenerandum se exhibeat, et legitima maturitate gratiosum fidelibus, perfidis uero horrendum et formidolosum: talia sunt enim episcoporum ornamenta legitima. Prudentem quoque .) Vt plebi commissae prospiciat, cui specu lator proponitur circumspectus, ut praeuenire malum, et sancta promouere quopacto possit intelligat. Sophrona pudicum dicunt, et quadrat praedictis, ut sit eatenus sobrius, ut et pudicus esse ualeat uisu, uerbo, signo, opere, et in omni conuersatione sua praeferat pudicitiae ac prudentiae spiritualis exemplar. Hospitalem.) Non sit publicus tabernarius, sed indigentium peregrinorum benignus su sceptor, et tractator beneficus, quales in ueteri testamento multos legimus prophetas et patriarchas: ne honesti uiri ad peregrinandum coacti, communi diuer sorio sub turba et clamoribus praepediantur a studio orationis, et meditationis sa crae, et quiete opportuna. Non solum hospitem se praebeat sacerdotum, sed piorum omnium indigentium. Ad quod sane conueniens est, habere non pellicem, sed uxorem fidelem, et matremfamilias, filliosque disciplinatos, in disciplina ecclesiasti ca eruditos, et exercitatos in operibus pietatis: ne toti ecclesiae cui sacerdos praefertur, passim inuratur diffamia, dum peregrini praeter Christianam charitatem tractantur. Ea ratione decimae illis debentur, non tam pro se, familia, et hospitibus, q̃ pro orphanis, pupillis, infirmis, egenis, et senibus, quorum omnium sacer doti episcopo studiosa incumbit diligensque cura: non solum uerbis, sed et factis con solandis et beneficiis iuuandis, nomine totius ecclesiae, ad talia exercitia obligatae: quamobrem et maioribus ecclesiae praefectis diaconi coniungebantur, ne quid in officio ecclesiastico negligeretur: non in cantu et processionibus, etc. Docto rem.) Qui aptus sit ad docendum ecclesiam, non aliam doctrinam quam canonicam et catholicam ex sacris scripturis, sapientiamutilem ad dei cultum, non promotum a facultate theologica in gymnasio, non empta facultate, non pecuniis prodi gendis in comessatores, ut interim fit, sed spiritu sancto edoctum ex sacris literis, qui instruat plebem, cui ex officio suo superintendere debet. Non uenator, non eques, non lusor, non auceps, uel piscator, sed uerbi dei doctor doctus. Hoc autem proprium et peculiare prae omibus officium est episcopale, Pascere Christi oues, quas amet plurimum, docere non philosophiam, non grammaticam, non poëticam, non constitutiones humanas, sed quae pertinent ad cultum uerum dei, ad pietatem, ad fidem et probatissimos mores, quorum exempla a Christo sumgrur.  \pend
\section*{AD TIMOTH. CAP. III. }
\marginpar{[ p.497 ]}\pstart Non uinolentum.) Qui nimium indulgeat uino potando, quo impediatur a modestia et doctrina in facto ac uerbo, utque dum opus est, cum opportunitate corri, piat, ac sine conuitiis, ebriosorum more. Non combibonem, qui huiusmodi abusus uehementer in aliis arguat ac corrigat. Obtundit enim uinum immodice sumptum ingenium, inepte laetificat, reddit consilii inopem, iracundum ac foedum pariter reddit, pudicitiae aduersatur ac castitati. Non percussorem.) Non promptum et propensum ad uindicandum se ad ulciscendum impraemeditato, sed ani mo tranquillum, longanimem, sustinentiam fortem conseruantem ac modestiam, etiam in medio iniuriarum. Non paratum ad arma inferenda, sed qui in bono ma lum uincere gaudeat: percutiat neminem, sed euangelica patientia sese common stret euangelicum praeconem, et doctorem mansuerudinis Christianae, exemplo et doctrina uere apostolica. Non enim resipit euangelicum et apostolieum spiritum tumultuari, non solum factis, sed et uerbis: exemplo magis utatur Christi et apostolorum, qui omnia matura grauitate egerunt, non maledicentia affectata, cum irrisione, quam ut apostoli et prophetae non obmiserint semper, tamen non sine magna modestia, et uirtute spiritus, fortissimaeque ueritatis: unde hoc non solum ad uiolentiam manuum pertinet, sed et ad acerbitatem linguae, ne saeuus et im probus sit obiurgator. Non auarum.) Vel turpiter lucri cupidum. Hoc quoque  audiant nostri episcopi, qui lucrum congerunt ex casibus reseruatis, ex sacerdotum concubinatu, ex sacramentorum et ordinum administrationibus, ex spurio rum redemptionibus: si non per se talia capiunt, subornant tamen Officiales, qui bus id negotii foede committunt, uel oblucrum consentiunt: communius tamen est, turpiter lucri cupidum, quam turpis lucri cupidum esse. Omne lucrum cupere, indignum est episcopo, cuius est officium expendere potius quam lucrari uel thesaurizare. Bona enim episcoporum, ecclesiarum sunt, hoc est indigentium de ecclesia, uiduarum, pupillorum et egenorum. Reponere thesauros non debent, ut ser uentur, nec ut multiplicentur in census asseruandi, et eisdem bonis per usuram quancunque lucrum turpe uenari et congerere. Sed modestum.) Aequum esse et fidum omnibus, aequanimiter ferre potens infirmitates et imperfectiones fragilium, donec roborentur in fide, et agnitione ueritatis et iustitiae. Non acutum, iracundum, uel ad uindicandum proprias inlurias et personales animum propensum gerens, sed modestia Christiana prae aliis ornatum. Non litigiosum uel contentiosum.) Vtboni consulat uerba ac facta aliorum, etiam contentiosorum: sed pacificum, non suscitet bella ac lites, sed tota diligentia componat. Quomodo autem carebunt discordiis et disceptationibus, qui diuitiis augendis inhiant, et possessionibus multiplicandis ? Dirimat potius contentiones, et pacificus sit, ut et beatus a domino iudicetur, si pacis quae sunt ut Christi discipulus studiose promo ueat. Non cupidum.) Quaestum nullo modo sectetur, ut Simoniaci, sed gratis accepta, gratis concedat. Liberalis in pauperes: non lanae potius abstractor quam uerbi dei administrator. Abomni specie cupiditatis abstineat. Sed domui suae bene praepositum.) Quomodo sine uxore bene praeerit, et cum spuriis, uel exter no famulitio? Vult Apostolus, ut familiam suam, uxorem, filios, famulos bene moderetur, ne scandalosa ac inordinata domus, toti ecclesiae fidelium obsit ac praeiu dicet. Si non huiusmodi episcopus familiam habeat, saltem canonicis suis religiose praesit antistes. Quos hodie nouimus irreligiosissimo iure exemptos authoritate Romana a iurisdictione episcoporum, et omnium scandalorum antesignanos in clero toto totius episcopatus uel dioecesis, et exemplaria insolentiarum in clero et populo. Filios habentem subditos cum omni castitate,) Disciplinatos, et  \pend
\section*{COMMENT. IN I. EPIST. }
\marginpar{[ p.498 ]}\pstart religiose cum reuerentia subiectos habeat filios suos, qui moribus et honestate reliquorum filios instituant, et exemplum honestatis exhibeant toti ecclesiae, patri obedientes, reuerentiam exhibentes. Quid hodie de filiis presbyterorum mundus iudicet, notius est quam̃; ut multis recitetur. Debet episcopus in propria sua domo di scere exercitio filiorum, quod doceat in ecclesia subditorum. Vnde et sequitur.  \pend
\phantomsection
\addcontentsline{toc}{subsection}{\textit{Quod si quis domui suae bene praeesse nescit, quomodo eccleliae dei diligentiam habebit? }}
\subsection*{\textit{Quod si quis domui suae bene praeesse nescit, quomodo eccleliae dei diligentiam habebit? }}\pstart Quomodo multos et quasi alienos curabit, qui propriis ac paucis prodesse non potest, et tamen ad id positus est utprosit. Qui in paruo infidelis deprehendi tur, quomodo super plurima constituetur Qui familiaribus non prospicit et necessariis, quomodo alienos dispensabit? Qui tam tenui sensu est, ut pauca cogitare nequeat, quomodo sapienter totam ecclesiam gubernabit ? Exercitari igitur in domo propria oportet, et in liberorum sancta educatione, ut regimen dioecesis totius ualeat sustinere.  \pend
\phantomsection
\addcontentsline{toc}{subsection}{\textit{Non nouitium uel neophitum, ne in superbiam elatus, in iudicium incidat diaboli. }}
\subsection*{\textit{Non nouitium uel neophitum, ne in superbiam elatus, in iudicium incidat diaboli. }}\pstart Nouus in fide conuersus nuper a quacunque mala uel inhonesta conuerlation ne, nondum exercitium spiritus Christi in se expertus, non est idoneus doctor ec clesiae. Cum enim alios docere et moderari instituerit, in his quae ad deum sunt, expedit exercitatum habere sensum in spiritualibus. Nemo autem repente fit summus, studium quoque sacrarum literarum non subito proficit. Verum nouitiatum illum non uult a baptismatis tempore computandum, quod multi olim tardius baptizabantur, sed a conuersionis tempore ad fidem. Ambrosius enim catechumenus diu Christianus fuerat, antequam baptizaretur, et episcopus nondum ba prizatus eligebatur spiritu dei promotus. Verum non expedit religione nouitium subito gubernacula religionis moderari, ne uel religio uituperetur, quasi deslituta idoneis pro huiusmodi functione personis, quam probi et sapientes omnes declinant, et cogatur aliunde pellicere promouendos. Vel ipse nouitius honoris sui sortem humiliter ferre non potens, insolentius se efferre incipiat, et statim sibi licere putet quicquid libuerit, unde externi calumniandi occasionem arripiant, quasi ea tantum ratione ecclesiae se immiscuerit, ut ab ecclesia honoretur. Denique facile honor obnoxius fit inuidiae, unde prodit temerarium iudicium et calu mnia. Cum quicquid a neophitis agitur, quamlibet religiose, insimulatores tamen patitur, etiam sancta inuertentes. Et cum aliud nihil obtendere possunt, saltem nouitiatum causantur. Sed et huiusmodi neophiti facile praesumunt maiora et insueta, audent iugum imponere, eleuanturque existimatione propria, sicque in iudicium incidunt calumniatorum, imo et sapientum nonnunquam.  \pend
\phantomsection
\addcontentsline{toc}{subsection}{\textit{Oportet autem et illum testimonium habere bonum ab his qui foris sunt, ut non in opprobrium incidat, et in laqueum diaboli. }}
\subsection*{\textit{Oportet autem et illum testimonium habere bonum ab his qui foris sunt, ut non in opprobrium incidat, et in laqueum diaboli. }}\pstart Ad quemlibet hoc episcopum referendum uidetur, ut scilicet etiam ab his qui foris sunt, et ab omnibus bonum habeat testimonium, ne religionem ob ministri improbitatem calumniam sustinere contingat. Verum specialiter, si a Iudaeis uel gentilibus conuersus sit ad fidem, omnino necesse est, ut ab illis nihil inhonestum et turpe impingi possit: ne dicant, Ecce quales sibi Christiani episcopos praeficiunt: quos nos improbos nouimus, hos ipsi tantis dignantur honoribus: nobiscum con temptibilis ob flagitia, nunc honoratur a Christianis. Talis religio, quales haben  \pend
\section*{AD TIMOTH. CAP. III. }
\marginpar{[ p.499 ]}\pstart tur ministri. Sicque in laqueum et in opprobrium calumniatorum incidit ecclesia ipsa, et non solum minister.  \pend
\phantomsection
\addcontentsline{toc}{subsection}{\textit{Diacones similiter pudicos, non bilingues, non multo uino deditos. non turpe lucrum sectantes. }}
\subsection*{\textit{Diacones similiter pudicos, non bilingues, non multo uino deditos. non turpe lucrum sectantes. }}\pstart Episcoporum ministri Diaconi sunt uocati, ut ex Actis discimus, non ut praedicarent uel legerent aut canerent euangelium: sed quia oblationes fidelium dia stribuebant ad ordinationem episcopi, toti seruiebant ecclesiae et ipso episcopo: deque et ipsi oblationibus uiuebant, sicut et episcopi, ut nunc sunt coadiutores pa roecianorum, ut liberius episcopi orationi, lectioni et praedicationi instarent, et ministerio uerbi diuini. Talis erat Laurentius olim, in primis autem Stephanus, cum suis collegis, qui tamen opportunitate monente non tacebant et ipsi a uerbo dei praedicando. Quia autem et ab ecclesia nutriebantur ob ministerium publicum, ut essent ecclesiasticae, hoc est publicae personae, ideo pudicos uult eos esse Apostolus, honestate conspicuos, uenerandae conuersationis in uerbo, facto et exemplo. Ne sinistra quaecunque opinio illorum dehonestaret officium, ne diuidendo collectas eleemosynas, suspicionem cuiuscunque incontinentiae incurrerent, si pauperioribus foeminis dona necessaria prae aliis largirentur. Non bilingues.) Ne quaestus gratia uarie loquantur de Christianis dogmatibus, ne aliud affirmarent coram episcopo, aliud coram simplicibus de ecclesia. Neque uerbum dei adulterent, ad gratiam nil loquantur, ob eleemosynas corradendas, ne uerbis fucatis duplicibusque mercari studeant a uulgo gratiam altioris status. Non multo uino deditos.) Ne luxu quocunque soluantur a debita grauitate, ne dedecori sint episco po ecclesiae officio suo, ne aboperibus pietatis fideles retrahantur, sicut ob filios Heli, ob dispensatorum improbitatem. Non turpe lucrum sectantes.) Non lucrum turpe in eleemosynam expendendum est, sed iuste acquisitae diuitiae, honeste acquisita. Praeuidebat enim in spiritu, bonorum affluentiam irrepturam in perniciem ministrorum ecclesiae, quam interim uidemus et sensimus creuisse in abu sus intolerabiles.  \pend
\phantomsection
\addcontentsline{toc}{subsection}{\textit{Habentes mysterium fidei in conscientia pura. }}
\subsection*{\textit{Habentes mysterium fidei in conscientia pura. }}\pstart Cum passim in domibus, et quasi in angulis instruere obligarentur eos, quos eleemosynis reficiebant diaconi, omnino scire probe debebant mysterium fidei Christianae, et rationem reddere eorum quae ab episcopis docebantur, respondere de articulis fidei, sicut Stephanus, assistere episcopo, defensare, ac ueritatis testimonium praestare pariter ac conuersationis, ubi id utile cernebant, et Christo, honorificum, euangelio et ecclesiae. Certi coram deo, se gloriam propriam non quaerere, sed profectum fidelium de ecclesia, implere quoque factis quod profitebantur uerbis. Quod si deinde minus promoueretur fidei negotium, ipsi tamen candide egissent functionem suam, et fideles in ministerio ecclesiae inuenirentur:  \pend
\phantomsection
\addcontentsline{toc}{subsection}{\textit{Et hi autem probentur primum, et sic ministrent, nullum erimen habentes. }}
\subsection*{\textit{Et hi autem probentur primum, et sic ministrent, nullum erimen habentes. }}\pstart Proximum in ecclesia post episcopum honorem habebant diaconi, tanquam fideles ecclesiae, et episcopi procuratores. Vt hodie fortassis dicerent Decanorum camerarios, sed longe miserabiliore comparatione. Oportebat ergo probari fidem corum, et probitatem, ac sanctam industriam, si illis committi possent tuto distribuendae collectae, si denique idonei uerbi ministri essent, episcopo concordes, fide, doctrina et moribus. Non ideo passim cuiuis haec functio commendari poterat, erant ergo probandi. Derogasset autem quodcunque crimen et inhonestatis diffa\pend
\section*{COMMENT. IN I. EPIST. }
\marginpar{[ p.500 ]}\pstart mia episcopo et ecclesiae, fideles absterruisset a fide, et a misericordiae operibus, ia cturam accepissent ecclesiae pauperes, facileque fraudem suspicari contingeret de oblatis in dispensatore et procuratore quondam notato de crimine: familiam quoque  et statum cum authoritate episcopi dehonestasset.  \pend
\phantomsection
\addcontentsline{toc}{subsection}{\textit{Mulieres similiter pudicas, non detrahentes, sobrias, fideles in omnibus. }}
\subsection*{\textit{Mulieres similiter pudicas, non detrahentes, sobrias, fideles in omnibus. }}\pstart De uxoribus nimirum diaconorum loquitur Apostolus, eadem ratione, qua et prius de episcopis eorumque familia praecipiens. Quia enim de ecclesiae bonis uiuebant et de oblationibus, maxime decebat collegas et consortes diaconorum pudicas esse moribus, uerbis et sensuum omnium honestate, conuersatione maturas, quietas, pacificas, modestas, sine quacunque arrogantia, et dissolutione quacunque in moribus, ne ob easdem ecclesiam turbari et scandaliz ari quocunque  modo contingat. Non calumniosas, rixosas aut discordiam seminantes, delatrices, susurrones uel mendosas, quo minus ecclesia conturbetur. Cum a talibus episco pi et diaconi summa semper diligentla publice dehortari debeant: sobrias quoque , ut dictum est de diaconis et episcopis. Foedissimum est siquidem in muliere uino lentia, et abusus in uictu et uestitu: uinum illis in primis mater est turpitudinis et libidinis. Non debebat uitiorum huiusmodi exemplar relinqui subditorum uxo ribus a foeminis ecclesiasticis. Fideles in omnibus, ut cum ministrorum ecclesiae adiutrices sint, instituto diuino, non secus ac diaconi ipsi, fidelissime dispensent omnia, nihil pro se per fraudem subtrahant, in fide omnia dispensent, ne ponant offendiculum fidelibus in ecclesia unde uiuunt, sed eidem fidelissime ministrent, quantum attinet ad suumofficium.  \pend
\phantomsection
\addcontentsline{toc}{subsection}{\textit{Diacones sint unius uxoris uiri, qui filiis suis bene praesint, et suis domibus. }}
\subsection*{\textit{Diacones sint unius uxoris uiri, qui filiis suis bene praesint, et suis domibus. }}\pstart Non notentur diaconi de quacunq incontinentia, ideo uxores plures quam unam minime illis permittit spiritus apostolicus. Vt enim licuerit eo tempore Chri stianis conuersis de Iudaismo uel gentilltate, uel per legem communem pluralitas uxorum, uel ad minus toleraretur, tanquam reliquiae conuersationis pristinae. Quanquam Romanis id legibus uetaretur, unde et cum tempore abrogata est merito huiusmodi uxorum pluralitas humanis legibus. Non tamen tunc quoque spiritus Christi in apostolis patiebatur in episcopo uel diacono hunc morem, maxime quod in grauamen non modicum ecelesiae cessisset tot personarum educatio, praeter notoriam incontinentiam in ministris dei et ecclesiae. Nihil autem hic de digamia uel trigamia praescribitur ab Apostolo: adolescens enim diaconus uel etiam integrae aetatis episcopus, nusquam praecipitur prima mortua uxore, secundam non ducere, ut perpetuo irremediabili molestia ureretur, Sicut nec iuuenculas uiduas uetabat a matrimonio, sed magis nubere illas praecipiebat Christus, qui in Apostolo loquebatur. Et hic est eadem ratio, ne satanas tentet et uincat. Bene autem praesint filiis, scientia, moribus, disciplina pietatis, erudiant suos filios diaconi, ut sint speculum honestatis caeteris adolescentibus de ecclesia, ad morum reformationem in toto populo. Suis quoque domibus.) Omnem uult familiam, ministrorum ecclesiae quorumcunque honestissime gubernari, eadem ratione qua et episcopos. Id quoque intelligendum de omnibus, qui quomodocunque praesunt ecclesiae, etiam in temporalibums. Quandoquidem laicorum et clericorum distinctio non legitur in apostolicis literis, sed Christianorum et gentillum. Maxime uero hic discimus, quales in moribus sint fide et exercitio, quotquot ecclesiae cen  \pend
\section*{AD TIMOTH. CAP. III. }
\marginpar{[ p.501 ]}\pstart sibus nutrluntur, qui scilicet in ministerium ecclesiae necessarii sunt. Nemo enim praeter ecclesiae ministrantes, abeadem sustentari debet, et omnino nemo otiosus, qui possit sufficere laboribus idoneis. Ex quibus enim prorsus nulla redit ecclesiae utilitas, qui nullo sunt usui in opere sancto, a Christo et apostolis instituto, nec usui futuri sperantur, hi non sustententur ecclesiae stipendio. Sunt autem ecclesiae necessarii episcopi, qui scilicet uerbum dei tractant et docent. Diaconi, qui discunt et cooperantur episcopis. Et hii scholastici, qui olim Lectores dicti, iuuentuti instituendae praesunt. Cantus horarum canonicarum, ut nunc fit, est prorsus inutilis: possent tamen pro eo commodius institui exercitia circa canonicas scripturas utiliter tractandas. Monachi omnes laborent manibus, olim magna illorñ pars libris scribendis sese nutriebant. Hodie apta essentmonasteria ad libros opti mos imprimendos: et si insuper succisiuis horis uerbo dei utiliter insistere uellent, ut deberent, uel unum et alterum psallere Psalmum, id laudabile foret. Tantum autem temporis die noctuque usq ad ingens fastidium murmurare Psalmos, et nescio quae responsoria canere, miror quis instituerit, nisi ob redimendum tempus et taedium orandi uel legendi. Et ne toti otiosi et inutiles ecclesiis, ecclesiasticorum bonorum uoratores iudicarentur, tandemque illis fraudarentur: temporales quidem possessiones in hoc saeculo pauperibus dispensasse, conforme fuerit euange lio, otium autem sectari, et omnem cuiuslibet rei uel minimam etiam egestatem memdicitatibus assiduis studiosius resarcire, quis probet? uel non dignissime rideat, imo damnet ex scripturis canonicis et uerbo dei. Laborandum est potius illis, eleemosynasque dare et facere studeant hi religiosi homines, quam requirere et mendicare. Crucem Christus sectatoribus suis imponit, non iners otium illud, et sine omni propria sed non aliorum solicitudine, omnium rerum necessariarum suppetentiam admirabilem, ex laboribus pauperum, et usura diuitum sectari. Haec certe sanctus Franciscus non docuit, sed laborem manuum, pro quo reciperent uitae sustentationem, sine affluentia, pretiositate uel curiositate, sed ut in omnibus non tam reluceret, q̃re uera esset summa rerum paupertas et humilitas, subiectio Chri stiana, et apostolica charitas ministrandi infirmis, etiam leprosis. Longe autem ali, ter cessit institutum boni uiri, q̃ prae se tulerat, sancto quodam spiritu instigatus, ad praedicationem suos hortabatur sancti euangelii, non solum uerbis, sed multo magis exemplis. Non tam denarios abhorruit quam auaritiam habendi, et luxum sectandi ac fastum. Inculcauit usque ad mortem suam laboritium, et humilitatem, fidem sanctae Romanae ecclesiae, et euangelii praedicationem, sine omnibus exemptionibus Papisticis et philosophicis studiis, quorum sectatores iusto zelo maledixit non semel. Quid autem profecerit, orbis indicat, etiam si sui fratres nondum uidere uelint, qui tamen aliorum comparatione minus fortassis aduersant ueritati, et resipiunt, ni fallar, copiosius. Si quis alius ecclesiae etiam in temporalibus ne cessarius minister decernendus sit, ille de ecclesia uiuat, nemo inficiabitur. Otiosorum autem prorsus nemo dignus habeatur, nec feratur. Qui non laborat, nec manducet, caeteris paribus, si scilicet laborare possit: quin et qui nesciunt, discant Iaborare, non pro pretio laboris reponendo, sed ad obediendum euangelio, non solum ad repellendam otiositatem, ut habet regula Minorum. Egens necessariis nemo sinatur in ecclesia, sed aeque non otiosi et ualidi mendicantes ferantur.  \pend
\phantomsection
\addcontentsline{toc}{subsection}{\textit{Qui enim bene ministrauerint, gradum bonum sibi acquirent, et multam fiduciam in fide, quae est in Christo lesu. }}
\subsection*{\textit{Qui enim bene ministrauerint, gradum bonum sibi acquirent, et multam fiduciam in fide, quae est in Christo lesu. }}\pstart Diaconi alicuius ecclesiae si ministerii sui functionem laudabiliter cum probi  \pend
\section*{COMMENT. IN I. EPIST. }
\marginpar{[ p.502 ]}\pstart tate, ac fideliter absoluunt, non solum laudem, sed et gradum sibi meliorem acquirent. Erunt proculdubio idonei ad succedendum episcopo, si non ei cuius diaconi sunt, transferentur saltem, eligendi in aliam ecclesiam. Vel ex cooperatoribus fient fortassis coepiscopi, uel in una ciuitate pro multitudine populi episcopi, id est parochi plures constituentur. Assequetur apud uniuersam Christi fidelium con gregationem authoritatem magnam, multamque fiduciam fideles collocabunt in eum, pro insigni fide, dexteritate, diligentia et uigilantia, qua in sua familia uide bitur sedulus. Libertatem ex hoc sibi permittendam ab episcopo acquiret, ut omnia uel saltem multa pro suo arbitratu sinatur agere, toties repertus integer et fidelis in Christi obsequio, ut nemo sibi fidem agendorum abrogare possit.  \pend
\phantomsection
\addcontentsline{toc}{subsection}{\textit{Haec tibi scribo fili Timothee, sperans me uenire ad te cito. }}
\subsection*{\textit{Haec tibi scribo fili Timothee, sperans me uenire ad te cito. }}\pstart Ex Laodicia scribebat Ephesum, nuper institutum ibidem eum episcopum consolando et confortando, ut sciret ad id temporis magis necessaria scitu, suo tem pore personaliter de omnibus coram edocturus, quae circa episcopi et diacono. rumofficia essentordinanda in ecclesia.  \pend
\phantomsection
\addcontentsline{toc}{subsection}{\textit{Si autem tardauero ut scias quomodo oporteat te in domo dei conuersari. }}
\subsection*{\textit{Si autem tardauero ut scias quomodo oporteat te in domo dei conuersari. }}\pstart Domus dei non templum intelligatur marmoreum, in quo cum thuribulo et cereis circumcursare debeat de ara ad aram, et pontificalia exercere muta uel ca nora: non est enim pontificale illud ceremoniarum ex traditione apostolica, sed ut patet collatione uetustorum pontificalium, secundum tempus semper auctum, quasi timuerint pontifices redargui, si non multa in domo dei agenda curarent, quo minus unum illud necessarium, uerbi dei praedicatio, et canonicarum scripturarum enarratio et inculcatio, ab illis indoctis hominibus exigeretur. Haecenim, ut probe Moria Erasmi admonuit, reiicienda erant in monachos otiosos, do nec et ipsi Suffraganei crearentur. Domus dei est quaelibet fidelium congregatio, Christo synceriter credentium, et uerbum dei toto pectore complectentium, et per opera fidei proficientium semper. Inter quos, non super quos, conuersatur episcopus apostolicus, eos docendo et instituendo in fide Christi, in doctrina san na et catholica.  \pend
\phantomsection
\addcontentsline{toc}{subsection}{\textit{Quae est ecclesia dei uiui, columna et firmamentum ueritatis. }}
\subsection*{\textit{Quae est ecclesia dei uiui, columna et firmamentum ueritatis. }}\pstart Domus illa dei in qua te ceu ministrum et superintendentem dominus constituit, omnis uidelicet fidelium congregatio: non tantum Ephesina, sed ubique eodem spiritu congregata sanctorum multitudo. Ipsa estecclesia, quam deus condi dit, gubernat, conseruat, ac docet spiritu oris sui, uerbo dei. Non enim per se, sed in Christo subsistit ac durat. Basisque est et sedes ueritatis, colitur em̃ in ea et agno scitur unus uerus solus et uiuus deus. Cum omnes aliae ecclesiae et hominum congregationes, deos falsos et nihili colant sua idola. Iudaei quoque a ueritate dei, quae est per Christum Iesum, aberrauerunt, et excaecati sunt, ac cecsderunt. Nostrae autem ec clesiae doctrina et fides fundamentum est, basis et columna immobilis et semper mansurae ueritatis. Ipsi enim errauerunt et ceciderunt, nos autem surreximus et erecti sumus, credentes in Christum dominum, et deum uiuum ac uerum, qui est uia, ueritas et uita:illam quoque ueritatem firmissime credimus, praedicamus, tene mus, et usque ad mortem defendimus.  \pend
\phantomsection
\addcontentsline{toc}{subsection}{\textit{Et manifeste magnum est pietatis sacramentum, quod manifestatum est in carne, et iustificatum est in spiritu, apparuit angelis, praedicatum est}}
\subsection*{\textit{Et manifeste magnum est pietatis sacramentum, quod manifestatum est in carne, et iustificatum est in spiritu, apparuit angelis, praedicatum est}}
\section*{AD TIMOTH. CAP. III. }
\marginpar{[ p.503 ]}
\phantomsection
\addcontentsline{toc}{subsection}{\textit{gentibus, creditum est mundo, assumptum est in gloria. }}
\subsection*{\textit{gentibus, creditum est mundo, assumptum est in gloria. }}\pstart Quanquam priora intelligi possint, ut diximus, de ecclaesia et columna. Sima plicior tamen, et aeque indubius sensus uidetur esse, et plane insignis, si sic literam ordinemus, ut quod praescriptum est, de columna et firmamento ueritatis, referatur ad ea quae illico sequuntur:ut sit is sensus, sic literam construendo: Vtscias, inquam, quomodo te oporteat conuersari in ecclesia illa tibi commissa. Columna autem et firmamentum ueritatis, et manifeste magnum, quod scilicet citra omnem controuersiam magnum est pietatis sacramentum, hoc est, quod nune tibi paucis prodam. Quasi dicat: Totius Christianae doctrinae et catholicae uerstatis columna, cui innititur, et firmamentum cui insistit, et mysterium summum pietatis et fidei: hoc unicum est, quod et in ecclesia praedices, fidei esse summam ac basim. Deus manifestatus est in carne: uel si in neutro genere interpreteris, myste rium pietatis diuinae in Christo nobis exhibitum est, hoc pacto. Cum olim multia phariam multisque modis apparuisset dominus deus, uel uerbum dei, illi consubstantiale et coaeternum, in forma angelica, specie humana, rubo igneo, spiritu, fa cibus et nube, etc. nouissime diebus istis nouissimis uel proximis loquutus est no bis in carne humana, assumpta in unitatem personae, sine commixtione naturarum, uel mutatione, in qua assumpta humana natura a uerbo, manifestauit se deus uerbum mundo, incarnationem suam miraculose perficiens, et habitu inuentus ut homo sermo dei, in Christo domino, ut inuisibilis deus, hominibus uisibilis apparens, cum hominibus conuersaretur in terra, carne protoplastis indutus, et ha bitu inuentus ut homo. Iustificatus in spiritu.) Per spiritum sanctum datum ho minibus, iustus deus apparuit, fide et gratia, per uerbum suum iustificans electos, qui in spiritu sancto uoluit omnes peccatores iustificari, quos omnes reos in primo parente damnauerat, hos in Christo per spiritum sanctum iustificat, et ipse iu stus innotescit ac pius, qui seuerus prius iudicabatur et durus, ut scilicet iustifice tur in sermonibus, et uincat, cum nostro carnali sensu iudicatur iniustus. Apparuit angelis.) Dei gratia, mysterium pietatis, ipsum euangelium gratiae diuinae ad homines, a saeculis magna ex parte ignotum angelis, per Christum hoc tempo re innotuit principatibus ac potestatibus in coelestibus, qui hactenus ignorauerunt tantam dei ad humanum genus pietatem, et in Christi mysterio didicerunt, quae prius abscondita illis erant et incognita. Quae nunc per euangelii praedicationem, et per ecclesiam, angelis quoque clarior innotuit, Ephes.3. Praedicatum est gentibus.) Omnibus populis, gentibus ac linguis hoc pietatis mysterium, et euangelium regni praedicatum est, delata est haec stupenda dei gratia omnibus hominibus, nemine excluso. Christus saluator datus est benedictio omibus gentibus, pignus et arra salutis, non solum Iudaeis, sed et gentibus coepit per apostolos praedicari hoc euangelium salutis: et praedicatum est nec solum per apostolos et apostolicos uiros, sed et Creditum est mundo.) Non solum Iudaeis, sed uniuerso orbi, contra carnem contranitentem et sanguinem, cui hoc mysterium non reuelatur, nisi a patre coelesti, qui trahit ad se per Christum quos elegit, et donat fide, quos iustificare decreuit per fidem in Christo. Assumptum est in gloria.) Glori ficatum est, et collaudatum ubique , et diffusum in omnem terram: et post agnitum mysterium etiam glorificatum est, Christo ascendente in coelum, et ad dexteram sedente, et euangelium mundo illucescente cum gloria, qua omnis mundana po tentia, sapientia et iustitia oppressa iacet, ut illud solum et unicum gloriosum sit in coelo et in terra. Hoc fidei tantum mysterium, et pietatis diuinae sacramentum, sedes et basis est omnis ueritatis, columna et fundamentum, cui innititur omne  \pend
\section*{COMMENT. IN I. EPIST. }
\marginpar{[ p.504 ]}\pstart quod sçimus, credimus, operamur, et speramus. Hoc si quis credat, ueneretur, et amplectat, iustus est et perfectus Christi seruus et Christianus. Hanc, inquam, summam doctrinae euangelicae, scias, teneas, doceas, et studeas inculcare, Christianae scilicet fidei breuissimos articulos, Deum scilicet incarnatum, uel in carne humana manifestatum, per spiritum adoptionis, in quo clamamus, Abba pater, in fide omnia iustificata, non lege carnalis literae uel scripturae. Angelis etiam admi rabilem, quos non apprehendit, sed semen Abrahae, collatam scilicet hominibus tantam dei gratiam, quae post factum illis innotuit, uel dum fieri primum coepit. Quod euangelium Christi, et reconciliatio tam perfecta, praedicata est gentibus omnibus. Mundo quoque creditum est, domino miserante, et hac credulitate ac fide iustitiam omnibus largiente. Assumptum quoque est hoc pietatis euangelium uel beneficium in gloriam dei, et patriarcharum sanctorum, qua nemo fraudetur fidelis in coelo, et toto orbe terrarum. Tale aliquid et prius aeque breuibus descripsit Apostolus, cum diceret, Vnus deus, unus etiam conciliator dei et hominum, homo Christus Iesus, qui dedit semetipsum pro nobis pretium redemptionis pro omnibus. Haec, inquam, est columna fidei, et euangelium regni: praeter haec alia non sunt quaerenda, non docenda, non imponenda fidelibus, ut necessario credantur. Si quis his con uenientiora uel consequentiora potest, abundet is suo in spiritu.  \pend
\endnumbering\beginnumbering\section{CAPVT IIII.}
\phantomsection
\addcontentsline{toc}{subsection}{\textit{\huge\textbf{S}\normalsize Piritus autem manifeste dicit, Quia in nouissimis temporibus discedent quidam a fide. }}
\subsection*{\textit{\huge\textbf{S}\normalsize Piritus autem manifeste dicit, Quia in nouissimis temporibus discedent quidam a fide. }}\pstart Spiritus Christi, qui in discipulis adhuc manet acceptus, uel in die pentecostes, uel in impositione apostolicarum manuum, per quem et futura prae nouerunt, ac prophetant de Antichristis, qui a Christo abstrahunt, et a fiducia in redemptore nostro nobis donata, is spiritus dei reuelauit prophetiam de futuris temporibus, post haec nostra tempora esse futurum, ut parum curaturi sint homi nes solo nomine Christiani, de fide in Christum, non agnoscent mysterium diuinae pietatis illius, putabunt se non tantum in Christi spiritu iustificatos per fidem. Nec erit gloriosum apud eos, quod deus euangelium gratiae suae uniuersalis transmiserit illis, omnibus scilicet electis: sed in operibus potius legis iustificari uolent, si non legem Mosis solam, sed et plus quam Mosaicae legis praecepta seruanda proponant, imo et contraria. Cum Christus fidem in deum et perfectam fiduciam docuerit esse medium et instrumentum uerae iustificationis, non ex operibus, quae nihil ualeant ad ueram pietatem. Docturi sunt enim obseruantiam in delectu ciborum, temporum, uestium, locorum, uotorum, meritorum, coelibatus, inuocationes quoque sanctorum mortuorum, opera etiam ieiuniorum, peregrinationum laborumque multorum et Missarum, quibus in rebus iustitiam collocantes, gratiam dei per Christum minus attendent, uel appretiabuntur. Posteriora tempora melius dicerentur quam nouissima, ne quis putet in fine mundi haec olim futura, quum etiam apostolorum tempore, et deinde semper pseudapostoli et sophistae supersti tiosique ac haeretici haec crediderint, seruauerint et docuerint. Desciscere autem a fide intelligitur, non a fide unius dei, sed a Christianismo, a uerbo dei, ab agnitione salutis, per Christum nobis paratae, et a libertate spiritus filiorum dei, a cultu dei, quo uere dominus uult coli, in spiritu, fide, spe et charitate, uereque sanctis operibus   \pend
\phantomsection
\addcontentsline{toc}{subsection}{\textit{Attendentes spiritibus erroris, et doctrinis daemoniorum, in hypocrisi loquentium mendacium, et cauteriatam habentium conscientiam. }}
\subsection*{\textit{Attendentes spiritibus erroris, et doctrinis daemoniorum, in hypocrisi loquentium mendacium, et cauteriatam habentium conscientiam. }}\pstart Cum ecclesiastici tam doctores que auditores ipsissimo deiuerbo, et Christi do  \pend
\section*{AD TIMOTH. CAP. IIII. }
\marginpar{[ p.505 ]}\pstart tirmae, quae ueritas est, adhaerere deberet, ipsi erroris spiritui et humanae praesum ptionis superstitiosis traditionibus attendunt. Quicquid enim a uia exorbitat, erae roneum est. Spiritus autem Christi uid est, cuius uitam, mores et uerba aemulantes, non errabimus unquam. Si quid praeter illius doctrinam et exempla unquam uel a quocumque tentabitur, sine dubio abducet a uera pietate ad morem gentilium impiorum, uel Iudaeorum nimiam superstitionem. Doctrinis daemoniorum)ait, Daemones legimus docuisse primum quidem, ut uerbo dei Eua non crederet, sed inuerteret sensum, augeret ac minueret, suo proprio sensu abundaret, praeter dei praescriptum sciret et ageret. Docuit idololatrarum morem, sacrificia, cultum et coelibatum, uestales uirgines, eunuchos suos, et gallos parauit, consultationes di uinorum, et caetera huiusmodi. Christo quoque suadere tentabat de cibo, de arrogantia, de superstitione, et idololatria. Mendacium in ore falsorum prophetarum, obstrepere ueris prophetis, ad uere sancta opera prouocantibus. Excelsorum et ararum multitudinem, perpetua uota deuotorum uariis idolis, ut is Marti, alius Veneri, Mercurio et Dianae, propriis quibusuis ceremoniis, et legitimis deserui rent, secundum loca, tempora, ac personas. Tales omnes doctrinae sunt daemoniacae et perniciosae, deoque odibiles, et contrariae fidei in deum per Christum. Diuinae autem et sacrosanctae sunt, discere mititatem, humilitatem, fidem, spem, charitatem, modestiam, tolerantiam, beneficentiam in amicos et hostes, uincere in bono malum, etc. Quicquid his obuiat, daemoniacum sit necesse est et spirituum men dacium. Loquentium mendacium per simulationem.) Qui religiosiores uideri uolunt, excogitantes de proprio doctrinas, et regulas uiuendi speciales, quasi Christus insufficienter docuerit, et cum apostolis non satis religiose et sancte uixerit, ut praeter euangelicam traditionem, alia quaedam illis ordinanda et instituen da commiserit, cum ea tales superstitiosi de suo corde confinxerint, foris uolentes uideri sanctiores, cum fuerint intra conscientiam repleti spurcitiis, hypocritarum more, qui plura ostentare uolunt sanctitatis indicia, quam iusserit Christus; interim necessaria obmittentes quae mandauit Christus, docentes, non litigare uerbis, nec alios iudicare, esse mites, pacificos, modestos, mansuetos, et humiles, et similia con silia esse euangelica, non necessario obseruanda, sed nudis pedibus uel scissis incedere calceis, sextis feriis ielunare suo modo, pecuniam non tangere, prodefun ctis orare, et similia, docent esse necessaria ad salutem suis professoribus. Sic enim summi pontifices Romani, Gregorius nonus, Nicolaus tertius, et Martinus quin tus declarant ac docent. Ad consilia euangelica obseruanda Franciscanos non ob ligari, nisi in ipsorum regula praeceptorie uel inhibitorie, seu sub uerbis aequipollentibus sint expressa, ut patet in Decretali, Exiit qui seminat, libro sexto. Similiter Clemens in Decretali, Exiui de paradiso. Consentientibus in his et similibus quatuor magistris Parisiensibus, etc. Similia et maiora confinxerunt sibi alii superstitiosi circa esum carnium. Sed de potu uini nihil uspiam legitur modificatum uel inhibitum. Cum itaque obmittant quae Christus docuit et praecepit, supererogationis tamen consilia sibi praefigunt, quae tamen nec ipsa complent. Vnice hactenus studuerunt perpetuo constitutiones figere et refigere, de euangelica autem perfectione nihil prorsus dignum sentire, uel etiam perficere. Et cauteriatam habentes conscientiam.) Conscii sibi sunt huiusmodi doctores, quod quae alios docent facere, ipsi non possuntuel nolunt, uel necessaria esse non iudicant, quantumlibet aliis perficienda iniungant quae Christus non docuit. Ea tantum ratione introducentes plurima, supra praeter uel contra Christi doctrinam, ut aliquid emungatur, ut territi uel angustiati in conscientiis, promptius ad uota obediant,  \pend
\section*{COMMENT. IN I. EPIST. }
\marginpar{[ p.506 ]}\pstart magis praelatos honorent, et ut sanctuli uideantur, fingentes se agere quod alios docent uel praeciplunt. Omnibus his modis conscientia talium cauteriata deprehenditur, iudicio autem percutientis serui, et mactantis in absentia domini conseruos suos, durissime iudicabuntur, ut portio eorum futura sit cum infidelibus, etc.  \pend
\phantomsection
\addcontentsline{toc}{subsection}{\textit{Prohibentium nubere, abstinere a cibis, quos deus creauit ad percipiendum cum gratiarumactione fidelibus, et his qui cognouerunt ueritatem, quia omnis creatura dei bona est, et nihil reiiciendum, quod cum gratiarumactione percipitur, sanctificatur enim per uerbum dei et orationem. }}
\subsection*{\textit{Prohibentium nubere, abstinere a cibis, quos deus creauit ad percipiendum cum gratiarumactione fidelibus, et his qui cognouerunt ueritatem, quia omnis creatura dei bona est, et nihil reiiciendum, quod cum gratiarumactione percipitur, sanctificatur enim per uerbum dei et orationem. }}\pstart Quid illi daemoniaci praedicatores, et cauteriata conscientia doceant per ordinem designat. Primo enim docuerunt non quidem ut nemo matrimonium suscipiat, hoc enim fieri non potuisset, et contra totam sacram scripturam. Nec docent peccatum esse, omnibus coniugem habere, sed tantum religiosioribus Chri stianis, quibus eatenus coelibatum persuaserunt, ut perpetuo seruandum uouerent, quasi rem turpem matrimonium habentes, et indignam tractantibus caetera sacramenta, quasi sanctum sanctiori couti non possit, quasi sacerdos sine peccatomagno coniugem habere non possit. Huius doctrinae fundamentum cum nullum sit in sacra scriptura, cumque multorum malorum causa sit, ut misera docuit experientia, et contra Christi et apostolorum doctrinas, merito dicitur diabolica, Cum nubere in domino, deus, natura, scriptura, ufus, necefsitas, et honestas omnia approbent, tanquam sanctum, et deo placitum, nullis prohibitum: neque aliquibus non necessarium, nisi specialissima et rara gratia spiritus praeuenientis, ut non urantur. Prohibitae autem nuptiae sub uoto, uel sine uotis, domino non placentibus, ineffabilium malorum causam praestant, quorum uotorum nemo nisi diabolus pater et patrator causaq esse potuit. Quod et apostolus Paulus manifeste sci uit, docuit, scripsit, ut anathema sit, qui contrarium docet. Abstinere a cibis etc.) A crapula abstinere et ebrietate, non onerare corpus et spiritum, non reddere imbecillum ad necessarias ac utiles functiones deus praecepit. Iudaeis denique genera ciborum quaedam prohibuit pro tempore legis Mosaicae, cum reliquis obseruan dae, ad pacificam cohabitationem, ad abominationes cauendas, ad praeueniendum aduersam corporum ualetudinem, quae alioqui non proficiunt omnibus complexionibus: nec enim omnia ideo usuunt, ut homini sint in cibum, sed etiam quaedam data sunt in remedium. Quaedam cibus ordinata sunt eorum animalium, quae homini data sunt in escam: quaedam nauseam naturalem adferunt de genere car nium: Quaedam ob idololatrarum superstitionem, quam remotissime abesse suis uoluit deus cultoribus, et seriose prohibuit. Nunc autem in nouo testamento spi ritus sanctus reuelauit, et Christus docuit, omnia munda mundis: nihil quod intrat in os, coinquinat hominem. Cibi itaque non insani, et nauseam non habentes annexam, de se immundi non sunt, si cum gratiarumactione accipiuntur, deoque ́ donante in escas hominibus, omnes comedi possunt, nemine prohibente. Si tamen ad corporis castigationem, et ad tuendam ualetudinem, sectandamque parsimoniam quisquam sine superstitione libere abstineat, tanquam ob nauseam uel medicinam, non quidem peccabit. Tantum ne immundum reputet et illicitum, uel coronae alicuius meritum arroget sibi, uel speret huiusmodi abstinentia. Non enim refert quid uel quando comederis, sed quantum. Possumus abstinentiam sectari non solum in piscibus uel leguminibus, sed etiam in carnibus: potest leguminum esca peccari, potest omnis generis carolicite manducari si libeat, cum fide et temperantia. Quicquid comedimus, cum gratiarumactione manducemus.  \pend
\section*{AD TIMOTH. CAP. IIII. }
\marginpar{[ p.507 ]}\pstart Agnoscentes largitorem omnium bonorum, qui omnia nobis in usum bonum con tulit, si non ad libidinem et superfluitatem utamur, et contra mandatum suum. A principio enim herbas concessit, post diluuium carnes addidit, sanguinem ne gauit, postea animalia quaedam immunda prohibuit, per Christum autem omnia mundauit, comessationem quancunque prohibuit, ac regnum dei dixit non escam et potum, sed fidem et charitatem. Quod addit, Fidelibus, et his qui cognoueruntueritatem, contra Iudaeos et Iudaizantes non solum loquitur, quorum iam haeresis ut nota damnabatur, sed contra futuros praedicatores, qui docturi essent aliquos cibos esse malos, religiosiusque esse uno cibo uesci quam alio: deum placari esu huius cibi, non alterius: offendi illum mortaliter, hoc tempore comedenti hunc ci£ bum, alio tempore non sic. Introducturi quoque uota contra huiusmodi escas, tanquam cultum dei, et deo grata, ueluti domino placeret magis huius cibi esus quam̃ al terius. Quanquam uero unum cibi genus magis impediat ingenium, et sacras meditationes uel eruditas q̃ aliud, nullus tamen cibus id facit, nisi immodice sumptus, id quod iure diuino uetat de omnibus cibis. Sunt igitur omnes tales doctrinae fal sae, contra deum, adinuentiones humanae, et daemoniorum. Omnis autem cibus et semper et omnibus bonus est, si cum gratiarumactione ad naturae sustentarionem, habita ratione ualetudinis, sumatur. Omnis autem cibus semper malus decernitur, si nimius et ad libidinem capiatur, ut uotum de hoc uel alio cibo perpetuum, uel ad regulare tempus, assumi non debeat. Non estut abstinentia uoueatur, quae semper est necessaria, ut semper illicita est superfluitas. Nec deo placere potest, ut necessariis naturae detrahatur, ad legitimas functiones impediendas: obligamur semper et omnes ad moderatum usum cibi, et omnium bonorum temporalium. Non oportet uouere uel abnegare, quod per uerbum dei nobis prohibitum alioqui scimus, quodque dei uoluntati resistit. Nunquam autem cibus sumen dus estsine gratiarumactione. Christus cum discipulis hymnum dixit et laudes post cibum sumptum, et exactam coenam. Si quid immundum uel insanum in ci£ bo sobrie sumptofuisset, oratione et uerbo dei audito et credito, sanctificar̃, fitque ́ remediabile ac innoxium: scienter tamen et temere cibum sumens perniciosum, et aduersum bonae ualetudini, peccato non caret. Non est deus tentandus. Non omnia omnibus conueniunt. Sed eo modo sacrae literae non loquuntur. Denique  non omnibus est omnis creatura bona, nec omnibus sanctificatur cibus per uerbum dei et orationem, sed tantum fidelibus, his scilicet qui cognouerunt ueritatem: ignorans enim omnem ereaturam esse bonam, et in conscientia sua iudicans malam, et non abstinens tamen ab ea, peccat, agens contra fidem. Nemo igitur contra conscientiam instigandus ad cibum, sed prius instruendus est, ut fidelis ueritatem agnoscat. Imo et cauendum est a talium scandalo, si ex infirmitate, et non ex contemptu aut temeritate credere nondum possint, ut non est in libero homis num arbitrio credere quicquid uelit. Si enim esca scandalizat fratrem, non manducabo carnes etc.  \pend
\phantomsection
\addcontentsline{toc}{subsection}{\textit{Haec proponens fratribus, bonus eris minister Christi lesu, enutritus uerbis fidei et bonae doctrinae, quam assequutus es. }}
\subsection*{\textit{Haec proponens fratribus, bonus eris minister Christi lesu, enutritus uerbis fidei et bonae doctrinae, quam assequutus es. }}\pstart Bonus minister Christi est et ecclesiae, qui bona et bene praedicat, nutritus in uerbis fidei, bonaeque doctrinae. Verbum autem fidei euangelium est, et quicquid in tota sacra scriptura et canonica continetur. Quae nos docet fidere, sperare et amare deum, Christum praedicat, uel futurum uel exhibitum, in quem credere iu stitia fidei sit. Bona itaque doctrina est, quicquid ex sacris scripturis credendum, spe\pend
\section*{COMMENT. IN I. EPIST. }
\marginpar{[ p.508 ]}\pstart randum ac operandum docetur, et in ipsius sanam ac claram intelligentiam diri git. Non illae gentilium tantum literae, quae abusiuo nomine bonae literae aestiman tur, quae nec bonum reddunt, nec ad bonitatem instituunt. Haec autem doctrina de cibis non diiudicandis, et generaliter bonis, fratribus proponenda iubetur, fidelibus scilicet, ut primo fratres sint, fideles sint, credant Christum legis perfectionem, tunc em̃ et non prius erudiri poterunt de cessatione legalium per euangelium. Non enim Iudaeis cessant legalia seruanda, nisi credant in Christum: alioqui stante infidelitate, magis peccabunt legem negligendo contra suam conscientiam, quam exuere non possunt uel uolunt. Sunt ergo primum de Christo docendi, postea de gratia et fide, sic tandem de libertate Christiana. Hoc ordine omnes superstitiosi docendi sunt, ut non ficte sese credere iactent. Bonus igitur minister est, qui bona et uera, iusto et apto ordine praedicat, ad capacitatem et fructum fidei plantandae. Malus autem minister erit, qui nondum assequutus est doctrinam fidei et pietatis, qui non est nutritus in uerbo dei, quantumlibet totam philosophiam et omnes linguas, ac alioqui plurima edoctus. Qui non Christi Iesu, sed suam gloriam et utilitatem sectatur, qui quaerit quae sua sunt, non quae lesu Christi. Qui docet non nubendum, qui suadet fouendam perpetuam castitatem, quod in uerbis fidei nusquam continetur: qui delectum ciborum magnifacit, caeterasque ceremonias necessario implendas iniungit, a Christo non institutas: qui fidei simpli cis uerba non docet: qui proiicit margaritas ante porcos scienter. Hi tales mali ministri non sunt audiendi. Quam, ait, assequutus es.) Quam a me didicisti, non nouum ali quid comminiscens, non nouas quaestiones inquirens, non noua mandata proponens, sed quae a me didicisti, uerba fidei, haec diligenter fratribus, hoc est piis hominibus inculca. Non enim uel ipse Paulus audebat docere quicquam, quod non in se docuisset uel reuelasset Christi spiritus. Nostri autem non sic agunt, Bre uiter, euangelii particula recitata uel praelecta, statim sua commenta ad longum ingerunt, nisi quod nunc dei gratia tales cum taedio audiri incipiunt a fratribus, qui semel gustare inceperunt doctrinam fidei.  \pend
\phantomsection
\addcontentsline{toc}{subsection}{\textit{lueptas autem et aniles fabulas deuita. Exerce autem teipsum ad pietatem, nam corporalis exercitatio ad modicum utilis est, pietas autem ad omnia utilis est, promissionem habens uitae quae nunc est, et futurae. }}
\subsection*{\textit{lueptas autem et aniles fabulas deuita. Exerce autem teipsum ad pietatem, nam corporalis exercitatio ad modicum utilis est, pietas autem ad omnia utilis est, promissionem habens uitae quae nunc est, et futurae. }}\pstart Didicerat a Iudaea matre sua Timotheus multas fabulas, quae et adhuc extant apud Iudaeos, quas et modo narrant filiis, praeter ea quae nunc quoque cabala aliqui dici uolunt, et pueris ac puellis legenda permittunt. Sicut et nos sanctorum legendas quasdam et apocrypha. Vult autem Paulus nihil horum in ecclesia publice praedicari, sed tantum uerba fidei, et sacras canonicas scripturas, populumque  dei certissima ueritate, non dubiis fabulis imbui. Non sunt pueris nudae et non explicatae poetarum fabulae proponendae, etiam Latinitatis discendae praetextu, sed fidei et honestatis fomenta ac disciplinae. Multo minus in ecclesia praedicandum aliquid est, nisi canonicum, probatum ac certum. Exercitium pietatis est in uero dei cul tu. Is autem uere deum colit, qui fide et charitate fert̃ totus in deum, et ex ea fide fert in omnia charitatis officia proximo impendenda. Is deum colit, qui dei mandatis obsequitur, uerbum dei reueretur in conscientia, et comprobat conuersatione religiosa. Non qui saepe et multum uersatur in humanis quibuscunque traditionibus, sed qui deum et proximum diligit. Haec uera est pietas, hic uerus aclegitimus dei cultus, haec religio Christiana. In his se exerceat episcopus suamque ecclesiam, discen do sacra, atque docendo uerbo et exemplis. Non in cantibus assiduis, spiritu scili\pend
\section*{AD TIMOTH. CAP. IIII. }
\marginpar{[ p.509 ]}\pstart cet et sono uel organis: non ceremoniis, uestibus, ornatibus, deo se placere arbitretur, sed in pietate. Si scilicet plenus fide, et sacrarum literarum peritia spiritua li, instruat populum. Si solicite incumbat, ne quis ecelesiasticus erret in fide et moribus, ne quis inuidiam seruans, charitatem laedat in ecclesia, nemo proximum in aliquo circumueniat, uel incommodet uerbo uel facto, nemo egeat necessarioui ctu et amictu. Haec sunt pietatis officia uere apostolica, ad omia ualidae, istae obla tiones domino placentes, his sacrificiis promeretur deus. Hic cultus dei uere reli giosus, et uita Christiana, digna ecclesiae, et necessaria in primis episcopo. Corpo ralem autem exercitationem ad modicum utilem, uult apostolus, castitatem tantum carnis, et uirginitatem corporis, ieiunium a cibo statis temporibus, et Iudyi co more superstitiosas uigilias, peregrinationes, et cantus, precularum numerum et ordinem: quae tamen in casu prodesse, nonnunquam possunt obesse, ut et alia multa. Pietas autem, de qua diximus, grata deo est et hominibus sacra, prodest nobis ac proximo, non est inanis et infructuosa. Quicquid autem in gloriam dei et proximi utilitatem referri non potest, uanum est, et exercitatio corporalis dicitur, non spiritualis. Promissionem huius uitae habet, Quia eadem mensura qua mensi fuerimus, remetietur nobis. De operibus misericordiae in extremo iudicio laudabimur, et praemia reportabimus. Fides in Christo hic nos beatos reddit, et in futuro. Charitas proximi impletio est totius legis, quae omnino non praeiudica re debet charitati. Nec exercitatio corporis ceremonialis quaecunque , nec uotum nec promissum uero diuino cultus praeferri debet nec praeiudicare.  \pend
\phantomsection
\addcontentsline{toc}{subsection}{\textit{Fidelis sermo, et omni acceptione dignus. In hoc enim laboramus et maledicimur, quia speramus in deum uiuum, qui est saluator omnium hominum, maxime fidelium. }}
\subsection*{\textit{Fidelis sermo, et omni acceptione dignus. In hoc enim laboramus et maledicimur, quia speramus in deum uiuum, qui est saluator omnium hominum, maxime fidelium. }}\pstart Attentum uult Paulus Timotheum suum, et omnem episcopum, ut certo sibi credat, utque doctrinam uerissimam amplectatur, non cunctetur si uideat alios aduersa docentes et agentes, uerosque apostolos et uerum euangelium hostiliter persequentes factis ac uerbis. Quanquam uero ad horam periclitari et laborare uideantur propter ueritatem, quia tamen sperant in dominum tota fiducia, omnia moderantem, et omnes sibi credentes saluantem, qui deus uiuus est, non idolum, non ignarus, non inuidus, non tyrannus: ideo et nos qui fideles sumus, sine dubio, suonon destituemur auxilio, nec ueritas per nos praedicata succumbet aliquan do, sed triumphabit, et nos cum illa in gloriam dei, qui omnium saluator est, maxime fidelium. Est autem peculiare uerbum Apostolo, aliquid asseuerare uolenti, ut dicat, Fidelis sermo, ne ab aliquo uel haesitetur uel negligatur, ut inutile uel casu dictum. Vult autem intelligi euangelium quod praedicauerat Ephesiis, et edocuerat Timotheum, propter quod sine intermissione multa patiebatur incommoda. In hoc, ait, et laboramus et maledicimur. Propter hoc euangelium multa passi sumus, et adhuc patimur continuo, tanquam falsi doctores, subuersores ue ritatis, pseudapostoli et apostatae a lege paternarum traditionum, multa quoque mala de nobis loquuntur, et imprecantur, in tantum ut etiam nobis id uitio detur et crimini, quod ausi simus tales deceptores et falsi doctores et subuersores uerborum dei, secundum ipsorum iudicium, adhuc tamen sperare in deum uiuum, cuius ipsi nos blasphematores uideamur. Ipse autem dominus deus noster, cui seruimus, cui obedimus in praedicatione euangelii, in quem confidimus, cui nos totos credidimus: ipse saluator est omnium hominum. Pro omnibus quidem solicitus est, omnes suas creaturas diligit, filic suum ex amore maximo pro omni  \pend
\section*{COMMENT. IN I. EPIST. }
\marginpar{[ p.510 ]}\pstart bus tradidit, Solem suum super omnes diffundit, salutis consilia et remedia nemi ni negat, citra condignum punit, omnes uult saluari, uocauit omnes, sese omnibus tradidit, uiam salutis aperuit omnibus: maxime autem fidelium toto corde sibi credentium, in se fiduciam omnem collocantium, qui uerbis suis credunt, qui dominicis praeceptis acquiescunt, uoluntatem illius faciunt. Si quando deficiant ac peccent, statim fideliter poenitent. Non poenitentes ex animo, dolent se non poenitere quantum uellent ac deberent. Orant pro peccatis suis dominum in nomine Christi, cuius iustitia et merstis saluari confidunt, de propriis meritis et iustiriis desperant. Hii tales fideles proculdubio sunt, et domino habentur: horum certiss simus saluator deus est, qui per Christum saluat fideles, tanquam filios et intimos amicos, quos dissimulare non possit.  \pend\pstart Praecipe haec, et doce. Nemo adolescentiam tuam contemnat, sed exemplum esto fidelium in omniuerbo et conuersatione, in charitate, in fide, in castitate.  \pend\pstart Haec, inquit, praecipe, Confidentiam scilicet in deum, in omni labore et perse cutione: pietatem ueram seruandam prae omnibus, nihil praeter scripturam docendum, fabulis non attendendum, matrimonium non interdicendum, cibos non respuendos, ut animae damnosos, et c. Similiter quae de Christo docenda praescripsi breuibus. Haec, inquam, praecipe ac doce, non traditiones humanas ac doctrinas externas. Nemo adolescentiam tuam contemnat.) Aetas senectutis uita immaculata. Deus mentes magis comprobat quam aetates, non tanquam inexpertus spiritualium loquaris, non adolescentium more agas et doceas, licet adolescens sis. Factis ostende doctrina et moribus, ut adolescentia tua promereatur senectutis honores. Quodque adolescentia detrahere uidetur authoritati, hoc morum et doctri nae maturitats accedat: ideo per authoritatem officii tui praecipias et doceas, ut non humana, sed ex spiritu proueniens doctrina credatur, ut non ob adolescentiam spernatur doctrina, sed ob maturitatem et soliditatem doctrinae adolescentia ueneretur, ut a nemine sperni possit. Sed esto exemplum fidelium, et e. Hoc tandem est utiliter praecipere et docere, si prior feceris, si exemplar honestatis in uerbo et conuersatione praestiteris, si uerba factis comprobaueris, si uita non doctrinam condemnet, sed persuadeat: non enim uerba sufficient, cum coeperit Christus facere et docere. Exemplar autem sit primo in sermone, utnihil praeter scripturas sacras doceat, et solidam doctrinam, semper de honestis loquatur, obloquatur ne mini, omnia mitius interpretetur ambigua, increpet detractores, sit apud eum, Est est, Non non. Deinde in conuersatione, ut humiliter, mansuete, pudice, sobrie, et amanter ad omnes sese sine hypocrisi exhibeat, paratus omnium commodis pro mouendis, solicitus de necessitate proximi, magis quam propria corporali ac spiritua Ii. Praeterea in dilectione, charitate scilicet dei et proximi propter deum, ut nedum uerbo et opere, sed corde quoque et animo faueat synceriter, non ob commodum priuatum, etiam inimicis, externis, immeritis, ingratis, ut emendentur et proficiant. In fide) quoque addit. Vt totus ex deo pendeat, sese totum deo committat: omnia quae eueniunt, deo accepta referat, non dubitet de gratia et clementia dei. Christi beneficia credat, ac iugiter meditetur. Dei gloriam quaerat in omnibus, uerbis dei assistat, credat, discat, doceat. Defendat ueritatem fidei catholicae et sa crae scripturae, non ferat infideles. Tandem castitatem uel puritatem addit. Fides enim Christiana cor mundum creat, uerba spurca non patiturl, opera turpia ex animo exhorrescit, nihil mixtum, uarium, impurum, duplex et insyncerumad.  \pend
\phantomsection
\addcontentsline{toc}{subsection}{\textit{}}
\subsection*{\textit{}}
\section*{AD TIMOTH. CAP. IIII. }
\marginpar{[ p.511 ]}\pstart mittit; omnia autem sint simplicia, dextra, certa et pura quae agit et loquitur, Castitatis uirtute libido regatur, et minime regat corpus uel animum.  \pend
\phantomsection
\addcontentsline{toc}{subsection}{\textit{Dum uenio attende lectioni, exhortationi et doctrinae, }}
\subsection*{\textit{Dum uenio attende lectioni, exhortationi et doctrinae, }}\pstart Non otiosus semper expectes aduentum meum, noli mihi perficienda absenti et uenturo reseruare, quae tu praestare poteris: interim lectioni instes scripturarum sacrarum, unde possis et subditis proficere. Non manibus uerses nisi sacros codices, non mundanarum uanitatum dogmata, non philosophos doceas, sed leagem et prophetas ac Psalmos, id est canonem scripturae ueteris. Euangelia enim eo tempore apostolorum praedicationibus discebantur, donec describerentur et euulgarentur per discipulos. Haec, inquam, tunc legebantur ab apostolorum disci pulis, non fabulae gentilium uel ludaeorum. Non autem simplicem lectionem indicit, sed ad intellectum, sensum, profectum et fidem, ne tempus perdatur. Non ut hactenus horae canonicae lectae sunt pro more et praecepto, sed ut primi illi pa tres ad minus in anno, totas sibi scripturas praelegebant. Quibus successores commentaria adscripserunt et annotationes: quae si extarent, superflua haberentur quae hodie conscribuntur. Vnde et Romanae ecclesiae consuetudo inoleuit, ut clerici et episcopi certis temporibus libros legere debeant, ueteris et noui testamen ti, iuxta dispositionem Gelasii Papae, ut habetur in Decretis, distinct.15. Canone, qui incipit, Sancta Romana ecclesia. Idque nominant Ecclesiasticum officium, sic per anni ciculum digestum, ut Pentateuchum legant a Septuagesima, per septem seilicet hebdomadas septem libros, quinque Mosis, cum losue et Iudicum. Deinde per quindenam Ieremiam usque ad Pascha, cuius diebus Homeliae sanctorum leguntur, et euangelia festi. Deinde usque ad Pentecosten Apocalypsim, Acta cano nicaeque epistolae. Deinde usque ad calendas Augusti, libros Regum et Paralipome non. Per Augustum Salomonis libros. In Septembri Iob, Tobiam, Hester et Esdram legi ordinantur. In Octobri autem Machabaeorum libri. In Nouembri, Ezechiel, Daniel, et minores prophetae legendi. Tandem in Decembri, Isaias. Post na talem domini, epistolae Paulinae omnes usque ad Septuagesimam. Haec habet Romanum ordinarium, et prisca ecclesiastica traditio: quae utinam seruata fuisset, intermissis Legendis dubiis sanctorum, et id non transisset in consuetudinem tan tum legendi, non etiam intelligendi sacrosanctas scripturas, melius certe interim ecclesia profecisset. Quod utinam uel posthac felicius agendum tentetur a minia stris ecclesiae, ut praemissa oratione ex Psalmis uel alia, intendant lectioni sacrae scripturae, canonicas scripturas horis canonicis ruminando, magis quam uocaliter mur murando et cum taedio. Nec tamen interim oportet obmittere etiam salsamentoa rum uice, saecularium authorum utiles libros, et liberalium disciplinarum studia. Semper tamen et ubique curato, ne quid nimis, et ut sapiatur ad sobrietatem. Nec ualet exemplum Christi, qui orasse frequentius legitur q̃ legisse. Discere enim Chri stus sacra non potuit, sapientia dei patris. Nos scripturas scrutari iubemur, et incumbere lectioni. Paulus eruditissimum Timotheum admonet, ac praecipit seriose, ut sana doctrina imbuatur, quam et doceat in ecclesia, ut et contradicentes re uincere possit, et redarguere. Exhortationi quoque iubetur attendere:non enim le ctio sufficit, sed exhortanda est ecclesia ad fidem in Christum et apostolicos mores, non cogendo sed exhortando, ut imperfecti instituantur. Vitia, infirmitates, et imperfectiones carnis dulcibus uerbis disuade, si locum inueniant: non terribiliter ac imperiose manda, non mineris, non dure increpes, sed exhorteris suauiter, sed instanter et sine intermissione, non leui capta occasione remittendo studium, uel desperando de deigratia. Doctrinae quoque instare iubet Paulus, ut ex  \pend
\section*{COMMENT. IN I. EPIST. }
\marginpar{[ p.512 ]}\pstart serinturis diuinis ostendat, euincat, doceat, sic se habere omnia iuxta del legem et sacras literas, ut in ecclesia praedicatur a ueris apostolorum discipulis, ne quae dicantur sint incerta et opiniones, sed scripturarum fide et authoritate solidata.  \pend
\phantomsection
\addcontentsline{toc}{subsection}{\textit{Noli negligere gratiam quae in te est, quae data est tibi per prophetiam, cum impositione manuum presbyterii. }}
\subsection*{\textit{Noli negligere gratiam quae in te est, quae data est tibi per prophetiam, cum impositione manuum presbyterii. }}\pstart Donum prophetiae, hoc est gratia interpretandi sacras literas, affluentius ei datum erat q̃ caeteris, ideo episcopus ordinatus est, cuius officium est docere legem dei. Nihilominus tamen lectioni instare iubetur, ne obmittatur locorum collatio ut certior interpretatio haberetur. Non enim quaelibet passim concinnata enarratio sufficit, sed quae desumitur ex scripturis, nisi manifesta sit lectio. Hanc in te prophetiae gratiam uide ne negligas, ne talentum creditum abscondas, sed usu et exercitio quotidiano usuras referas domino tuo debitas et exigendas. Cum tibi presbyteri et seniores imponerent manus, ut esses episcopus, prophetiam quoque  habebant, te futurum idoneum huic officio, quia donum spiritus et gratiam suscepisses docendi ecclesiam. Eo igitur prophetiae dono utere solicite, et ne negligas, sed doce, exhortare et insiste pro uiribus, iuxta acceptam gratiam, per manuum impositionem in profectum ecclesiae tuae.  \pend
\phantomsection
\addcontentsline{toc}{subsection}{\textit{Haec meditare, in his esto, ut profectus tuus manifestus sit omnibus: attende enim tibi et doctrinae, insta in illis. Hoc enim faciens, et teipsum saluum facies, et eos qui te audiunt. }}
\subsection*{\textit{Haec meditare, in his esto, ut profectus tuus manifestus sit omnibus: attende enim tibi et doctrinae, insta in illis. Hoc enim faciens, et teipsum saluum facies, et eos qui te audiunt. }}\pstart Admonui te legis diuinae praedicandae, et Christianae fidei articulorum, instare lectioni, exhortationi, ad fidem, spiritum, dilectionem, puritatem, conuersatio nem et sermonem. Haec, inquam, meditare, ut apposite, sufficienter et fructuose loquaris de eis, non impraemeditate et temere, deque fabulis. Nondum commendaui tibi quicquam de Missis, consecrationibus, confessionibus, exorcismis, defunctorum animabus, et officlo ceremoniarum. Quae, quia non accepi a domino, ideo eorum nihil tradidi tibl, quamlibet insignis apostolicae ecclesiae episcopo, fidelis apostolus. In his itaque quae tibi proposui totus esto, studium omne confer, et cura solicite ut fiant: et ita fiant, ut per haec proficias patenter in subditis, et in ec clesia tua tanta, ut prosit ouibus pastoralis cura tua, ut sciant ex spirituali commo do suo se habere episcopum, doctorem et pastorem idoneum. Vnde et proficiant semper in melius, et crescat indies fides et conuersatio Christiana. Attende quoque tibi et doctrinae.) Ne culpa doctorem redarguat, ne actio orationi repugnet, sed sana et irrefragabilis sit doctrina tua. Contra et praeter quam nulla introducatur in ecclesiam, sed ea ut certa ac firma statuatur et propugnetur. Insta in illis docendo et agendo, ut proficiat populus. Illa, inquam, non alia propone, inculca, defende. Quae si pro uocatione et gratia dei feceris, officioque tuo incubueris, teipsum saluum facies, euades iudicium nequam et infidelis serui, et negligentium pastorum, sua quoque tantum quaerentium, gregem deserentium uel perdentium. Non audies comminationem et grauem damnationem pastorum sese pascentium, eorum te periculo eximes, etiamsi non acquieuerint, contemnant autem et reiiciant, non credant, et damnentur. Tu tamen liberabis animam et conscientiam tuam, principi pastorum gratam referes rationem. Si uero audierint et fuerint obsequuti, si monitis iustitiae et sanae doctrinae tuae dociles animum apposuerint, ut ea perficiant, tu quoque illis causa salutis eris, et una tecum saluentur, et salutis suae fidum te fuisse ministrum testabuntur fideles.  \pend
\section*{AD TIMOTH. CAP. V. }
\marginpar{[ p.513 ]}
\endnumbering\beginnumbering\section{CAPVT V.}
\phantomsection
\addcontentsline{toc}{subsection}{\textit{\huge\textbf{S}\normalsize Eniorem ne increpaueris, sed obsecraut patrem, iuuenes ut fratres, anus ut matres, iuuenculas ut sorores, in omni castitate. }}
\subsection*{\textit{\huge\textbf{S}\normalsize Eniorem ne increpaueris, sed obsecraut patrem, iuuenes ut fratres, anus ut matres, iuuenculas ut sorores, in omni castitate. }}\pstart Modestia iuncta authoritati, necessaria est episcopo, neminem non obseruet, nullum negligendo praetereat, omnes curet, proficiat omnibus, et medeatur. Etiamsi grandioris aetatis sint, si quid contra ea quae docuerit episcopus, ab illis fieri et committi uideat, hoc est, contra legem dei comiter, amice, pia admonitione, sine supercilio et fastu, et sine seueritate corrigat. Non facile seniores ferunt a iuniore moneri, multo minus increpari seuerius, uel saeuius obiurga ri. Et tamen insolentia senum neutiquam dissimulanda est, ne ob senium et annos ipsorum delicta et mores in authoritatem et exemplum trahantur a iunioribus, procliuioribus ad mala sectanda. Corrigi ergo senes debent, sed non saeuitia, non uerborum duritia, sed blanditiis, lenitate grata et matura, ac quasi reuerentiali insinuatione admonendo, non contumeliis irrogandis, quae ferri non possunt a senibus, quin magis exasperantur. Seruanda igitur reuerentia est, et ratio aetatis, sed tamen non intermittenda admonitio transgressionis, uerum non uelut imperiose, sed cum modestia. Si uero senes omnino stulti sint, et inueterati dierum ma lorum, et publico scandalo patrocinantes, senectutis suae authoritate iuniorum insolentiis exemplo uel uerbo, iam Danielis exemplo utendum est, et pro rei indignitate rigide tractandi. Semper tamen decet iuuenem praelatum modestia in uerbis, de quibus hic agitur, ad senes, quatenus magis condolentia et dolor mali tiae et rigoris necessitas urgere uideatur episcopum, quam contemptus senectutis, uel praesumptio praelationis, quam iuuenes declinent. Iuuenes utfratres.) Solemus uariis modis alloqui et corrigere fratres insolentes, iuxta cuiusque indolem, ut unum durius alloquamur, alterum familiarius affectuque suauiori: aliquando et de industria irritamus, etiam intendentes turbari aliquem, sed ut citius resipiscat. Nunquam tamen infamia publica afficiendi sunt iuuenes, nec contemptu uili afficien di coram multis: praeuenienda est iactura famae et honoris, dum salua omnia ada huc sunt, et honestatis spes in adolescentibus, ne semel confusi, animum despondeant, et desperatione reddantur peiores. Ideo contra externos accusatores studiose potius excusandi sunt iuuenes, defendendi et uindicandi: uerum familiariter et amice conueniendi sunt cum digna duritia, ut intelligant fidem et amicitiam: si publice propugnentur, priuatim obiurgentur cum modestia. Sic fratrum animi conciliandi sunt, sic emendandi iuuenes: sic proficitur, si omnia ex syncera naturalique fraternitate tententur et aptentur: caeteris tamen paribus, ut non sint procaces et filii Belial. De Christianis em̃ fit sermo et coaetaneis, ad quos pro cha ritate afficiamur, ut consulere cupiamus, dum scilicet agnouerint se amoris tantum causa admoneri et corripi, non odio. Anus ut matres.) Honorifice alloquendo, commonendo, dedecus ostendendo, quod illas maneat, erubescentiam incutiendo, ex proposita rei indignitate, non contemptu, odio, maleuolentia, sed uerbis modestis, ne magis irritentur quam emendentur, et ad desperationem adigan tur, et tristitia obruantur, bis contemptae, ex uitio scilicet et aetate uel fama. Iuuenculas etc.) Amanter ac syncere sorores alloqui consueuimus, ac sinae formae uellibidinis quacunque habita ratione, tantum ut consulamus honori et rebus agen dis. Sic dum episcopus uel parochus necessitate adigitur ad alloquendum iuuen culas, breuiter id faciat, syncere, grauiter et caute, omnia obmittens leuitatis signa, uerecunde, composite, et ordinate ad disciplinam, ne quo pacto sinistrae su\pend
\section*{COMMENT. IN I. EPIST. }
\marginpar{[ p.514 ]}\pstart spicionis indicium proferatur, et castitatis fama periclitetur, etiam si nihil mali in tercedat. Abomni enim mala specie abstinendum, cum et lubrica sit fama pudoris, maxime episcopo inter iuuenculas et adolescenti, quantumlibet grauis sit moribus: nemo enim satis et omnibus horis sapit, hac potissimum in causa.  \pend
\phantomsection
\addcontentsline{toc}{subsection}{\textit{Viduas honora, quae uerae uiduae sunt. Si qua autem uidua filios habet aut nepotes, discat primum domum suam regere, et mutuam uicem reddere parentibus, hocenim acceptum est coram deo. }}
\subsection*{\textit{Viduas honora, quae uerae uiduae sunt. Si qua autem uidua filios habet aut nepotes, discat primum domum suam regere, et mutuam uicem reddere parentibus, hocenim acceptum est coram deo. }}\pstart Ex consequentibus contextuque literae uel sermonis satis intelligitur, is honor de quo hic dicitur, deferendus uere uiduis, hoc estomni auxilio uiduitatis et sola tio, quae scilicet marito, filiis, nepotibus, et diuitiis destitutae, despiciuntur, et misere animam trahere permittuntur. Non patiatur episcopus huiusmodi male audire, indigere, odiri, maleque tractari, sed eam, si uere sit uidua, ecclesiae nomine defendat, suuet, consoletur, de sumptibus ecclesiae, et collectis in synaxi fidelium, prouideat talibus uiduis in communi de necessariis. Idem autem intelligendum de omnibus ecclesiasticis personis, eadem miseria laborantibus, pupillis, orphanis, infirmis, peregrinis, et qui sua pauperibus erogarunt, ut postea necessariis egere cogantur. In summa, curet summopere episcopus, ut nemo inter Christianos sibi commissos indigeat, uel mendicare cogatur. Si qua autem etc.) Si aliqua non sit uere uidua, sed marito tantum uiduata, non autem filiis, nepotibus et diuitiis, istam non tuebitur episcopus sua solicitudine, sed qui id possunt filii et nepotes, uictum illi administrantes et amictum: ipsa uicissim eos moderetur, instruat, mo neat, et, si potest, nutriat, foueat ac doceat: perinde ac se quoque simili humanitatis officio nouit a parentibus educatam. Debent enim patres thesaurizare filiis, non contra. Nec huiusmodi uidua suscipiatur ad collegium uerarum uiduarum, sed domui et familiae suae potius praesit et incumbat. Hoc enim domino placet, ut prouectiores aetate et prudentia praesint iunioribus: et qui aliis benefacere possunt, non patiantur otiosi sibi benefieri cum iactura egentium, et officium exigere, qui debent praestare et possunt. Discant primum etc.) Non se talis uidua intromit, tat ministerio ecclesiae, dimissa familia fillorum nepotumque suorum, istis suis potius praesit et inuigiler, benefaciat, et regat propriam domum: quam ubi bene rexerit, regendoque exercitatior euaserit, postea ecclesiae deseruire legitime poterit, si ad hoc per episcopum sexagenaria deligatur.  \pend
\phantomsection
\addcontentsline{toc}{subsection}{\textit{Quae autem uere uidua est et desolata, speret in deum, et instet obfecrationibus et orationibus noctu die: nam quae in deliciis est, uiuens mortua est. }}
\subsection*{\textit{Quae autem uere uidua est et desolata, speret in deum, et instet obfecrationibus et orationibus noctu die: nam quae in deliciis est, uiuens mortua est. }}\pstart Describitur hic uere uidua ea dici, quae solatio huius mundi caret, marito, filiis, nepotibus, possessione: sperat autem in domino, et omnem spei suae fiduciam collocat in deum, qui pupillum et uiduam suscipiet. Non utitur hoc mundo, non sectatur gaudia huius mundi, et carnis libidines, sed uacat orationibus et audientiae uerbi dei, seruitio indigentium, consolatione afflictorum, sine intermissione cogitat quae domini sunt, exonerata molestiis huius uitae, et domesticae solicitudinis, quomodoque placeat marito, sed deo uacat soli. Talis estuere uidua, mun do neglecta, deo autem coniunctissima et commendatissima. Die noctuque ́, ait. Si cut de Anna legimus filia Phanuelis, quae Christum suscepit oblatum in templo. Vidua uero quae in deliciis uersatur, non seruiens deo, nec uacans orationibus, sed libertate abutens, delicias sectatur, otio uacat et multiloquio, uinolentiae, tra\pend
\section*{AD TIMOTH. CAP. V. }
\marginpar{[ p.515 ]}\pstart dita libertate in occasionem carnis. Talis, ut maxime suauiter iuxta carnem uiuat. mortua tamen est spiritu, gratia dei priuata, Christo non uiuit, quia non pietats studet. Tales ab ecclesia nutriri minime debent, cum sint opprobrium et diffamia ecclesiae, sed corrigi potius ac puniri oportent, si quo modo reuiuiscant in Christo per poenitentiam.  \pend
\phantomsection
\addcontentsline{toc}{subsection}{\textit{Et hoc praecipe, ut irreprehensibiles sint. }}
\subsection*{\textit{Et hoc praecipe, ut irreprehensibiles sint. }}\pstart Officii quoque tui est, circa uiduas curandas, ut eas doceas maturitatem, et cir cumspectionem tantam, ut nihil sinistri de eis suspicari uel dici possit, ne laxius uiuant ob libertatem suam in uiduitate, ne scandalo sint ecclesiae: uideanturque non ob religionem et castitatis amorem abstinere a secundis nuptiis, sed ad fatisfaciendum libidini exoneratae maritis.  \pend
\phantomsection
\addcontentsline{toc}{subsection}{\textit{Si quis autem suorum, et maxime domesticorum, curam non habet, fidem negauit, et est infideli deterior. }}
\subsection*{\textit{Si quis autem suorum, et maxime domesticorum, curam non habet, fidem negauit, et est infideli deterior. }}\pstart Quanquam haec sententia uel regula de uidua dici uideatur quae non est uere uidua, quae curare iubetur familiam, potius quam ut uel ecclesiae commendetur alenda, uel eidem ministratura, sed potius docetur praeesse suae domui. De quibus uiduis re uera Paulus haec quoque retulit. Veruntamen generalis gnome uidetur, pronunciata ab Apostolo uniuersaliter de uiris et mulieribus, uiduis et maritis. Volens neminem sic curare aliena, ut propria familiae suae necessaria et domesti corum suorum, quibus dei prouidentia rector ac prouisor constituitur, curam negligat uel reiiciat, ut ea tantum ratione orationi uacare eligat, ne cogatur proximorum necessitati inseruire. Quin hac sententia Pauli iubetur paterfamilias quicunque et materfamilias, domesticorum administrationem diligenter gerere, non secus quam si a domino specialiter eorum curam demandatam suscepisset. Quod si abiiciat instituendae prolis curam, regendae domus, sustinendi molestias proximorum, de quibus deus mandauit cuique : eligatque obmissis istis, uacare propriae deuotioni, ecclesiae ministrare, religiosus ad speciem fieri, deoque secundum proprium suum sensum seruire, hic tantum abest quo deo placeat, ut magis infideli deterior definiatur a spiritu dei. Qui uero legitime curat suam familiam, ut Abraham, Isaac, Iacob, Ioseph, et caeteri, id sane agit quod natura iussit, et dominus praecepit. Qui uero pro arbitrio suo agenda eligit, et quae deus mandauit obmittit, dominum men dacem et mutabilem facit, aliudque uelle, aliudque iubere credit: id quod fidei aduer satur. Maxime autem quod infideles quoque naturae ductu optimo et minime improbo, proli prospiciunt et proximis. Idipsum deus praecepit, natura bestiis quanlibet ferocibus indidit, et ratio naturalis aequissimum iudicat. Illud hii suo sensu abum dantes despiciunt, propterque traditiones proprias et humanas statuendas domini dei mandatum transgrediuntur. De his ex Plinio et Petrarcha aliisque ́, ut ubique , et ex scripturis sacris multa adduci possent, nisi clarior esset haec Pauli sententia Sole, et uerior statutis humanis.  \pend
\phantomsection
\addcontentsline{toc}{subsection}{\textit{Vidua eligatur non minus sexaginta annorum, quae fuerit unius uiri uxor, in bonis operibus testimonium habens: si filios educauit, si hospitio recepit, si sanctorum pedes lauit, si tribulationem patientibus subministrauit, si omne opus bonum subsecuta est. }}
\subsection*{\textit{Vidua eligatur non minus sexaginta annorum, quae fuerit unius uiri uxor, in bonis operibus testimonium habens: si filios educauit, si hospitio recepit, si sanctorum pedes lauit, si tribulationem patientibus subministrauit, si omne opus bonum subsecuta est. }}\pstart Ecclesiae apostolicae primitiuae ecclesiastici, tunc et deinde dicebantur non tantum clerici, sed et omnes pauperes, orphani, pupilli, et in primis uiduae, hoc est, orbatae uel desolatae uiris filiis et possessione, nonnunquam ob fidei confessionem  \pend
\section*{COMMENT. IN I. EPIST. }
\marginpar{[ p.516 ]}\pstart spoliatae. Bona autem ecclesiae erant, quae sabbathis uel dominicis diebus, post au ditam fidei doctrinam colligebantur, nemine uacuo comparente. Collecta autem siue ad esum pertinerent, siue pecuniae, benedicebantur ab episcopo, et per diacones, id est ministros ecclesiae, ad episcopi fidelis dispositionem ecclesiasticis talibus diuidebantur, prout cuique opus erat: de quibus quoque episcopus et diaconi cum sua familia, uxore et filiis sustentabantur ad uictum et amictum, quibus et contenti erant: non ad abundantiam uel luxum. Peregrini etiam fratres, boni testimonii, reficiebantur, uirgines pauperes, et infirmi. Nemo egere uel mendicare permittebatur. Temporis autem successu testamentis fidelium, et munifica diuitum largitione huic usui prouisum est, et indigentium consolationi. Alicubi etiam ex decimis et possessionibus agrorum ac domuum fortassis. Tandem deuotio prin cipum, et simplex religionis amor, Christianaque religio regum, cum non possent ipsi prospicere pauperum necessitati, occupati rebus suis et officiis, huiusmodi onus et studium inseruiendi et distribuendi necessaria egentibus commendauit religiosioribus personis, et probatae sanctimoniae uiris, abbatibus et monachis, qui ut tunc praeter arctum uictum et amictum, in summa parsimonia uictitabant aliquandiu: nec suspicio aliqua fraudis locum haberet, ne non dispensarent illi serui dei necessaria indigentibus illis sibi commissis. Eis quoque personis ecclesiasticis sacris literis incumbentibus, inque spiritualibus populo seruientibus, orphanis quoque et uiduis, hospitibus et infirmis. Vnde et usque hodie in uetustissimis monasteriis et collegiis bene ordinatis, ordinationes de lectionibus quotidia nis sacrae scripturae, infirmorum curandorum officiales, hospitum susceptores, scholastici, qui pueris praesunt instituendis, sacerdotes sacramentorum dispensa tores. Episcoporum uero uel abbatum proprium munus erat proponere uerbum dei ad ecclesiam. Semper tamen durauit festis diebus oblationum collecta proec clesiasticis. Temporis autem lapsu, proque morum mutatione, deprauatis omibus, charitate frigescente, quaestu autem luxuque irrepente paulatim in ecclesiasticos, bona ecclesiastica et diuitiae in usum pietatis datae cesserunt in maximum abusum, continuoque magis inualust in dedecus cleri et monachorum. Episcopi enim principes facti sunt, imo ex principibus facti sunt episcopi. Monachi et abbates satrapas se gerere coeperunt, suisque priuatim usibus applicare bona ecclesiastica, et pauperes negligere, quos in sua mensa habere iussi erant, uel aliis rationibus aliquantisper prouidere, extructis hospitalibus. Religiosi caeteri plus q̃ mundani facti sunt, quos suis coloribus et dignis depinxit ante quadringentos annos D. Bernhardus, et ipse monachus, qui dum reformationem monachorum meditatur et incipit, longe ditiores, et ceremoniis addictiores, orbique Christiano inutiliores post se reliquit, et diuitiarum corrasores, sub pietatis specie auariores, quam ante eos fuissent Benedictinenses, quos Bernhardus obotium et diuitias digne traduxerat in Epistola uel Apologia sua ad Gulielmum abbatem, digna quae legatur a piis et impiis monachis, speculum sane nouissimi monachatus. Omnino autem qui spirituales nunc dicuntur et religiosi, toti facti sunt animalia uentris, praeter cantum et murmur suum, otio tantum incumbunt: quod quas sequelas relinquat, scriptu. ra et tota philosophia docet, et rerum usus cum experientia. Hodie non sexagenariae uiduae eliguntur uel uirgines domino educandae, ad ecclesiae obsequium et honorem, sed filias quas odiunt, et per quas posteros non dignantur, quas sub pal lio religionis malunt meretrices in irreformatis coenobiis, quales maximo passim numero habentur, uel anxie ac turpiter in conscientiis urendas et desideriis, quam matresfamilias fieri. Filios quoque suos in monasteriis nebulones et lenones agere  \pend
\section*{AD TIMOTH. CAP. V. }
\marginpar{[ p.517 ]}\pstart imuni, et nuuriunt ôblationibus adulteros, ubi et impinguantur et luxuriantur. de bonis ecclesiae. Non est hic locus excusationis multorum pessimorum hominum, obpaucos bonos, sub arcta disciplina ad honestatem coactos. Siquidem cancellis tantum regularibus, et timoribus poenarum continent sese apud homines, quicquid secreto in corpore et conscientia, naturali necessitate et usu eueniat, cum dominus pectus, conscientiam, animum et cor exigat pacatum et pium, totaque fiducia sibi sana conscientia inhaerentem. Non quidem sine omnibus peccatis, sed sine diffidentia emendandi errata, quae uotis tantum impediri non possunt, nec quacunque abstinentia et rigorosa corporis castigatione. Quae naturae peccatricis propensio nititur tantum in uetita, et a uoluntate hominum quamlibet san cta non pendet, gratiae dei promissionem non habet ex uerbo dei. De quibus mul ta dicipossent, iamque dici inceperunt a multis bonis uiris, qui ecclesiasticum statum cuperent non abolitum, sed uerbo dei reparandum in melius, quod huiusmodi perpetua uota ignorat et improbat. Interim dum uentris haec animalia nutriuntur in collegiis quibuscunque , pauperes caeteri pupilli et uiduae negliguntur a clericis: quicquid corradere possunt, reponunt in census, unde augeatur numerus hoc pacto deum uentrem colentium, cantorum et otiosorum hominum, ut in immensum uideamus succreuisse multitudinem, solo nomine spiritualium, re uera autem carnalissimorum scortatorum cum suis spuriis et meretricibus. Omnia sub praetextu nominis cultusque diuini. Susceptores autem hospitum, infirmo rum saeruitores, plagarum horribilium et grauium infirmitatum curatores, pua pillorum et uiduarum patroni, cum re uera ecclesiastici sint, hodie saeculares uocantur, nisi ubi pro diuitiarum collecta, religiosorum domus nominantur Hospi talia, ut nihil sit integrum. Heae autem uiduae, de quibus hic Apostolus, non omnes assumebantur ad stipem uel annonam ecclesiae, uel etiam ad tutelam, sed eligebantur ad praeceptum apostolorum et apostolicorum uirorum: tales autem qua les hic describuntur. Non minoris aetatis quam sexagenariae, ne semel susceptae in cu ram ecclesiae, redeuntes ex carnalitate ad coniugium: uel per aetatis lubricitatem peccantes, probro magno afficerent ecclesiam, neue impudicitiae suspicio locum haberet, quia ad istam aetatem prouectae, et uiribus destitutae manuum labore se nutrire amplius non possent: essent autem uerae uiduae, hoc est solatio et patrimo nio destitutae, integrae autem famae. Additur autem, Quae fuerit unius uiri uxor, de qua scilicet non sit suspicio intemperantiae per aetatem et pristinam conuersationem, qualis autem notaretur, si saepius nupsisset, deque maritorum possessionibus ditior euadere potuisset. Laudi enim dabatur non solum in ecclesia prima, sed et apud gentiles, si nuptiis secundis casta abstinuisset mulier, et non parua nota censebatur uel dedecus denuo nupsisse, uel saepius. Hoc quidem quanquam apo stolus non omnino probet, placuit tamen ut continenter uiuant quae possunt, et non iterent matrimonium, si sint prouectioris aetatis, caeteris paribus. Mandat au tem et uult uiduas iuniores nubere, quod tutius sit, et continentiae necessarium remedium, quae nuspiam purior, sanctior et tutior inuenitur et facilior, quam in sacro matrimonio, quod et suam continet claritatem casta generatio, exemplo sanctissimorum hominum. Interim ingens cauetur periculum et scandalum exstiale et intolerabile. Quod autem hic ab Apostolo ueluti peccatum notari uidetur bigamia uel secundae nuptiae, siquidem secundo nuptam indignam iudicat stipe eccle siastica, hoc agit, consulere uolens honestati, adactus spiritu apostolico, et fidelium famae apud tunc temporis infideles, nolens id honoris deferri foeminis, quae ob im pudicitiae notam a gentilibus reprobarentur, cum hic tantum ageretur de subsidio  \pend
\section*{COMMENT. IN I. EPIST. }
\marginpar{[ p.518 ]}\pstart corporali, quod talibus poterat abalns, et modis ains compeniari. Interimuero consulebatur paci et tranquillitati reliquarum uiduarum ecclesiasticarum, quae non sine murmure huiusmodi admisissent suo contubernio, uel sibi aequari sustinuissent. Nihil tamen syn ceritati conscientiae adimebatur, quam ut liberam seruaret, uel inuitis uel damnantibus gentilibus, qui secundum mundum sapiebant, secundas nuptias improbantes in mulieribus, non in uiris: ipse fidelis existens Aposto lus non solum permittit, sed praecepit secundas nuptias iunioribus uiduis, ne deterius contra deum peccetur, quod hoc indulto praeueniri possit. Arcet autem hu iusmodi a stipe ecclesiastica, consulens famae, paci et castimoniae: non ut apud deum ream, sed infirmorum infamationi obnoxiam apud gentes. Sicut et sterilitas uxo ris apud Iudaeos probro non carens, eadem ratione arcetur, quae tamen noxa carere constat apud deum et homines, sed non apud Iudaeos, quibus sicut benedictio erat filios generasse, et promissionis donum: sic caruisse maledictio putabatur, et ut infausta et misera plangebatur. Sterilis autem aliquandiu mulier, si peperisset tardius, ingens gaudium protulisset caeteris uicinis mulieribus. Dicitur consequenter. In bonis operibus testimonium habens.) Exigebatur charitatis et misericordiae exercitium omnibus notum, ut eius honor et fama religiosa nota fieret toti ecclesiae, a qua sustentanda erat uidua, ne forte fauore et falsa laude promotam dicerent, sed conspicua essent multis testibus, pietatis suae digna opera et studia sancta. Si filios educauit.) Non dicitur, Si filios peperit, licet id uide ri possit significasse, liberos scilicet generasse et educasse, fortassis fuerat supplen dum religiose, honeste et laudabiliter: uel ut sterilis arceatur ab ecclesiastica stipe in gratiam Iudaeorum. Quanquam probabilius uideatur de non sterilibus inten disse, quandoquidem quae sterilis tanta aetate perstitisset, molestiis liberorum carens, sibi prouidere laborando poterat, de sustentamento senectutis suae, Similiter si generans filios statim amisisset nondum educatos. Quin et si improborum filio rum uelfiliarum mater, potuit detestabilior uideri contubernio sanctarum uiduarum. Negari denique non potest, infelicifsimum apud Iudaeos habitum, si uxor sterilis cuiquam contigisset, quae et ut infausta, admittenda non uideretur ad stipem.  \pend\pstart Si hospitio recepit.) Hospitalitas in ueteri et noua lege impense laudatur, et apo stolorum tempore suscepisse praedicatores euangelii, et pro ueritate extorres patria, egentes, uagos et peregrinos, magnum pietatis officium non indigne habebatur: nec fidelis aestimari poterat, qui huiusmodi aliquando exclusisset hospitio, uel indigne tractasset. Nec digna erat ecclesiastico nutrimento, qui euangelistas Christi et ec clesiae doctores ac plantatores opulenta domo repulisset, uel negasset hospitium.  \pend\pstart Si sanctorum pedes lauit.) Sine calceamentis ambulantium primum et facile beneficium est, lauari pedes. Hoc Abraham exemplo docuit et Lot, ut fatigatt suauius quiescant. Christus quoq in ultima coena idem officium discipulis commendauit, charitatis monimentum, et praecepit mutuum exhibere ob humilitatis exercitium, quae charitatis fomentum esse consueuit. Sanctorum dicit, pro praedicatione euangelii discurrentium, et fidelium in causa fidei et salutis laborantium, uel sustinentium exilia et expulsiones pro uerbo dei, et uera fide praedican da in mundo. Subdit quoque : Si tribulationem patientibus subministrauit.) Si enim ex doctrina euangelica omnibus petentibus tribuendum est, si praedicatoribus nihil possidendum, si pro euangelio asserendo omnia dimittenda sunt, si mul ta patienda, si de ciuitate in ciuitatem fugiendum, si non claudenda uiscera abege no, si in opera misericordiae profundenda sunt omnia, si nihil proprium habendum quod non sit paratus impartiri egenti Christianus, certe ecclesiaslica disciplina  \pend
\section*{AD TIMOTH. CAP. V. }
\marginpar{[ p.519 ]}\pstart Ipiritu sancto ab apostolis edocta est, ut huiusmodi pauperes uel sic depauperati necessariis non fraudentur, sed Christi praeceptis obedientes, Christi quoque prouidentia necessariis non indigeant. Non enim fieri potest, si iuxta euangelium ui uatur, ut desit cuiquam uitae necessarius prouentus et suppetentia, quibus centuplum etiam in hac uita promittitur. Dicitur autem tribulation em patientibus, non oblectamenti, curiositatis, instabilitatis, fastidiosi otii, auaritiae, superbiae uel consolationis gratia discurrentibus subministranda hospitalitas huiusmodi, sed pro ueritate et euangelio patientibus gratia impertienda laudatur et muneratur. Fuerunt enim tunc quoque ut nunc pseudapostoli, qui domos uiduarum deglube bant, seducebant uariis doctrinis, ueros apostolos persequebantur, commoda, glo riam et solatia propria sectabantur. Non uult Paulus talibus, sed tribulationem pro iustitia patientibus ministranda esse uitae necessaria. Si omne opus bonum subsecuta est.) Si studiosa fuit aemulatrix operum bonorum, quae per eam fieri po tuerunt: si dedita officiis pietatis, si assidua in omni bono opere. Non sufficit quae dam egisse bona, nullum potius obmisisse laudabile est, ubi adfuisset facultas. Hu iuscemodi uiduas pauperes ecclesia nutrire debet: quin et uiros his dignos laudi bus recitatis, de bonis ecclesiae debetur diuino iure sustentatio. Non sunt enim bo na ecclesiastica templorum et aedificiorum, pinguiumq ac otiosorum clericorum, sed ministrorum ecelesiae, et honestorum pauperum necessitatem habentium. Illorum maxime, qui profuderunt sua propter euangelium in usus pauperum, quique  toti propensi in opera pietatis, neglectui habuerunt temporalia bona sua.  \pend
\phantomsection
\addcontentsline{toc}{subsection}{\textit{Adolescentiores autem uiduas deuita. }}
\subsection*{\textit{Adolescentiores autem uiduas deuita. }}\pstart Viduas intelligit omnes iuniores foeminas, extra matrimonium uiuentes, uel absentibus maritis habitantes solitarie, ut sunt etiam sanctimoniales, et de quibuscunque scandalum uel uiolenta suspicio nasci posset. Non enim tutum est, nec suspicione uacat, si episcopus frequentet uiduas iuniores, a qua suspicione omni bus modis antistiti et doctori abstinendum. Nonhic Apostolus meminit negorii confessionis, ut uel excipiat, uel in manifesto fieri, et breuiter expediri cum multa maturitate ea colloquia ordinet: proque conscientiarum scrupulis secretis explicandis leges non decernit cum mulierculis agendi. Quod ut maxime necessarium, a monachis et plebanis diligentissime praedicatum est, quod tamen uidemus in omnibus apostolorum epistolis et ordinationibus neglectum, hauddubie ut minime necessarium, alioqui si tam necessaria ad salutem fuisset haec omnium peccatorum confessio auricularis, ex qua nostris temporibus et antea tot personae periclitatae sunt, in suis conscientiis audiendo et confitendo occultissimas tentationes et crimina, certe apostoli tale aliquid scripsissent. Et quod infelicissimus Antoninus in suo confessionali docuit, id dudum sancti episcopi et doctores, Cyprianus, Augustinus, Ambrosius, Chrysostomus, et Hilarius minime obmisissent, qui nec unicum tractatum ea de re scripsisse inueniuntur. De qua Summistarum authores mendicantes illi fratres maxima uolumina aediderunt, et scandala conscientiis inextricabilia posuerunt animabus piis.  \pend
\phantomsection
\addcontentsline{toc}{subsection}{\textit{Cum enim luxuriatae fuerint in Christo, nubere uolunt habentes damnationem, quia primam fidem irritam fecerunt. }}
\subsection*{\textit{Cum enim luxuriatae fuerint in Christo, nubere uolunt habentes damnationem, quia primam fidem irritam fecerunt. }}\pstart Iuniores uiduae adepta per mortem mariti libertate, discussoque timoris iugo et uerecundiae, expertaque carnis uoluptate, dum licentius uiuere incipiunt, occa sionesque impudicitiae cibo, potu et otio non uitant, periclitari incipiunt, et libidine tentari ac uinci incauta familiaritate et cohabitatione. Sicque sensim prolabuntur  \pend
\section*{*COMMENT. IN I. EPIST. }
\marginpar{[ p.520 ]}\pstart ad petulantiam, luxum et carnis lapsum, scandalizantque et infamant Christi sacrosanctam ecelesiam, et nomen Christianum, quod subinde male auditpropter eas. Tandemque ́, quod secretius perperam egerunt, publica diffamia detegitur, et earum luxuria propalatur. Quam suam ignominiam, dum euadere, contegere et excusare cupiunt, nubere deliberant, matrimonium praetexunt, uel de nouo con trahunt. Habentes damnationem, primum quidem in conscientia, et penes semetipsas, consciae sibi turpitudinis, confusae apud semetipsas, deinde publice reiecisse fidem comprobantur uiuenti marito praestitam, sed tantum per hypocrisim seruatam innotescit signis certioribus, dum fictam fuisse suam castitatem testantur, et non ueram. Nunquam ex animi sententia maritum amasse solum osten dunt facto, uel non ex fidei synceritate continuisse uidentur, sed ex necessitate, dudum aliud facturae uel tentaturae, si per timorem zelumque mariti, et timorem opprobrii licuisset. Si tamen non secreto simile tentasse maritata credetur, facile enim deteriora hominibus de talibus arrident, ac persuadentur eredenda temere, tanquam quae fieri potuerint et facta sint.  \pend
\phantomsection
\addcontentsline{toc}{subsection}{\textit{Simul autem et otiosae discunt circuire domos: non solum otiolae, sed et uerbosae et curiosae, loquentes quae non oportet. }}
\subsection*{\textit{Simul autem et otiosae discunt circuire domos: non solum otiolae, sed et uerbosae et curiosae, loquentes quae non oportet. }}\pstart Duplo delectat has uiduas, et placet licentia, post coactam continentiam. Gau dent enim se a nemine impediri, quo minus exeant domos: solitudo domi displicet, lectioni em̃ et orationi non instant. Inuitatio passim grata fit, dum non esteis domi curanda familia, otiosae desident, parta dilapidant donec quicquam extat, uel stipe ecclessastica contentantur. Otiosae circumcursant per alienas aedes, circumferunt rumores, comportant nugas, augentque suo more mendaciis: insigniter uerbosae, loquaces, et nugaces, uanae et leues, neminem formidantes, libertatem suam in carnis delicias conuertentes. Curiose perquirentes quicquid passim dicitur uel agitur, tempus otiose fallunt uerbis, uanae, leues et dissolutae domum redeunt, in confusionem reliquarum uiduarum, demortuique mariti et propriam. Tandem autem abnegata fide, infidelium foeminarum more uiuentes, aduersus Christum instituunt blasphemiam. Confundunt ac diffamant ecclesiam, insolentiam suam hactenus palliatam produnt tandem, timore tantum poenae dissimulatam.  \pend
\phantomsection
\addcontentsline{toc}{subsection}{\textit{Volo autem iuniores nubere, filios procreare, matresfamilias esse, nullam occasionem dare aduersario, maledicti gratia. }}
\subsection*{\textit{Volo autem iuniores nubere, filios procreare, matresfamilias esse, nullam occasionem dare aduersario, maledicti gratia. }}\pstart Ad praecauendum tot malis muliercularum istarum, uolo magis ut uiduae iuniores non tam libere uiuant, sed denuo marito nubant, ut quae sese regere nesciunt, in officio et honestate per uiros seruentur, amputentur ipsis lasciuiae causae, lenones procosque euadant. Habeant quod domi curent, mariti copia et austeri, tas libidinis medeatur ardori, et ut petulantia compescatur, opprobriumque ac diffa miam effugiant. Filios et filias gignendo carne humilientur, et solicitudine regen di familiam. Habeant occupationes dignas iuuenculis, filios educando, curando, instituendo. Prae illis honestissimis uoluptatibus, alii pruritus ac libidines cessabunt. Occupatae erunt circa domus curam et necessitates. Familiam regant mul ta solicitudine. Matrisfamilias officium negotiosum obeant, die noctuque ́. Exenplo praesint, iubeantque ancillis et filiabus. Vni uiro decenter inseruiant, cui casti tate et uerecundia sanctoque timore placere meditentur. Haec apostoli Pauli uoluntas, non dubium est praeceptum esse, si periculum timeatur incontinentiae, otii, suarumque sequelarum, diffamiae et indisciplinatae conuersationis. Cui concordat ratio, natura, aequitas, et uerbum dei, quod non omnes capiunt. Sic autem nuben\pend
\section*{AD TIMOTH. CAP. V. }
\marginpar{[ p.521 ]}\pstart tes euadent, nedum incontinentiae periculum, luxuriae diffamiam, feruoremque ́ urentis libidinis: sed etiam declinabunt oblocutiones, ne occasio praestetur aduersariis et detractoribus, male uel suspicandi uel loquendi de eis, interimque eccle, sia nihil confundetur. Amputanda enim quantum fieri possit causa est omnis iustae suspicionis, maxime mulieribus, quibus ob sexus fragilitatem facile imponitur: sunt enim propensioribus affectibus communiter et leues, difficile reluctantes cupiditatibus. Vix illibata fama defenditur et euadit periculum diffamiae, etiam dum studiose omnia dispiciuntur, multo minus dum libertati nimium indulgetur.  \pend
\phantomsection
\addcontentsline{toc}{subsection}{\textit{lam enim quaedam conuersae sunt retro post satanam. }}
\subsection*{\textit{lam enim quaedam conuersae sunt retro post satanam. }}\pstart Admonet omnibus nota necessitas, dum uidemus ex huiusmodi uiduis iunigribus quasdam turpiter uiuere, et scandalosam diffamiam incurrisse, uiduarum ecclesiasticarum sancte institutum collegium foedasse, quod licentiosa conuersatione sua propositum castitatis abiecerint, ecclesiam confuderint, postque manifestatam libidinem, et proditam fidem, nuptiis sese excusare, uel in totum post sata nae tentamenta abiisse, experientia docuit. Christo enim cui sese in castitate serui turas instituerant, orationibus et bonis operibus incumbendo seruire proposue, rant, turpiter et cum probro ecclesiae desciuerunt. Quod autem a uoto emisso cuiquam, et rescisso desciuerint, et uoti transgressores fuerint, non scribitur: quod coactae fuerint redire ad pristinam professionem suam, iussae perpetuo seruare quod uouerant uel instituerant, minime habetur in apostolicis scriptis. Non sunt excommunicatae, non incarceratae, nec uiolenter compulsae ad reditum: id enim tacetur ubique . Suo periculo domino suo stabant uel cadebant. Non erantobseratae cancellis, nec, ut nunc, muris clausae, alioqui discurrere et nugari, ut praedicitur, non potuissent, uel abire post satanam. Omnia libere seruabantur. Apostolica tem pora et doctrinae uota omnia ignorabant. Sine uotis etiam primi monachi religiosius uixerunt, praeter uotum obedientiae, quod semper temporale erat, donec aliquoin loco, uel sub aliquo abbate uiuere religiose, et sine animarum periculo uo luissent uel poruissent. Post mille annos ab exordio monachatus, qui tamaen etiam apud gentiles fuisse inuenitur, quando pro stabiliendis examinibus monachorum et monasteriis, pro ceremoniis continuandis, et augendis censibus ac ampliandis uinculum uotorum excogitatum est, quo constringerentur ad perpetuam mansio nem. Sicque definierunt primores ordinum, et obtinuerunt per tyrannidem, laqueos huiusmodi conscientiarum, sine uerbo dei, in impedimentum euangelicae conuersationis et uitae. Comminiscentes sibi friuolas rationes, et detorquentes scripturas, ad probandum longe diuinius et sanctius esse, ex uoti necessitate bene agere, quam ex spiritus libertate. In uotis substantialibus, quae uocant, nullam cadere dispensationem. Subsequuta est crudelitas perpetuorum carcerum, pro asseruandis traditionibus humanis, et excommunicatio Papalis, et infinita huiusmodi fri uola commenta, omnia sine uerbo dei, et in detrimentum uerae pietatis ac religio nis, quae consistit in obseruantia mandatorum dei et uerbi dei. De quibus nunc passim, gratia deo, multa leguntur, ut posthac non tam late pateat nassa capiendi simplices in claustra, et illaqueandi conscientias traditionibus mere humanis et diabolicis.  \pend
\phantomsection
\addcontentsline{toc}{subsection}{\textit{Si quis fidelis habet uiduas, subministret illis, ut non grauetur ecclesia, ut his quae uere uiduae sunt, sufficiat. }}
\subsection*{\textit{Si quis fidelis habet uiduas, subministret illis, ut non grauetur ecclesia, ut his quae uere uiduae sunt, sufficiat. }}\pstart Si quis uir fidelis uel mulier cognatam habeat uiduam, quam de suis faculta\pend
\section*{COMMENT. IN I. EPIST. }
\marginpar{[ p.522 ]}\pstart tibus commode nutrire possit, et uitae necessaria ministrare, non eam patiatur ob trudi ecclesiae, quae solet uere uiduas, id est destitutas omni subsidio alere de communibus prouentibus ecclesiasticis: ne forte dum alendae fuerintomnes ecclesiasticae personae, tam uiduae quam aliae, donationes uel collectae non sufficiant. Iterum uidemus ad quem finem ecclesiasticae diuitiae, decimae et oblationes ordinandae sint, et quae olim fuerit consuetudo fidelium, et quid nostris temporibus emen dandum sit, et uere ecclesiastici qui dicantur.  \pend
\phantomsection
\addcontentsline{toc}{subsection}{\textit{Qui bene praesunt presbyteri, duplici honore digni habeantur, maxime qui laborant in uerbo et doctrina. }}
\subsection*{\textit{Qui bene praesunt presbyteri, duplici honore digni habeantur, maxime qui laborant in uerbo et doctrina. }}\pstart Sunt qui male praesunt presbyteri, episcopi scilicet et diaconi, nulla cura occupati pauperum uiduarum, uel de uerbo dei praedicando, sibi solis uiuunt, otiantur, libidinantur, fastum gerunt, ambiunt honores et quaestum, bona tantum absumunt ecclesiastica. Quid talibus debetur aliud, que utcum sint sal infatuatum, et ad nihil ualens ultra, eiiciantur foras de sedibus suis, excommunicentur, confundantur, publice iudicentur, qui suis insolentiis sese iamdudum publice accusarunt scandalosa actione, et otio, uita et moribus. Conculcentur autem pedibus, id est, honore destituantur, commodis, diuitiis, officio et publico dedecore afficiantur quod merentur: ne si permittantur in scandalosa uita sua, negligatur popu lus, obtendatur peccantium authoritas, et patrocinium iniquitatis, fiatque ; sicut sa cerdos, sic et populus, irruatque lupus in oues, et dispergat, cum ex pastore factus fuerit lupus uorax, mercenarii quoque inierint foedus cum lupis, in gregis perniciem. Qui uero bene praesunt episcopi, qui hic dicunt et fere ubique a Paulo presby teri, nisi ubi manifeste aetatem solam indicat nomen presbyteri. Tales, inquit, duplici honore digni habeant. Episcopi uero ut bene praesint, necesse est esse presby teros, id est seniores, non tam aetate, q̃ morum grauitate, sapientia, et multarum rerum experientia, non iuuenilibus affectibus obnoxii. Quales erant Timotheus atque Titus, tales episcopi sensu cani et sacris eruditi sermonibus, duplici honore digni habeantur. Honor dupliciter sumitur, pro uulgari reuerentia, ut estuia cedere aut loco, nudare caput, flectere genu, quae gratitudinem indicant et appretiationem, agnitionemque donorum dei in honorato:estque magis honorantis q̃ ho norati. Aliter, ut hoc loco, honor est munus, eleemosyna, sustentatio, necessariorumque prouisio sufficiens et honesta, quae debetur doctori: non ut profundat in usus mundanos, sed ut honeste uiuat et modeste, ut sit exemplar omnis parsimoniae, maturitatis, et contemptus rerum mundanarum, et speculum uitae Christia nae, hoc est euangelicae, imo apostolicae. Hoc duplici potius honore, quam duplici pro uisione et annona honorentur, cum nihil superflui illos deceat, sed eatenus uictum et quibus tegantur habentes, exercitia pietatis praestare possint, et familiam necessariam educare, minime etiam otiosam, sed ad sacra officia instituendam pro uirili. Si qui uero fuerint presbyteri et antistites praefectiue ecclesiae, qui non soIum praesunt integritate uitae, consilio prudentiaque administrandi ecclesiastica bo£ na, sed etiam doctores se exhibent, ut laborent legendo, docendo, exhortando, in culcando, et dispensando doctrinam ecclesiasticam, euangelicam et apostolicam, uerbo et doctrina prosint insigniter uti praesunt. Istis duplex hic honor omnium maxime referatur, tantoque diligentius, quanto sibi suisque sacris illis addicti exercitiis spiritualibus, minus familiae suae educandae prospicere possunt, et de uitae ne cessariis prouidere. Duplici etque honore dignus accipi potest, quia eo magis quae sufficiant simplicem ac honestam suppetentiam assequant. Quasi dicere uelit, Hii ta\pend
\section*{AD TIMOTH. CAP. V. }
\marginpar{[ p.523 ]}\pstart les diligentius excolantur, et fidei laborumque suorum fructum uberiorem assequanf. Non enim duplum pretiosarum uestium, uel ferculorum sibi comparare docentur, qui frugalitatis specimen esse debent. Hospitalitatis enim et eleemosynarum materia aliunde sumenda foret, de bonis scilicet ecclesiasticis, cuius insuper dispensator sit idoneus.  \pend
\phantomsection
\addcontentsline{toc}{subsection}{\textit{Dicit enim scriptura, Non infrenabis os boui trituranti. }}
\subsection*{\textit{Dicit enim scriptura, Non infrenabis os boui trituranti. }}\pstart Mos est prouinciis multis, bobus uel equis circumactis trituram facere, quibus sic magnopere laborantibus, crudele fuerit negare cibum, quem aliis parant tanto studio ac sudore. Omnino enim qui ecclesiae opus exequitur, qui docet, instruit, hortatur, sacramenta ministrat, priuatim seruit uel publice, is dignus estut ab ecelesia nutriatur et foueatur. Quid autem de illis sentiendum, qui nihil laborant huiusmodi, nunquam sacra leguntut doceant, non praedicant, non monent, non consolantur conscientias? Quorum potius omne opus est, in ligandis inuol uendis conscientiis, praeceptis et excommunicationibus: ad populum non loquun tur, nisi suas leges: toti sunt in colligendis, recensendis, augendis et absumendis sumptibus, et pecuniis numerandis. Non solum non praedicant euangelium, sed nec praedicari sinunt, si impedire possunt, nisi coacti timore amittendi luxum, quae stum et fastum, Paulus hic tractat et haec citat de dispensatoribus uerbi dei, non de cantoribus, organistis, monachis, uel nonnis. Quod si in tanto numero hodie ecclesiarum presbyteri sacras in collegiatis ecclesiis scripturas, utolim institutum erat, legerent, ad aliquem legitimum fructum, uel audirent ad eruditionem capessendam, ut aliquando idonei surrogari possent ad officium pastorale, digni certe iudicarentur suo uictu. Nunc autem boues triturantes, id est agricolae ipsi, toto anno laborantes, fraudantur durissimo labore suo, uinitores aquam bibere coguntur, cum uxoribus, filiis et famulis dura inedia macerantur, ut uix uiuere et spiritum trahere possint. An non est hoc boum ora constringi et alligari ? Interim clericos otiosos et scandalosissimos fouere, in perniciosissimum scelus mundi. Quin et inter clericos hi maxime egent qui doctiores sunt, et magis in ecclesia laborant, quorum fructus propagatos deuorant maiores Canonici, praelati, aut Curtisani. qui uix digni sunt ut a terra sustineantur, animalia uentris, gloriae, et otio turpissimo scandalum inferentes ecclesiae, de quibus multo plura dicenda, sed uerbum domini omnia huiusmodi declarabit et iudicabit.  \pend
\phantomsection
\addcontentsline{toc}{subsection}{\textit{Dignus est enim operarius mercede sua. }}
\subsection*{\textit{Dignus est enim operarius mercede sua. }}\pstart Indignum graueque piaculum est, laborantem fraudare sua mercede, uel usque  etiam ad noctem. Operarium autem dicit, non otiosum, non stertentem prae pigritia, non aliorum laboribus, partis inhiantem, et solis uoluptatibus operam dantem. Non sunt operarii, qui tota die in foro stant otiosi, sed qui pondus ferunt laborando diei et aestus. Non qui ludunt alea, et colludunt mulierculis, ac illudunt uino, et lautitiis dissoluuntur: qui otiosi, ad nihil utiles, laborantium sudores exigunt, et necessariis fraudant. Operarios dicit, qui laborant in uerbo et do ctrina, qui bene praesunt, qui gregi aduigilant, qui in ecclesiae utilitatem proficiunt. Hodiernus autem cantus et soni, uigiliae et Missae emptitiae, nec deo placent nec prosunt hominibus. Impediunt magis a melioribus, aboratione, a uerbo dei ediscendo, audiendo, docendo, nec populo nec clero quacunque in re utilia, in eum usum a patribus nunquam institutum fuisset hoc officium. Merces autem respon detobsequio et operi. Fraus itaque dicetur et dolus, rapina quoque , mercedis nomi ne exigere, cum nihil egeris utile uel gratum cuiquam. A plebe catholica susten\pend
\section*{COMMENT. IN I. EPIST. }
\marginpar{[ p.524 ]}\pstart tari uelle, cui non seruias, nihil prosis, sed obsis maxime et impedias profuturos, aliorum potius tibi rapias mercedes laborum, de sanguine quoque et sudoribus aliorum sagineris. Vt hodie cernere est in praelatis ecclesiarum et beneficiosis, qui cum sint sceleratissimi, cleri et populi oppressores, et Christi opprobrium, uiolen ter deglutiunt et praeripiunt mercedes pauperculis pastoribus ecclesiarum, et ad ructum uorant. A quibus, si quando uicarii populo et ecclesiis, pauperes et idonei praeficiuntur sacerdotes et praedicatores, hii non diutius feruntur in officio, quam ab euangelio pure, libere ac fideliter praedicando destiterint, statim abiiciendi, dum Christum et Christianam disciplinam ac doctrinam, ipsum scilicet uer bum dei, simul ac fidem pietatemque coeperint asserere, uolentes impedire, ne populus sapiat aliquando, et mercedem deneget otiosis indignum, tantumque de factonullo iure debitum. Digniores plane, ut alligati asinariae molae, abiiciantur iussu Christi domini, in profundum maris, et ut fatuum sal conculcentur ab hominibus. Ligatis quoque manibus et pedibus, mittantur in tenebras, ut partem habeant cum infidelibus in saeculo futuro.  \pend
\phantomsection
\addcontentsline{toc}{subsection}{\textit{Aduersus presbyterum accusationem noli accipere, nisi sub duobus aut tribus testibus. }}
\subsection*{\textit{Aduersus presbyterum accusationem noli accipere, nisi sub duobus aut tribus testibus. }}\pstart Non omnino certum, an de episcopo loquatur, an de seniore. Videtur autem de utrisque accipiendum. Episcopi enim sunt seniores de ecclesia. Seniores prospicere debent, ut prudentiores toti ecclesiae pro authoritate senii. Vnde et Senato res uel Senatus dicuntur spectabiliores de ecclesia, etiamsi non uncti sint et rasi. Denique tot populi tam magnae ciuitatis, unius prudentiae et spiritui niti non debent, sed de seniorum consilio, quasi senatum habeat, cum quibus iuxta dei uerbum et regulam omnia tractet. Sintque unius ciuitatis hoc modo plures episcopi, imo presbyteri non tam aetate quam sensu et moribus, qui omnes in authoritate habeantur, et sint in congregatione fidelium: etiamsi non omnes doceant ac praesint uerbo et doctrina, assistunt tamen episcopo, et presbyteri sunt plebis honorabiles et seniores ecclesiae. Non enim legimus discrimen fuisse apostolorum temporibus intra ecclesiam, ut pars una religiosa, clericalis, spiritualis, consecrata, et ec clesiastica fuerit: alia autem ecclesiae Christianae pars maior, prophana, saecularis, mundana, carnalis et laicalis habita et dicta fuerit. Fuerunt enim omnes Christiani religiosi, spirituales, utentes hoc mundo, tanquam non utentes, uacantes piis operibus, et in sua lingua quique sacras scripturas legentes uel audientes congruis temporibus, et insuper labore manuum uictitantes. Id enim docuit Paulus apostolus, uas electionis, et insignis apostolus uerbo et exemplo, Vt qui non laborat in ecclesia, non manducet. Quod si qui saeculariter, mundane, carnaliter, et sine timore dei, et sine ecclesiastica disciplina uiuere praesumpsissent, Christiani non habebantur, sed post satanam ire iudicabantur, et euitabantur tanquam haeretici, et ad sacrosancta mysteria fidelium non admittebantur. Seniores itaque Christianorum bonorum presbyteri dicebantur, inter quos qui uerbo dei et doctrina ua Iebat, et donis sancti spiritus nobilior et conspectior, is aliis prospiciebat, praeferebatur, et episcopus dicebatur et erat: magis officio docendi praecellens, que pote state et imperio, fastu uel pompa: honorificout senes decet uestitu, secundum loca et regiones, ornati procedentes, sine curiositate, pretio uel superfluitate notabiles. De istis presbyteris ecclesiae honoratioribus, ex authoritate, probitate et sapientia praestantioribus dicitur, Aduersus presbyterum et grandaeuum, maturum uirum et spectatae uirtutis personam, accufationem non recipias, nisi sub duobus  \pend
\section*{AD TIMOTH. CAP, V. }
\marginpar{[ p.525 ]}\pstart aut tribus testibus. Poterant insolen tes, a talibus presbyteris correpti et admonis ti, facile malignari, et comminisci mendacia, redundaret autem hoc in probrum ecclesiae, eleuaretur authoritas uirorum grauium. Ideo non facile admittendi calumniatores uel audiendi, ne ansa praebeatur leuibus personis, temere lacescendi seniores ecclesiae. Omnes Christiani pie religioseque uiuentes, ecclesia nominabantur. Licet non abessent hypocritae et pseudochristiani, sed de secretis cordium deo non ecclesiae iudicium defertur. Non uno uel duobus testibus, sed pluribus tam spectabiles uiri, a facie et conuersatione accusari debebant. Ne facile concin netur mendacium contra honestos ab indignis et leuibus, uindicare se uolentibus de castigatione digna. Sicut de Quadrato legimus in historia Alexandri Papae san cti:ideo probabiliore testimonio opus est contra praelatos et magistros aliorum, conuincendos de inhonestate.  \pend
\phantomsection
\addcontentsline{toc}{subsection}{\textit{Peccantes coram omnibus argue, ut caeteri timorem habeant. }}
\subsection*{\textit{Peccantes coram omnibus argue, ut caeteri timorem habeant. }}\pstart Hoc eadem ratione ad presbyteros referendum, ut si quis senex impius publice insolenter peccet, ut fide dignis testibus et multis crimen testari possit, cum id cedat in scandalum plurimorum, et totius ecclesiae et Christianitatis. Hii sic con uicti, publice et coram omnibus argui debent, ut timorem habeant caeteri, et per sonarum acceptio nullum locum prorsus inueniat in ecclesia dei. Neque omnia hy pocritice in ecclesia agi uideri possint ab infidelibus, utque criminum castigatio scan dalum tollat ecclesiae innocentis: quae non tam ex innocentia uitae, quam ex zelo quoque  diuinae gloriae, et Christianae honestatis dinoscitur. Non sunt quidem omnes in ecclesia innocentes, criminum autem osores, uindices et zelatores, omnes esse debent uere ecclessastici. Cumque uiderint extranei, etiam presbyteris et maioribus non parci uel conniueri, sibi quoque minores non parcendum nouerint, si fuerint deprehensi in ecclesia criminosi. Nihil hic praescribit diuus Paulus de correctione episcopi in specie, tum quia Timotheus probatae erat et habebatur sanctimoniae, tum quia in primitiua ecclesia episcopi assumebantur ex pluribus sanctis san ctiores. Verum presbyteri nomine satis indicatur, coram pluribus et omnibus arguendos etiam episcopos, si scandalosi fuerint. Quantoque lues contagiosior, si episcopus obnoxius sit criminibus, tanto adhibitis sufficientibus attestationibus, durius ac seuerius de eo iudicandum erit. Hic iterum et semper episcopi et presbyteri nomine, ecclesiae praefectus et parochus quilibet intelligi debet, non magna rum tantum urbium, et praefecturarum antistites. Ex horum enim doctrina et moribus pendet profectus detrimentumque ecclessae:longeque maturiore consilio deligendi sunt testimonio plebis piae, et non tumultuose suffragio, q̃ magistratus ci uitatum. Vbi non modica adhibenda diligentia est, utquam integer fama ac moribus, rerumque experientia, et spectata probitate probatus et conspicuus deligatur: nisi maioris aestimare uelint rectorem ciuium, quam doctorem animarum. Verum sicut superius Paulus de Viduis, quae uere uiduae essent, tractans, eas specialiter docuit, et tamen generalem regulam scripsit de domesticis omnium maxime fouendis: sic hoc loco uidetur, loquens de presbyterorum arguitione et contestatione, sub inferre scilicet regulam generalem, de omnibus publicis et scandalosis in ecclesia peccatoribus, ut tales coram omnibus ab episcopo arguantur, quo caeteri timorem habeant similia committendi. Sicut et gnome est, crimen non facile imo nun quam recipiendum, nisi appareant et conuincant illud clare testes idonei. Sicut et in lege Mosis, et passim in sacris literis, generales regulae notantur, cum de ne gotiis et personis spiritualibus sermo in cidit, a quibus nemo episcoporum exclu\pend
\section*{COMMENT. IN I. EPIST. }
\marginpar{[ p.526 ]}\pstart ditur, omnium autem maxime is qui uult haberi et dici maximus: tantum abest ut feratur ab ecclesia, qui ecclesiam tam grauiter scandalizat.  \pend
\phantomsection
\addcontentsline{toc}{subsection}{\textit{Testificor coram deo et Christo lesu, et electis angelis eius, ut haec custodias sine praeiudicio, nihil faciens in alteram partem declinando. }}
\subsection*{\textit{Testificor coram deo et Christo lesu, et electis angelis eius, ut haec custodias sine praeiudicio, nihil faciens in alteram partem declinando. }}\pstart Sub contestatione generali immensae prouidentiae dei, omnia cernentis et gubernantis: et domini lesu Christi, qui ut uerbum dei patris, aeque immensus, nobiscum semper est, et ubique praesens usque ad consummationem saeculi, ut sit caput ecclesiae et rector, omni praeditus potestate in coelo et in terra, patri consubstantialis: et quasi in prospectu angelorum, et spiritualium ministrorum dei, qui fide lium sanctis negotiis iugiter intersunt, Ego te, ô Timothee, adiuro et obtestor, ut praeceptor discipulum, ut non patiaris aliquod fieri praeiudicium uerbo dei, quod me referente et docente audis ac suscipis, quin praecise insistas illi, tanquam regu lae ueritatis, et Christianae doctrinae, ne quomodolibet detrimentum patiatur pura ueritas, labefactetur et adulteretur, sicut fit a pseudo discipulis. Nihilque committas declinando in alteram partem, hoc est, praeter fidel regulam et scripturae: nihilque doceas unquam iuxta humanum sensum, et iudicium carnis. Vel secundum alios, ut nihil iudices, etiam sub testibus paucis uel pluribus, cum praecipitatione iudicii, inconsiderate, friuole uel inconsulte, utpote aetate iunior, et feruentiore ingenio, et praeuolatili. Ne iudicando indulgeas affectionibus, et inclinationi na turae, ac propensiore animo in ea quae libent et placent, sed circumspecte iudices, mature, praehabita matura deliberatione, non iuxta naturae sensum, et personarum acceptionem, sed suste, aequaliter et sidellter pronuncies. Sicut constitutus in oculis domini dei et Iesu Christi, et angelorum, ut similiter nos quoque a domino et Christo Iesu iudicandos cogitemus. Discimus hoc loco primum, quam inuitus ueniat Apostolus ad mandandum ac praecipiendum, obtestatur potius, et precibus contestatur: deinde semper deum habere prae oculis, ut iudicaturum nostros cogitatus et opera. Nihil praecipitanter, nihil praeter regulam uerbi dei, nihil ad gra tiam uel displicentiam personae, sed luste iudicet episcopus. Angelos quoque del di scimus mitti in ministerium et testimonium actuum nostrorum, quos et reuereri semper debemus, et cauere a peccatis, non solum deo et Christo domino aspicientibus, sed etiam angelicis exercitibus, qui assistunt semper nobis, et contenplantur actus et uerba, interim tamen in coelis patrem aspiciunt, et beati fruuntur. Id uero episcopo faciendum non solum in iudiciis, sed in omnibus administrationibus, ut prouide, fideliter, syncere et candide agat omnia, is qui praeest eccle siae episcopus.  \pend
\phantomsection
\addcontentsline{toc}{subsection}{\textit{Manus cito nemini imposueris, neque communicaueris peccatis alienis. }}
\subsection*{\textit{Manus cito nemini imposueris, neque communicaueris peccatis alienis. }}\pstart Apostoli mittentes ad fidem praedicandum aceuangelium aliquos, uel officium uerbi imponentes, utebantur manuum impositione, in signum approbatio nis acceptum a Christo, dabaturque talibus spiritus sancti donum, ad hoc idoneum, ac sufficiens in ecclesia. Quia uero tunc subintrabant multi pseudapostoli, et ma li Christiani sese admiscebant, obque quaestum et gloriam, opus apostolicum tenta bant, ideo Apostolus praemonet Timotheum, ut prospiciat sibi, ne his officium praedicationis mandet ac delegat, nisi prius fidem, mores studiaque digna probe co gnouerit, ne malae et erroneae doctrinae illius et seductionis, ipse causa uideri possit, qui tales ad id ordinauerit: neque error humanus et falsitas spiritus sancti doctri na aestimetur, qui manuum impositione uel dabatur, uel dari significabatur. Erat enim attestatio quae dam sanctitatis in doctrina ac uita, quae non cito exhiberide.  \pend
\section*{AD TIMOTH. CAP. V. }
\marginpar{[ p.527 ]}\pstart bebat. Qui enim impium doctorem instituerit, errorum illius et seductionis in populo causa erit, taliumque peccatis et erroribus communicabit semper. Summo pere igitur discutienda fuerit praedicatorum doctrina, studium, mores et uita, hoc est praelatorum ecclesiae, quorum primarium officium est, docere uerbum ueritatis.  \pend
\phantomsection
\addcontentsline{toc}{subsection}{\textit{Teipsum castum custodi. }}
\subsection*{\textit{Teipsum castum custodi. }}\pstart Caste hactenus et pure uixisti, ô Timothee, iuuentutemque tuam disciplinis se uerioribus instituisti, ad facilius castitatem conseruandam, eam posthac quoque di ligentius serues: non ut tuapte uoluntate castus sis, cum dei solius donum sit casti tas uera, sed in seruitutem redigere, ac abstinentiis crucifigere, et comprimere fo£ mitem studeas ut hactenus, ut ea ualde conducunt castitati seruandae, in integra potissimum aetate, dum sanguis feruet et caro. Verum sic serua abstinentiam, ut ne quid nimis, ne per corporis imbecillitatem, et nimiam attenuationem periclitetur corporis ualetudo, et impediantur opera necefsariae charitatis. Ideo subdit consequenter,  \pend
\phantomsection
\addcontentsline{toc}{subsection}{\textit{Noli adhuc aquam bibere, sed modico uino utere, propter stomachum tuum, et frequentes tuas infirmitates. }}
\subsection*{\textit{Noli adhuc aquam bibere, sed modico uino utere, propter stomachum tuum, et frequentes tuas infirmitates. }}\pstart Ne posthac, ut hactenus consueuisti, solam aquam bibas: quae licet non omni bus, tui tamen corporis obest ualetudini, ne dum forte nimium castigas corpusculum, sub onere magnopere portando succumbat. Non propter se bona et eligibilis est maceratio corporis, sed propter uirtutum exercitia. Aqua quidem, non asue tis multis rationibus incommodare potest. Et insuper tui corporis conditio non fert nimiam. Vino tibi opus est, non omnibus: nec id ad libidinem gustus, sed ob stomachum debilem. Eo igitur posthac utere ad necessitatem, sed modico: nisi enim debilitas id exigeret corporis, uini abs te potum neutiquam exigerem. Nos, etiam religioss professione et habitu, puelli etiam uinum bibimus, adolescentes facti po tamus, uiri inebriamur. Monachi et clerici nefas et miseriam putant aquae potum, et poenam intolerabiliorem. Episcopi non nisi exquisita uina ac electa potant. Sic praestamus apostolicos. Nimius quicunque potus, etiam aquae, in uitio est, mul to magis uinum. Non quid bibas uel comedas, sed quantum curandum prudentibus et religiosis, imo et medicis. Modus est potationi ponendus, ut quantum seruiat incolumitats bibatur. Non mittit Apostolus ad medicas purgationes et con fortationes, sed dietam indicit et parsimoniam, optimam studiosorum medicinam.  \pend
\phantomsection
\addcontentsline{toc}{subsection}{\textit{Quorundam hominum peccata manifesta sunt, praecedentia ad iudielum: quorundam autem et subsequuntur. Similiter et facta bona manifesta sunt: et quae aliter se habent, abscondi non possunt. }}
\subsection*{\textit{Quorundam hominum peccata manifesta sunt, praecedentia ad iudielum: quorundam autem et subsequuntur. Similiter et facta bona manifesta sunt: et quae aliter se habent, abscondi non possunt. }}\pstart Praeuenire dicuntur iudicium, et non indigere contestationibus, ea hominum multorum peccata, quae manifeste, inuerecunde et audacter committuntur, in prospectu hominum, et totius ecclesiae. Huiusmodi peccata ad publicum etiam iudicium delata punĩantur: nec ibi multis opus fuerit testimoniis, et inquisitionibus scrupulosis. Illa quoque tu solus punire potes, et ecclesia debet. Quae uero peccata sequuntur iudicium, id est nesciuntur nisi idoneis et paucis testibus coram iudicibus comprobentur: haec possunt abscondi, et secreta admonitione uel correctione corrigi ab episcopo et aliis, quae alioqui deo iudici et extremo iudicio seruabuntur. Non omnia hic puniri et emendari possunt. Non sunt de foro episcopi uel ecclesiae nisi manifesta, Similiter et bona opera.) Quaedam manifeste  \pend
\section*{COMMENT. IN I. EPIST. }
\marginpar{[ p.528 ]}\pstart siunt, non egentia multa discussione: nec sinistre criminari possunt, ob suspicionem occultae imperfectionis, quasi scilicet, licet bona sint, tamen quia mala conscientia fiant, ideo ueniunt improbanda uel abscondenda. Licet enim omnia nostra opera re uera sint defectuosa, et si diuino sistantur iudicio expendenda, mala deprehendantur. Quia tamen speciem mali non habent, non sunt obmittenda uel occultanda. Quae uero aliter se habent, bona uidelicet sunt in apparentia, penitius autemomnino mala; illa suo tempori et iudici abscondi non poterunt. Nisi quis sic intelligere hunc locum uelit, ut dicatur, Homines quidam nunc sic uiuunt, ut ex eorum operibus, manifeste improbis, antequam moriantur, uel a deo iudicentur damnandi, a nobis peccatores agnoscantur. Quorundam peccata ma la quidem sunt, sed in hac uita sic fiunt, ut uel dubia uel occulta cum sint, non iudicari possint ab hominibus, deo autem soli post hanc uitam iudicanda reseruantur: de quibus nihil ad nos, suo domino stant uel cadunt. Vicissim hominum quo rundam opera clare bona uidentur, a nemine nisi improbo criminanda, ea a nobis laudanda sunt, de eis nemo ob uerecundiam praepediri, uel propter sinistrum oculum obmitti: quandoquidem sunt qui omnia bona opera infamare gaudent et student. Sunt tamen et opera in speciem tantum bona et hypocritica, quae ut deo manifesta sunt et ignota hominibus, suo tamen tempore apparebunt qualia fuerint. Nihil enim opertum, quod non reueletur. Abstinentia autem a uino, et similia, nec manifeste bona sunt, nec mala, sed pro iudicio del et rationibus secre tis, necessaria uel indigna, grata uel ingrata, suo tempore apparebunt, et discutientur a iudicio diuino.  \pend
\endnumbering\beginnumbering\section{CAPVT VI.}
\phantomsection
\addcontentsline{toc}{subsection}{\textit{\huge\textbf{Q}\normalsize Vicumque sunt sub iugo serui, dominos suos omni honore dignos arbitrentur, ne nomen domini et doctrina blasphemetur. }}
\subsection*{\textit{\huge\textbf{Q}\normalsize Vicumque sunt sub iugo serui, dominos suos omni honore dignos arbitrentur, ne nomen domini et doctrina blasphemetur. }}\pstart Emptitii seruuli infidelium dominorum abidololatria conuersi, ad unius ueri dei fidem, patienter subinde ferant obedientiae iugum ac laborem, et sortem conditionis suae, et statum miserum. Nec quia ad Christum conuersi sunt, fraudare tentent dominos suo seruitio, ut transfugae fiant et profugi ob fidem dei, quast contemnentes dominos ob infidelitatem, quibus se Christianos seruire secundum corpus indignum putent, sed magis obsequiis humilibus seruituti satisfaciant, ab omni offensa temperantes sibi, honore debito colentes dominos suos, et dignos arbitrantes, a deo sibi praepositos cogitantes, non casu obtigisse uel for tuna. Sed sic esse domini uoluntatem, ut boni exerceantur a malis, ut mali, cum deo placuerit, conuertantur a bonis. Vt inclarescat uirtus seruorum dei, et Chri sti gloria a suis sectatoribus illustretur, ut seminarium seditionis non detur ethnicis, et occasio malignandi infidelibus contra pios. Ne nomen dei optimi maximi et solius, per idololatras blasphemetur, quasi non aequus sit, qui iubeat dominos suis destitui seruis, et passim gentiles irritari a Christianis seruis profugis, sicque pa cem infringi, obedientiam cassari, maiores sperni, promissa negari, commercia rumpi, debita non reddi, deumque uerum insolentia seruorum illius apud impios blasphemari. Similiter et doctrina euangelica infamabit, si eius professores mun di ordinem, et magistratuum dominorumque in temporalibus obedientiam detre ctauerint, et omis religio Christiana propter eos male audiet, et in Christi doctri na scandalizabuntur idololatrae, Qui potius ad mentem meliorem, exemplofidelium aedificandi fuerant, alliciendi et conuertendi, unaque nobiscum aliquando dominum deum nostrum glorificaturi in seruis suis. Voluit et Iudaeos dominus  \pend
\section*{AD TIMOTH. CAP. VI. }
\marginpar{[ p.529 ]}\pstart obedire regibus infidelibus, pro eis rogare, tributa pendere, pacem fouere. Et Chri stus Caesari iussit reddi, quae debita sunt Caesari. Petrus quoque in sua epistola multa huiusmodi. Aufert quidem baptismus, id est professio Christiana, peccati domi nium, sed abheri dominatu et saecularibus obligationibus non eximit, siue id graue seruitium sit et tyrannis, siue leue, et secundum saeculum licitum, non est tunc disputandum, sed potius humiliter obsequendum. Si uero seruus commode possit et honesto modo liber fieri, non obmittat. Si dominus seruum manumiserit, ut et ipse a seruitute peceati liberetur a deo, optime fecerit, et ad id est inducendus opportune, non contemptu, uiolentia, temeritate et diffamia euangelii, iugum erit excutiendum.  \pend
\phantomsection
\addcontentsline{toc}{subsection}{\textit{Qui autem fideles habent dominos, non contemnant, quia fratres sunt, sed magis seruiant, quia fideles sunt et dilecti, quia beneficii participes sunt. }}
\subsection*{\textit{Qui autem fideles habent dominos, non contemnant, quia fratres sunt, sed magis seruiant, quia fideles sunt et dilecti, quia beneficii participes sunt. }}\pstart Si qui conditione mundana, et diuino iudicio imperscrutabili serui facti sunt. et seruili sorte oppressi, si dominos sortientur uel habuerint Christianos, non ideo cosdem minus reuerenter tractent, non segnius obediant uel seruiant, quasi indi gne ferendo parem, Christianum uidelicet Christiano obsequi. Pares equidem sunt in Christo, et fratres in fide, sed non in huius uitae et saeculi fortunis, sicut et in dei beneficiis non omnes sunt aequales. Sic enim unicuique homini deus prospexit, et praeordinauit uitae ordinem, statum et exercitium, in quo creatori suo seruiat, commensurauitque grauamina, ut incertum habeatur, seruorum ne an dominorum conditio generaliter tolerabilior fuerit, pro dotibus naturae tam differenter passim impertitis. Quin potius tanto promptius et alacrius eisdem seruiant, quanto ob Christi religionem mitius et mansuetius tractantur a Christianis dominis serui Christiani, et ad prophana et sacrilega seruitia minus compelluntur, sed pacifice quae Christianae religionis sunt exercitia sancta permittunt. A quibus et diligun tur in Christo ut fratres, instruuntur ut discipuli, et tanquam filii beneficiis afficiuntur, et sunt insuper omnium bonorum participes in Christo. Duplici nomine promptius Christianis dominis seruiant quam ethnicis, tum quia domini eorum sunt, tum quia socii religionis seruulos suos non contemnunt. Si enim fideles sunt domini, seruos suos tractabunt ut fratres potius quam ut seruos: diligent magis, contemnent minime: beneficiis afficient, et immoderatis grauaminibus minus oppriment: bonorum suorum eos participes aestimabunt et facient. Quod si crudeliter et absque humanitate tales seruos tractauerint domini, ostendent facto non se tam Christianos esse, quam Domitianos et perfidos. Tolerandi tamen nihilominus sunt domini etiam discoli, quandiu contra deum et Christianam iustitiam, et Christum uniuersorum dominum nihil imperauerint. Alioqui deo oportebit magis obedire quam hominibus, quamlibet mundana potestate contra deum turgidis.  \pend
\phantomsection
\addcontentsline{toc}{subsection}{\textit{Haec doce et exhortare. }}
\subsection*{\textit{Haec doce et exhortare. }}\pstart Talia episcopum docere decet, non solum maiores et dominos, sed etiam infimae conditionis et sortis homines. Seruos hominum non negligere. Docere autem ea quae pertinent ad pietatem, ad tranquillitatem, publicam pacem, ac ad iustitiam etiam humanam, ne suscitetur dissidium in ecclesia, ne per inobedientiam subditorum, ordo ciuilis et humanus confundatur. Exhortari quoque debet episcopus eos quos uiderit cessantes et cunctabundos, quo crebrius admoniti et consolati subditi atque serui, leuius portent onus a deo impositum, pro quo dei fauorem ac beatitudinem consequi possint. Non docet eos libertatem carnis et mundi, quae  \pend
\section*{COMMENT. IN I. EPIST. }
\marginpar{[ p.530 ]}\pstart saepe praecipitat, sed quae in Christo est, patientiam et longanimitatem.  \pend
\phantomsection
\addcontentsline{toc}{subsection}{\textit{Si quis aliter docet, et non acquiescit sanis sermonibus domini nosti Iesu Christi, et ei quae secundum pietatem est doctrinae, is inflatus est, nihil sciens, sed insaniens circa quaestiones et uerborum pugnas. }}
\subsection*{\textit{Si quis aliter docet, et non acquiescit sanis sermonibus domini nosti Iesu Christi, et ei quae secundum pietatem est doctrinae, is inflatus est, nihil sciens, sed insaniens circa quaestiones et uerborum pugnas. }}\pstart Aduerbium Aliter refertur ad apostolicam doctrinam, quae tam in superioil bus quam subsequentibus describitur. Si quis, inquit, iuxta proprium sensum, et leges nationum, iuxta erroneam hominum rationem ac philosophorum opiniones docuerit, et non credit uel appretiatur aut acquiescit ex conscientia doctrinae salutari et ueritatis, sanis, immaculatis aequissimisque sermonibus domini nostri Iesu Christi, sed introduxerit alios etc. Tantum enim istis sermonibus adhaerendum est, Christi scilicet domini, haec tantum docenda, haec sola doctrina est pietatis, secundum quam domino seruire et placere poterimus fideles, de qua sola certi sumus, quia Christi est, qui ueritas est, qui in apostolis et sanctis loquitur, a quo ac cepi quod et trado. Eam scire ac sapere, pietas est:eam negare, impietas. Is (ait) inflatus est.) Superbia siquidem est, sensum proprium praeferre dei uerbis. Omnis enim eruditio alia a fide Christi inflat, et tumorem mentis excitat:in se confidere, de proprio iudicio et sensu praesumere facit, contemptui habet omnem simplici, tatem spiritus, scripturae et stili. Cumque crasso carnis intuitu eam attingere non potuerit, despicit, et stultam pronunciat. Non enim quae spiritus sunt carnales ca piunt, sed dum credunt se sapientes esse, stulti nimirum fiunt. Nihilsciens.) Hoc unum sciet, si quicquam scit, quod nihil sciat. Qui deum ignorat, qui in Christum scopum sacrae scripturae ueteris et noui testamenti finem non credit, qui nititur humanae rationi erroneae, cum sit omnis homo mendax, et comparatus iumentis insipientibus, quid scire poterit ? Scientia non est nisi certae ueritatis agnitio. Chri stus autem ueritas est, et lux mundi, praeter quemomnia tenebrae sunt, et ignorantia atque stultitia. Sed insaniens circa quaestiones et pugnas uerborum.) Lon go quidem tempore post apostolorum excesfum, quaestiones et pugnas uerborum paucas fuisse, licet nunquam defuerint, satis clarum est: sed et haereticorum dein de discordias et pugnas nemo nescit. Verum perinde ac nunquam fuerint per Vniuersitates nostras, haec disceptandi stultitia disseminata est, ab illarum infelicissima hactenus matre Parisiensi. Quae et insanire fecit omnia illustria ingenia, et languere iamdudum et turpiter circa quaestiones futiles, et insanas uerborum pugnas, ex quibus tot miseri quaestionarii et sophistae prodierunt, nihil certum uel integrum relinquentes, sed dubiis quaestionibus, ac tot suis sectis omnia inuol uentes. Iusto itaque iudicio dei palpabiles istas tenebras inciderunt, quia sanos sermones domini neglexerunt, sacras literas pro minimis duxerunt, subsannauerunt, et philosophicis suis stultitiis postposuerunt. Si quis ex ea scriptura sacra quid dam in Sorbonicis suis sordibus allegare uoluisset, contra suum Aristotelem uel Lombardum, prophetalem uel apostolicam sententiam, contra sua decreta et articulos, statim scrupulos innectere studuerunt suis distinctionibus, et sophisticis diuerticulis, ad inuoluendum claram ueritatem, uere Scoticis tenebris: insuperque  ridere consueuerunt, uelut indoctum theologum, qui sacras scripturas diligentius tractare uoluisset. Non est ut id inficias ire possint iamdudum sic actum, cum solenne illis disputatoribus sit, cauere ne authoritate scripturae utantur in suis Sorbonicis disceptationibus: quandoquidem in dubium sensum trahi consueuerint in istis scholis. Quae si theologicae dici mereantur, et non potius diabolicae, pruden tis esto iudicium. Et tamen sic esse et fuisse ante haec nostra tempora homines testabuntur  \pend
\section*{AD TIMOTH. CAP. VI. }
\marginpar{[ p.531 ]}\pstart ltabuntur et illorum scripta, longe diuersa ratione et fide a bonis istis et sanctis patribus, qui in monasteriis Regularium et monachorum olim ecclesiasticam doctrinam promouere conati sunt, ut fuerunt Beda, Alcuinus, Rabanus, Haymo, et similes multi. Indicium sane horribilis irae diuinae in huiusmodi Scholasticos, ob perfidiam illam suam, qua uerbum dei neglectui habere non timuerunt, et hu manis traditionibus pro omni uirili oppugnauerunt simplicem ueritatem, cui prae tulerunt suum sensum, fastum et quaestum.  \pend
\phantomsection
\addcontentsline{toc}{subsection}{\textit{Ex quibus oriuntur inuidiae, contentiones et blasphemiae, suspicione malae, conflictationes hominum mente corruptorum. }}
\subsection*{\textit{Ex quibus oriuntur inuidiae, contentiones et blasphemiae, suspicione malae, conflictationes hominum mente corruptorum. }}\pstart Vt de tot haeresibus taceamus, ante tempore Augustini ecclesiam molestanti bus, et inuidiam episcoporum tempore Arrianorum, de contentionibus et blasphemiis synodorum quarundam. Quid nostris uidimus temporibus? tot simultates inter Vniuersitates, quas uocant, in eadem Vniuersitate, inter uias Antiquorum, Modernorum, Realium, Nominalium, Praedicatorum, Minorum, ob quaestiones et pugnas uerborum, quot et quanta scandala publica, quot imposturas et insanias, usque ad sacrilegia. Tot scriptores in sententias Magistri, quorum nullus non impetit alios omnes, uellicat et improbat. Quin et consentiuntomnes falsos articulos sui Magistri Parisiensis, solenniter improbantes, ut in illis Magister non teneatur. Hodie autem fere omnes concordant mirabili et execrabili consensu aduersus tractantes diligentius sacras literas et Euangelium Christi, proq uiribus conantur scripturam uel trahere ad suos sensus et abusus propugnandos, uel ad labefactandam sacrosanctam authoritatem eius, quam non maiorem conantur asserere, quam suos articulos uel patrum sententias. Sed ut omnes conspirent contra canonicam scripturam, nihil tamen efficient, frustra conabuntur contra fortissimam ueritatem et uerbum domini. Infatuauit hodie sapientia dei magistros illos nostros, et in agnitionem suae uoluntatis transtulit, ut uidemus, paruulos, plebeios, piscatores et agricolas, mulierculas quoque et uirgines, et quotquot non sunt inflati spiritu Aristotelico, et superbia mentis suae, ut honorem et ludicium suum praeponant clarissimae ueritati et fidei. Suspiciones malae) dicuntur, dum alterum alter haereticum iudicat, ac proclamat excommunicatum, de haeresi suspectum, in ira dei, in potestate diaboli, in excommunicatione Papali, ob solas nega tas traditiones humanas, et Scholasticos articulos, quae tamen ex sacris literis pro bari non possunt, et quae nunquam certam apud omnes obtinuerant authorita tem, sed quibus semper fuit reclamatum sacris literis. Superuacaneae conflictationes hominum (ait) mente corruptorum.) Fortassis dicet quispiam, hodie conflictationes huiusmodi augeri, eo quod receditur a placitis magistrorum nostrorum, et ab Vniuersitatum articulis, dum mare et Alpes non sinuntur transcendere determinationes Parisiensium. Sed quis probabit superuacaneas eas esse uel habendas conflictationis causas? cum sint de lege dei discenda et docenda, de euangelio, de fide, peccato, sacramentis, libero arbitrio, deque solis asserendis diuinis scripturis, non de opinionibus philosophorum, et logicis uel metaphisicis. Omnia quoque tantum contra theologastros geri et quaestionarios opinatores, magis quam contra doctores catholicos. Quis dicet superuacaneas conflictationes, si de linguarum proprietate et dialectis disputemus, circa literam interpretandam simpliciter ueteris et noui testamenti, ad inquirendum germanum sensum sacrae scripturae? Cui soli inniti debent dogmata ecclesiastica, et fidei Christianae articuli. Num dicentur esse humanae quaestiones, quae solae uersantur circa literas diuinas,  \pend
\section*{COMMENT. IN I. EPIST. }
\marginpar{[ p.532 ]}\pstart et organa sancti spiritus, in quibus est uita et ueritas? Quis mente corruptos dicet eos, qui ex uerbo dei, cui diligenter et cum fide toti incumbunt, digladiantur tanto zelo contra luxum, quaestum ac fastum ecclesiae satraparum, contra totpublicos, notissimos pestilentissimosque abusus, eorum qui sibi arrogant regimen ec clesiae, quam potius conculcant uerbo et exemplo? Quomodo corrupti sunt memte, qui ante omnia fidem inculcant in deum, spem erigunt, charitatem suadent, opera uere bona, quia a domino mandata et approbata commendant, urgentque ́ fortissime?  \pend
\phantomsection
\addcontentsline{toc}{subsection}{\textit{Et qui ueritate priuati sunt, existimantium quaestum esse pietatem. }}
\subsection*{\textit{Et qui ueritate priuati sunt, existimantium quaestum esse pietatem. }}\pstart Quid illis iustius obringere potuit, quam quod abeis ablata est ueritas, tam indigne tractata, tot studiose dubiis inuoluta, tot opinionibus conspurcata, et omni sua specie et gratia destituta, ut ludibrio haberi coeperit Christus ipse, et ueritas sacra scriptura ? Quis em̃ eam aliquot iam saeculis digne tractauit, quibus lingua rum studium nullum curae fuit, sine quibus nihil agitur in sacris libris ?Sed defue runt iudicio dei libri, praeceptores, imo potius deerant sacrarum rerum admiratores et studiosi, dum interim omnes humana ac uana admirarentur, dum honores, titulos inanes, et diuitias ambiunt, dum fides uera non inuenitur in mun do, iuxta praedictionem Christi, dum nemo audet contradicere magistris nostris, dum ueritatis Apostoli concremantur, dum homo peccati sedem occupat Christi, et humilitatem Christi opprimit maiestas Antichristi, hominumque traditiones praeferuntur diuinis legibus. Et ex tot magistris nostris Parisiensibus uix unus et alter inuenit, qui sacras literas totas tractarint suo more, pius Nicolaus Liranus, Albertus magnus doctus, et Carreñ. Hugo laboriosus. In Magistrum autem Parisiensem commentatores ac quaestionarios supra ducentos nominare possimus, quos extare comperimus nobiles et famatos. Existimantium, inquit, quaestum esse pietatem.) Dei uerus cultus, id estpietas, est, fide, spe, charitate deum colere, nihil antiquius ducere que deo placere, uerbo dei meditando et perficiendo insiste re, die noctuque legem domini meditari, scrutari testimonia del, promissionesque di uinas, et toto corde confidere, proximo bene uelle, benefacere, et ut seipsum dili gere, de seipso diffidere, peccata propria plangere et odire, dei misericordiam im plorare ac sperare, inque Christi mediatoris et saluatoris nostri solis meritis praesumere, et c. Hanc ueram pietatem alii existimauerunt implorare sanctos, colere pa tronos, ornare reliquias eorum, peregrinationes instituere, canere uerba dei, et cantica hominum, in honorem sanctorum, Horas canonicas irregulariter legere, dicere uel orare. Quis autem unquam orauit Genesim (Quis intellexit ex eis prophetas ? Denique sacerdotiorum census augere, Missas multiplicare, emere, uende re, lumina incendere mortuis, imagines colere rebus pretiosis, offerre argentum, aurum, animalia, homines, uinum, frumenta, fructus, extruere basilicas grandes, dedicare altaria, uotaque et ordines instituere infinitos. Quis negabit haec nunc et hactenus a multis et magnis existimata esse ipsissimam religionem et pietatem? Quis nesciat haec esse quaestum ? Nónne his modis deus, imo sancti dei homines, contra uoluntates suas coluntur a Christianis nomine ? Quis hunc cultum ferat subuerti ? Nondum diximus de quaestu confessionum, confessionalium, dispensa tione uotorum, redemptione purgatorii, sacramentorum omnium Simoniacis con tractibus, ieiuniorum indictorum redemptionibus, indulgentiis plenariis, episcoporum casibus. Et quis haecomnia commemoret, quibus adauctus est quaestus in immensum ? Nónne religiosissimus dei cultus est, mendicantibus quatuor ordinum in primis, maxime autem sacerdotibus eleemosynas largiri, cumillis sepe\pend
\section*{AD TIMOTH. CAP. VI. }
\marginpar{[ p.533 ]}\pstart liri, et in eorum habitibus, eo ipso quoque tertiam partem poenarum in purgatorio expiari, confraternitates illorum emere, uel eleemosynis promereri? Certe haec homines existimauerunt, sic edocti ab auaris, acceptissimum dei cultum haberi, peccatorum propitiationem, ampla merita, arram felicitatis aeternae, uiam salutis. Et si quis omnibus istis contradixerit, excommunicatus et haereticus habebitur, dignus morte corporis et animae. Hic appello omnes nostri saeculi homines, pios et impios, clerum et laicos, si non haec sic dicta, credita et facta fuerint generalis ter toto orbe Christiano, ut posterius saeculum uix, nisi extarentlibri, talia credituri sint. Et quid in uniuerso clericorum statu non resipit quaestum? Nihil de hoc idolo abominationum, stante in loco sacro, de hoc infernali quaestu, et satanae cul tu satis atrociter dici potest, pro tanto abusu nostrorum temporum. Sed uerbo dei interficietur propediem hoc idolum auaritiae et mendacii, domino propitiando ecclesiae suae.  \pend
\phantomsection
\addcontentsline{toc}{subsection}{\textit{Est autem quaestus magnus pietas cum sufficientia. }}
\subsection*{\textit{Est autem quaestus magnus pietas cum sufficientia. }}\pstart Si quis omnino quaestuosus esse uelit, et ueras consequi diuitias, is disçat sua sorte esse contentus, praesentis uitae necessariis tantum rebus fruatur, dominum deum suum colat in spiritu et ueritate, diues in fide et charitate, ac uerbo dei. Mun dum cum sua concupiscentia contemnere, congerendis diuitiis studeat non instare. Multa conquisiuit, et abundanter thesaurizauit sibi, qui pietatem obtinuit.  \pend
\phantomsection
\addcontentsline{toc}{subsection}{\textit{Nihil enim intulimus in hunc mundum, hauddubium quia nec auferre quid possumus. }}
\subsection*{\textit{Nihil enim intulimus in hunc mundum, hauddubium quia nec auferre quid possumus. }}\pstart Duuitiae huius mundi, ut aurum, argentum, agri, possessiones in hoc mundo, nos tempore praecesserunt, nondum nobis natis etant, nihil harum rerum intulimus, Animam nostram creauit dominus, corpus nobis aptauit, fauor dei praecessit generationem nostram, a nobis nihil uel nobiscum huic mundo intulimus, omnia quoque post nos relinquemus. Corpus terra mater suscipiet, Spiritus ibit ad deum qui dedit illum. Nihil ergo soliciti simus de terrenis, nec etiam in crastinum, mul to minus quam aues coeli. Qui autem hodie insatiabiliores colligunt diuitias huius mundi, quaestumque sectantur et turpe lucrum, quam illi ad quos loquitur hoc loco dei uerbum, Christus in Apostolo, regula pietatis, ad episcopos et presbyteros eccle siae Christianae? Nihil quidem inferunt, nec exportant, sed interim comportant semper diuitias iniquioribus usibus profuturas. Mundum nunc illi possident, et a mundo possidentur, alieni a Christo. Efferent autem secum ex hoc mundo con scientiam criminum, dei terribile iudicium, diffamiam ad posteros, exemplum pestilentiae, animae damnationem, successorisque subuersionem et scandalum, et cer tum periculum damnosi luxus et fastus.  \pend
\phantomsection
\addcontentsline{toc}{subsection}{\textit{Habentes autem alimenta, et quibus tegamur, his contenti simus. }}
\subsection*{\textit{Habentes autem alimenta, et quibus tegamur, his contenti simus. }}\pstart Quanquam hoc pertineat dictum ad Christianos omnes, et Christi discipulos, omnium tamen id maxime episcopos attinet, ut usum omnium rerum suarum naturae et necessitatis finibus metiantur, ut corpus alatur, non oneretur, pinguae scat, luxurietur, non tractetur tenerius, sed ad commoditatem et parsimoniam alimenta suscipiat. Tegatur necessario amictu, minime pretioso siue superfluo aut curioso. Hic locus necessario tractandus semper nostris praedicatoribus, ante omnia ad hoc nostrum saeculum, cum uidemus tantum luxum in uictu et amictu, ut ditissimus hic nobis mundus minime sufficiat. Quomodo nunc adaucta sunt omnia, augenturque indies, sine discrimine opum uel personarum, ad intolerabilem iacturam patriae, euacuationem totius terrae fomentum omnium uiciorum,  \pend
\section*{COMMENT. IN I. EPIST. }
\marginpar{[ p.534 ]}\pstart oppressionem pauperum, perfidiam popularium, tyrannidem diuitum, fidet Christianae ludibrium, scandalum infidelium, Christi domini nostri blasphemiam. Haec, inquam, iugiter sunt et strenue inculcanda, nemo hanc sedulitatem non probabit, negabit nemo, uerissima atque damnosissima esse et pessima omnes norunt, de testantur ut malum quisq in alio, nemo autem emendat in seipso, praecones eccle siae tacent, arrident, aemulantur, canes muti, quorum animae requirentur, de scan dalo iudicabuntur, quin et domini sacerdotes et pontifices maximam sibi gloriam ducunt contrarie uiuentes huic apostolico praecepto, quum haberi uelint maxi, me apostolici.  \pend
\phantomsection
\addcontentsline{toc}{subsection}{\textit{Nam qui uolunt diuites fieri in hoc saeculo, incidunt in tentationem, et in laqueum diaboli, et desideria multa inutilia et nociua, quae mergunt homines in interitum et perditionem. Radix enim omnium malorum est cupiditas, quam quidam appetentes, errauerunt a fide, et inseruerunt se doloribus multis. }}
\subsection*{\textit{Nam qui uolunt diuites fieri in hoc saeculo, incidunt in tentationem, et in laqueum diaboli, et desideria multa inutilia et nociua, quae mergunt homines in interitum et perditionem. Radix enim omnium malorum est cupiditas, quam quidam appetentes, errauerunt a fide, et inseruerunt se doloribus multis. }}\pstart Qui ditescere desiderant uel student, aut impensius cupiunt diuitias, tentantur grauius. Non loquitur autem hic Apostolus de his qui non magnopere opes uolunt, sed si affluant diuitiae, non cor apponunt, utuntur tanquam non utantur, grati deo de donata affluentia, magis q̃ parta multa solicitudine et cura. Sunt em quibus parum etiam et frigide annitentibus occurrunt diuitiae: sunt qui multis laboribus nihil assequuntur, uane moliuntur, quos fugit mammona. Diuites fieri, cum pauperes fuerint, qui euincere conantur sua fata, qui persequuntur diui tias fugientes, praestaret tales aequanimiter ferre suam sortem, ex necessitate fora mare uirtutem. Non enim omnibus diuitiae prosunt. In hoc saeculo.) Si in futu rum reponere conarentur thesauros, ubi nec erugo nec tinea corrumpit, sancte agerent, tutoque uiuerent ac sine periculis. Si diuitiis spiritualibus, fide, patientia, pace, charitate et sapientia uera ditescerent, Christiani essent. Incidunt, inquit, in tentationes.) Volentes nolentes permissione dei, et diaboli machinatione, contra suam conscientiam, famam, uitam, salutem, contra deum et scripturam aliquan do tales moliri quaedam contingit, ex affectu et desideriolucri turpis. Cum prius quam ditescere conarentur uel inciperent, ea ut mala, perniciosissima abhorruissent. Sic autem fit plura beneficia uel sacerdotia possidentibus, et theologis magistris, si Curtisanorum praxim inciderint, et bonis uiris, si multa proles accidat, quam ad gloriam ampliorem prouehi cogitent huius mundi, uel ad statum suo similem. Hii enim statim incipiunt uehementer tentari per auaritiam. Et quis ne sciat in huiusmodi studiis tentationum molestias et maxima pericula maris, uiarum, latronum, praedonum, infidelium ac mendosorum, principum quoque et coaequalium, ut nuspiam illis tuta fides occurrat ? Quanto magis enim diuitiae affluunt, tanto maior illis conflatur inuidia plurimorum. Figulus enim odit figulum, et mercator comparem. Deficiunt quoque paulatim in timore dei, laxa conscientia. Spernunt primo uerbum domini, deinde odiunt, tandem persequuntur. Nemo enim simul deo seruire potest, et mammonae. Sequitur autem tentationem trepidatio et casus. Qui enim diligit periclitari, facile perit. Qui incaute ambulat, impingit. Mundus plenus est laqueis. A mundi amatoribus facile incautus et auarus inescat et capitur. Cum iustus uix saluatur, et oculatissimus uix euadit, quid patitur desideriis languens, ubi unum peccatum trahit aliud, et inuoluit inexpli cabiliter ? Semper autem adest tentator diabolus, semper quaerens quem deuoret rugiens leo: magis tamen unusquisque a sua concupiscentia trahitur ac ilicitur.  \pend
\section*{AD TIMOTH. CAP. VI. }
\marginpar{[ p.535 ]}\pstart Iuaque uoluntate et deliderns propriis ac terrenis inuoluitur et damnatur. Et de sideria, ait, multa inutilia ac nociua.) Quanto magis auari ditantur, tanto magis desiderant. Facilius est totum deserere, que desideriis habendi nimium imperare. Subsequitur autem habita, inutilis cura ac solitudo, dum omnia timent, semper attoniti, nunquam lae ti, nihil enim illis sufficit, nihil satis fortunatum. Optauerant enim homines illi longe plura, iam fata accusantur, somniis inquietantur, nihil cogitant nisi mammonam deum suum, ut uero deo seruire non possint, uel libeat. Nociua autem desideria sunt, quaecunque contra legem dei et charitatem proximi cogitantur in communem iacturam, dum priuatim nobis consulimus, proximum grauamus. Quae mergunt homines in interitum et perditionem.) Desideria no ciua dum augentur, impellunt ad nimis propense uolendum, etiam contra dei le gem, ad satisfaciendum auaritiae, ad saeculares pertrahuntur tumultus, inuoluun tur litibus, iudiciis ac solicitudinibus, diuitiarum scopulis alliduntur. Quid autem restat aliud, q̃ ut mergantur in profundum auaritiae, et intereant diuitiarum mole obruti, hydropisi desideriorum mundanorum perditi, deum negligant, odiant, uendant et tradantut ludas, quem auaritia infelix reddidit inuidum, sacrilegum, impium, traditorem, et parricidam? Quot hodie uiri religioso statu clari, deumque ́ alioqui timentes, sed sic mammonae auaritia demerguntur, ut proximi amorem ac necessitatem nihili ducant, tantum ut sui census augeantur et possessiones, quantumlibet certi sint quam nihil ualeat omnis uita et conuersatio sua, etiam in suis conscientiis redarguti atque conuicti, omnia sua geri fucata religione, deumque offen di talibus certo sciant, sed auaritia irretiti sopiuntur ? Ideo Christus animarum medicus et saluator, paupertatis amorem docuit, pecuniarum diuitiarumque contenptum: non tantum spiritu, quae necessaria est in primis, uera paupertas omnibus Christianis. Rebus itaque et possessionibus immunes esse uoluit apostolos et apostolicos uiros, ne in diuitiarum admirationem incidentes, ut facile fit, a fide aberrent in uerbo dei praestanda semper. Vnde enim tot abusus, torluxus, tot errores, et scandala in moribus, q̃ ex auaritia clericorum? Qui ubi scripturam sacram sibi ubique aduersari sentiebant, deprauauerunt eandem fucatis interpretationibus, ad suum quaestum, et distinctionibus friuolis succurrerunt suae auaritiae et impie tati. Inde tot regulae, tot dispensationes, tot priuilegia: de quibus iam alii satis clamant, ut audiri iam coeperint. Patet em̃ late, quomodo cupiditas et amor suiipsius fidem synceram et catholicam ac apostolicam, e terra exterminauerint. Verisimi le est et tempore apostolorum statim subintrasse pseudapostolos, qui quaestum magis quam pietatem inquisierint, quandoquidem de his saepe et multum queritur Apostolus. Radix omnium malorum cupiditas est, dum quae sensui applaudent, affatim sectamur, et amorem proprium non moderamur, sed quicquid libet inquirimus, si uacat, etiamsi minus liceat. Vnde et incidunt mancipia cupidinis in multas angustias, solicitudines, dolores ac simulationes, fictamque uitae austeritatem prae se ferunt, et hypocriticam iustitiam coguntur amplecti, qui pecuniam cupiunt, ac ditari uolunt. Sieut autem suauissime uiuunt, qui sua norunt esse sorte con tenti, rerum mundanarum parum soliciti, quibus pauca desiderantibus, pauca ab sunt, modicis contenti: sic qui succumbunt affectibus colligendi pecunias, omnium fiunt miserrimi, plenique perpetuis angoribus, quibus semetipsos inserunt, nemine compellente, ut sine doloribus assiduis mundo huic uiuere non possint, qui Christo placere et acquiescere, meliora suadenti, ac praecipienti non eligunt. Sua culpa bis miseri, cum fide et obedientia Christi possent esse beati.  \pend
\section*{COMMENT. IN I. EPIST. }
\marginpar{[ p.536 ]}
\phantomsection
\addcontentsline{toc}{subsection}{\textit{Tu autem, o homo dei, haec fuge, sectare uero iustitiam, pietatem, fidem, charitatem, patientiam, mansuetudinem. }}
\subsection*{\textit{Tu autem, o homo dei, haec fuge, sectare uero iustitiam, pietatem, fidem, charitatem, patientiam, mansuetudinem. }}\pstart Quae praedicta iam sunt, agunt pro more suo mancipia huius mundi, suiipsius ac saeculi amatores, ditescere cupiunt, quaestum sectantur, ueritati resistunt, falsitatem et errores praedicant, superbiunt sensu, et habitu. Tu uero homo dei, a deo electus, donatus multiplici gratia dei, spiritu dei ornatus, per prophetiam ordina tus, totus deo dicatus, et ecclesiae minister datus, tu haec omnia fuge, quae deo aduersantur, tibi et ecclesiae tuae obsunt, sectare uero iustitiam, tota uita tua iusta sit et uirtuosa, ut sis sine quaerela coram hominibus, nemini iniuriam aliquatenus ir rogando. Soli deo placere cupias, omnem improbitatem caueas in opere et opinione, ut sit tibi saluus honor et authoritas integra in doctrina. Hic, ut uidetur, generaliter intelligenda est iustitia, omnes uirtutes complectens, quarum subinde insigniores recenset. Pietatem) inquit. Virtutem equidem, qua deo placere possis et hominibus, habet enim remunerationem insignem et certam in uita quae nunc agitur, et futura. Haec hic dicitur tractabilitas ad homines, et insignis modestia, cum timore filiali ad deum ut ad patrem. Fidem) dicit totam fiduciam, qua ex deo totus pendeas, ut caetera omnia paruifacias prospera uel aduersa. Quicquid egeris in dei gloriam, certus sis deo placere, deque eius bonitate et gratia sine cunctatione semper bonum sperare. Charitatem.) Ad inseruiendum proximo dum possis, affectu in primis, et per exhibitionem operis in signis, accepta oppor tunitate. Omnibus ut sis dulci, pio synceroque corde coniunctus, paratusque benefacere aliis, ut tibi fieri iuste rationabiliterque uelis. Patientiam) addit, quae piis in hac uita necessaria est, quos manent persecutiones in hoc mundo, qua et animam possidemus, et Christi exemplo maxime commendatam, uiam ad regnum coelorum regiam, omnibus quoque sanctis familiarissimam. Mansuetudinem) quoque , ut peccantibus facile indulgeamus, omnia quae bene fieri possunt, in bonam par tem interpretando: non reclamando, si contra q placeat quidpiam in nos uel dicatur uel fiat. Benefacere etiam nobis malefacientibus, et in bono malum uincere, oppressis pro uiribus condolere et succurrere. Haec episcopo diligenter cogitanda sunt semper, sed longe diligentius exercenda in hoc mundo.  \pend
\phantomsection
\addcontentsline{toc}{subsection}{\textit{Certa bonum certamen fidei, apprehende uitam aeternam, in qua uocatus es, et confessus bonam confessionem coram multis testibus. }}
\subsection*{\textit{Certa bonum certamen fidei, apprehende uitam aeternam, in qua uocatus es, et confessus bonam confessionem coram multis testibus. }}\pstart Docui te in primis ueram in Christum fidem, ut tu quoque eandem doceas om nes, utque in et per eam apprehendatur uita aeterna. Hanc doctrinam ac fidem sem. per confitere. Quandoquidem autem fidei illius hostes multos sustinebis, aduersarios ueritatis, qui per obseruantiam Mosaicae legis religionemque dierum ac ciborum, uestium caeterarumque obseruantiarum carnalium deo placere uolunt, et idipsum docent, ignorantes fidem in Christum. Ideo fortiter et strenue certamen ineas cum illis, atque sustineas. Hic enim uincendi sunt, ad hanc pugnam conficien dam uocatus es. Hanc coram meipso et multis testibus professus es, esse eam apprehensionem et consecutionem uitae aeternae, et consummatae iustitiae, ut in ea sola saluandi sint omnes qui saluantur. In isto certamine si fidem perdideris, simul quoque aeternam uitam amiseris. Certa itaque pro hac fide, apprehende uitam aeternam, ad hanc enim uocatus es. Quaestum proculdubio magnum assequeris, si fidei illius pietatem obtinueris et seruaueris. Foris et intra ecclesiam daemonum insultus experieris, ut tuam fidem labefactent, terreant, anxiumque a fiducia in deum deiiciant. Sed resiste in fide hac unica, apprehende promissam tibi et fidelibus uitam aeternam, ad quam uocatus es, et per fidem obtinebis, si pro ea perse\pend
\section*{AD TIMOTH. CAP. VI. }
\marginpar{[ p.537 ]}\pstart ueranter dimicaueris. Confessus es fidem illam, et me docente credidisti, et docendam suscepisti, dum episcopatum, id est munus praedicandi, et ecclesiam instituendi obiuisti. In hac confessione perseueres atque proficias. Quomodo autem uita aeterna, nuda fide credendo deum esse et ueros articulos fidei, posset apprehendi, cum id etiam diabolus possit et impii ?Sed tora mente deo fidere, fiduciamque bonitatis in deum comprobare per obedientiam praeceptorum, uel de fragilitate admissa humiliter ueniam orare et sperare, quis neget esse uitae aeternae uiam, arram, ianuamque ex qua nemo excluditur? De qua alias.  \pend
\phantomsection
\addcontentsline{toc}{subsection}{\textit{Praecipio tibi coram deo, qui uiuificat omnia, et Christo lesu, qui testimonium reddidit sub Pontio Pilato, bonam confessionem, ut serues praeceptum sine macula, irreprehensibile, usque in aduentum domini nostri lesu Christi. }}
\subsection*{\textit{Praecipio tibi coram deo, qui uiuificat omnia, et Christo lesu, qui testimonium reddidit sub Pontio Pilato, bonam confessionem, ut serues praeceptum sine macula, irreprehensibile, usque in aduentum domini nostri lesu Christi. }}\pstart Deus creator omnipotens et immensus omnia uiuificat, in quo uiuimus, mouemur et sumus: non modo autem nos homines, sed et omnes creaturae a deo crea tore accipiunt perpetuum suum esse, uiuere, moueri, operari, et quicquid possunt. Cui subiecta sunt omnia, cui obedire debent omia, inque gloriam et honorem suum ordinantur. Eidem obedies, seruies, placebis: si quae de fide in eum praedicanda committo, compleueris. Iesus autem Christus dominus, cum ante Pontium Pilatum iudicandus astaret, testimonium reddidit ueritati, nec mortis timore cuncta bundus, contestatus est ueritatem se oportere testari: quam et tu, si, ut obligaris, praedicaueris, attestatusque fueris praeceptum illud meum de fide et uita aeterna praedicanda, diligenter cu stodies, et fidelibus, ut omnium maxime necessariam inculcabis, per quam et iustificentur, et uitam sine dubio assequantur aeternam, Sine macula, ait, uolens ut fidei praedicationem, et caeterorum quae imposuit et commisit discipulo, sine macula seruet, dum quaestum non admiscet, dum sua non quaerere auditur, dum non tam laudari quam proficere cupit. Maculam autem inurunt non solum praedicatori, sed etiam doctrinae, quaestus, ambitio et insynceria tas, dum proditur uel notatur. Refertur etiam ad ipsum uerbi ministrum, ut non seandaliz et opere, quos instruit doctrina. Irreprehensibile) quoque uerbum fidei dicitur, cum sacrarum scripturarum authoritate declaratur, ostenditur, roboratur, ne sensu proprio docere uideatur, quod reprehensionem non euaderet, cum omnis cogitatio et sensus humanus mendacium operentur. Traditiones quoque  humanas uidemus sine reprehensione doceri non posse, sunt enim mutabilitati obnoxiae. Lex autem domini irreprehensibilis est, si bene percipiatur, et mutari non potest. Secundum quam tu quoque et omnis episcopus praedicans ac docens, irreprehensibilis eris. Vsque in aduentum, etc.) Quousque Christi fides gratia et gloria toto orbe clarescat, et agnoscat totus mundus eum esse dominum et redenptorem, cui deus pater omnia dedit in manus, subiecit omnia, ut saluator sitomnium, per Christum in deum credentium et sperantium. Qui solus colendus est cum filio et spiritu sancto, et inuocandus, per quem deus omnia restituit, et recreat orbem. Per eum et in eo necesse est conciliari deo omnem hominem in hunc mundum uenientem. Haec, inquam, agnoscere, et de Christo confiteri, utque redem ptorem glorificare et colere, ut regi regum et domino dominantium seruire, uer bosuo credere, in eum sperare, omne bonum de eo consequi, longe nobilior est susceptio Christi aduenientis cum maiestate, feliciorque cultus, quam si rex ac princeps Christus corporaliter adueniret in mundum, denuo gloriosus, et conspicuus omnibus, sicut uenerat in carne passurus ac moriturus. Donec aliquando in gloria  \pend
\section*{COMMENT. IN I. EPIST. }
\marginpar{[ p.538 ]}\pstart maiestatis suae orbem uniuersum, uiuos et mortuos, manifeste ac iustissime iudicaturus adueniat in fine mundi.  \pend
\phantomsection
\addcontentsline{toc}{subsection}{\textit{Quem suis temporibus ostendet beatus et solus princeps, rex regnantium, et dominus dominantium, qui solus habet immortalitatem. }}
\subsection*{\textit{Quem suis temporibus ostendet beatus et solus princeps, rex regnantium, et dominus dominantium, qui solus habet immortalitatem. }}\pstart Ostendet dominus Christi domini apparitionem et gloriam toti orbi, quum euangelium suum ubique praedicabitur, acceptabitur, intelligetur, credetur et seruabitur. Haec erit nobilis Christi apparitio, cum non iam in carne corruptibili, sed in uirtute et gloria spiritus sui subiiciet sibiomnia. Id autem fiet gratia et dono dei patris, tempore destinato et praefinito, omnibus ignorato modo ineffabili, dum subito sicut fulgur ab Oriente emicat et apparet in Occidentem, qua illustratione uultus sui orbem iudicabit. Beatus.) Vnico hoc uerbo ineffabilem maiestatem ueluti balbutiendo exprimit. Beatus est enim, qui possidet omne bonum, is solus deus est. Sicut et supra dixerat, Beati dei, quo scilicet maior et felicior sanctiorque cogitari non potest, nec nominari. Et solus potens.) Annititur humano loquendi modo uel paululum aliquid clarius et specialius de deo eloqui. Dicitque ́, Solum potentem: quod plus est q̃ omnipotens, et magis incomparabiliter, quo tamen re uera maius explicari nihil potest. Vt non possit tantum omnia, sed et quicquid uspiam est potestatis, uel potentiae, id sit dei solius. Nulla enim creatura aliquid potest, nisi in et per deum. Omnia enimoperatur in omnibus, nihil hic distinguitur de potentia dei ordinata uel inordinata, absoluta uel conditionata, non de malo uel bono. Omnia enim ipse solus potest, solus ipse operatur, nihil ma lum quod facit. Vidit enim deus omnia quae fecit, non solum ea quae creauit, et erant ualde bona. Ipse itaque solus omnia porest, nos possumus nihil, nisi in eo. Sic angeli quoque et coeli, sicomnia quae sunt, uiuunt, mouent, et operantur in deo, ut nihil potentiae nobis adscribamus. Peccare uero potentia non est, sed defectus in opere bono, et impotentiae indicium, quae de deo sicutomne quod imperfectionem connotat, dici uere non possunt. Solus enim beatus est, solus potens, solus do minus uniuersorum, a quo omnis creaturarum uirtus et actio proficiscitur.  \pend\pstart Rex regum.) Ipse rex solus est, omnia gubernans, disponens et prouidens. Nihil enim in mundo quicunque domini possunt, nisi per deum, nec uelle bonum nec fa cere. Non solum cor regis, sed quicquid uiuit et habet, potestque et agit, in manu domini est, et quo uoluerit, illud uertet. Vnicus iste rex regum, deus solus potens et faciens omnia, nobis orandus est, ut propitius, clemens et misericors dirigat actus nostros in beneplacito suo, dispensationemque dei in hoc mundo, humiliter et patienter modestissimeque amplectamur, quicquid tandem illud fuerit. Interim tamen uoluntatem illius ex suo uerbo discentes, omni studio et uiribus annitamur, ut quae praecepit dominus, haec fiant, tam a nobis quam ab aliis. Hoc enim nos conari debere, scimus deum ipsum uelle et praecepisse, ut quicquid in manus dederit agendum nobis, hoc instanter et fideliter operemur, scientes uoluntatem domini dei nostri. Quod si prouidentia sua moderante, secus cesserit ac uolebamus, minusque successerit bonum, quod etiam domino inspirante instituimus, certi tamen simus, quod haec nos deus facere uelle uoluerit et tentare, etiamsi non perficiendum concesserit. Si peccatum est quod accidit, pro culpa admissa oremus, eiusque ́ iudicium deprecemur. Si melius processit, domino acceptum referamus. Sic igitur rex regum, et dominus dominantium agnoscetur a piis, digneque coletur. Si nos et nostros actus iuxta suam uoluntatem pro uiribus instituerimus, quas uires, et quod uelle ac perficere suae tantum gratiae adscribamus, nihil nobis, qui sineeo  \pend
\section*{AD TIMOTH. CAP. VI. }
\marginpar{[ p.539 ]}\pstart nihil possumus neque sumus. Solus habet immortalitatem.) Solus est fons uitae, et ipsa uita, qui non potest non esse, a semetipso habet esse, et alia omnia suum esse accipiunt ab eo, et participatione sui subsistunt. Nunquam desinet, ut nec incepit. Aeternus est, deficere non potens nec destrui, contrarium enim non habet; nec est compositus, sed simplicissimum bonum, sine principio et fine existens. Immor talitatem solus habet, quam donare possit cui uoluerit, a quo solo eam sperare, consequi et habere possumus, et a nullo alio. Cuius regnum et ministri angeli, animae et spiritus immortales sunt per dei gratiam.  \pend
\phantomsection
\addcontentsline{toc}{subsection}{\textit{Qui lucem habitat inaccessibilem. }}
\subsection*{\textit{Qui lucem habitat inaccessibilem. }}\pstart Lux est deus, et in eo non sunt ullae tenebrae, nihil enim ignorat, omnia nouit quae sunt, fuerunt, erunt, et esse possunt, omnia praesciuit et praeordinauit. Cumque ́ re uera sit maxime cognoscibilis, quia perfectissimum uerum, imo ipsissima ueritas, tamen a se solo cognoscitur perfecte, et ab unigenito suo. A nulla autem crea tura cognosci potest, nisi quando et quantum sese dignatur ostendere. Voluntarium est enim speculum, si uult uidetur: si non uult, non uidetur. Incomprehensi bilis est omni creaturae, etiam angelicae. Accedere eum nemo potest, sed dignatur admittere uel attrahere amicos et filios, quos dignos habuerit. Lucem habitat, quia sibi soli sufficit, suapte luce et bonitate beatus, nullius indiget, accipere et consequi nihil potest, in omnia se diffundit, et bonitate sua omnia replet, quam omnes participant creaturae. Sicut lux splendorem diffundit, et a remotis uidetur, ac per reflectiones: sic deus in creaturis cognoscitur quantum uult, quandiu, et quibus.  \pend
\phantomsection
\addcontentsline{toc}{subsection}{\textit{Quem nullus hominum uidit, sed nec uidere potest. }}
\subsection*{\textit{Quem nullus hominum uidit, sed nec uidere potest. }}\pstart Angelos uiderunt patriarchae et prophetae, figurasque a deo formatas, per quas dei mandato, angeli uoluntatem summi imperatoris ad homines multiformiter retulerunt. Ipse autem in sua essentia a nemine uisus est unquam, nec uideri potest, incorporeus et purissimus spiritus, imo spiritualior omni spiritui. Verbum quidem patris et filius in diuinis, in fine temporum humanandus, carnemque ex impolluto uirginis utero et sanctificato assumpturus, mundum liberaturus, deique ́ amorem maxime commendaturus humano generi. Is in ueteri testamento patribus san ctis loquutus est, apparuit, promisit, terruit, docuit, monuit, omniaque peregit quae deus cum patribus egisse describitur. Creasse seilicet mundum, et loqutum fuisse ad Adam, Cain, Noë, Abraham, Isaac, Iacob, Mosen, Iephthe, Gedeonem, Samuelem, Dauid, et prophetas. His siquidem omnibus quaecunque deus locutus scribitur, et a sanctis patribus, non sine scripturarum authoritate egisse dicitur, dei filium, sapientiam, sermonem, et uerbum dei egisse intelligatur, sicut et scripturis certissimis comprobatur. Sed nec, inquit, uidere potest.) Non uidere potest ho mo nisi lucidum et coloratum, et obnoxium accidentibus, sui speciem a se diffun dere ualens, imo sola ipsa accidentia, nullam autem substantiam, quae nullo humano sensu cognosci potest a nobis, sed tantum per accidentia. Deus autem pura et spiritualissima ac supersubstantialis substantia est: ad cuius intuitum oculus noster instar noctuae, ad manifestissima naturae occaecatur. Spiritu autem et amore deus percipitur ab homine, dum eum attraxerit, seseque dignabitur inclinare, et fa uore suo complecti.  \pend
\phantomsection
\addcontentsline{toc}{subsection}{\textit{Cui est honor et imperium sempiternum. Amen. }}
\subsection*{\textit{Cui est honor et imperium sempiternum. Amen. }}\pstart Honor inter homines professio est dignitatis et bonitatis in honorato. Cum itaque solus deus bonus sit, a quo omnia bona, ideo soli deo honor debetur et glo\pend
\section*{COMMENT. IN I. EPIST. }
\marginpar{[ p.540 ]}\pstart ria: peccatque uehementer, qui uel honorari cupit de donis non suis, sed dei, que a deo suscepit, cui ea non refert accepta, quasi non receperit: uel, qui honorem ei defert cui non debetur, fraudando eum a quo bonum omne est quod honorat. Talis est apud multos sanctorum honor, qui a se nihil nisi peccata habent, sola autem des gratia sunt id quod sunt, quem et solum honorari uoluissent, quia sancti sunt. Imperium) quoque omnium, dei solius est, qui imperat uniuersis, qui facit et potest quicquid uult in coelo, terra, et in omnibus abyssis. Cuius nutui subiecta sunt omnia, ad cu ius maiestatem Caesaris Augusti imperium uix est ceu eleuatio capitis minimi uermi culi, nisi quantum suo imperatori deo subiectum se Caesar agnoscit, et obsequitur studiose, ad domini dei uoluntatem de hominibus promouendo. Nemo quamlibet in hoc mundo potens sibi contra deum quicquam uendicet, cogitet potius deum omnibus praesidentem et imperantem, cuius minister sit, ut sit fidelis fami liae sibi commissae a domino, quem in horas expectet iustissimum iudicem. Amen.) Veritas est, ut credenda discamus quae sunt praedicta, et infallibilis ueritatis uerba ac sententiae. Quae et cum reuerentia ac timore dei omnia audienda sint, ut ue re profiteamur nos credere de deo summum honorem summumque imperium. Nos uero digne subiectos et obsequentes exhibeamus summo bono.  \pend\pstart Diuitibus huius saeculi praecipe, non sublime sapere, neque sperare in incertum diuitiarum.  \pend\pstart Collecta pro pauperibus, qui semper nobiscum sunt, nunquam intermittenda docetur his uerbis. Vult autem praecipue diuitibus praecipi eleemosynae largi tionem, quae tamen parum ualet, nisi ab humili profecta, et non a iactabundo uel superbo. Ideo primum hortatur diuites ad humilitatem, ut meminerint nihil se habere a semetipsis, dona dei esse quae assequuti sunt, ideo non contemnant pauperes, non moribus, uestitu et sensu superbiant. Agnoscant plura se debere deo, qui plurima acceperunt, dum honorantur, quiete uiuunt, nullius indigent, domini recordentur, in cuius manu sunt omnia, alioqui fastus solet diuites occupare: cum tamen quae difficile plaerunque obtinentur, facile etiam amittunt: caduca sunt enim bona diuitum, et incertae diuitiae, ideo nulla in eas spes est reponenda. Imminent fures et fata, nihil fixum et certum in illis. Spes uero certa est expectatio futurorum, non incerta sit caducorum. Exempla habent historiae sacrae atque prophanae fugatium diuitiarum. Gentes sperant temporalia, fideles spiritualia et aeterna.  \pend
\phantomsection
\addcontentsline{toc}{subsection}{\textit{Sed in deo uiuo, qui praestat nobis omnia abunde ad fruendum. }}
\subsection*{\textit{Sed in deo uiuo, qui praestat nobis omnia abunde ad fruendum. }}\pstart Semper nostra curat uiuens deus, et semper uigilans, nihil ignorans, cuius fi lii sumus, negligi a tanto patre non possumus, si eidem confidamus. Obtinetur illius bonitas et beneficium fide, conseruatur spe, perficitur charitate. Nemo confidat in mortuos. Sanctos non nobis deus praefecit in hac uita constitutis, de eis nobis deus nihil commisit. Solus deus anchora sit nostrae spei per Christum. Non in opera bona quaecunque nobis confidendum, sed deeis magis timendum, si ad iu stitiae diuinae lancem appendi contigerit. Deus autem pater est gratiarum et con solationis, qui clementiam et gloriam pollicitus est, et potest uultque praestare, si firmam illi patri nostro filii habuerimus fiduciam. Praestat autem sua dona deus abun de, diuitias, honorem et sanitatem, ut frui quidem eis liceat aliquandiu, cum timore et gratiarumactione: nec id diutius dabitur, quam domino concedente. Ideo senper ambiendus dei fauor et gratia operibus misericordiae, quae exigit a diuitibus. Quales enim proximis suis fuerint, talem sibi dominum experientur. Omnia autem praestat, quod nihil nostra prudentia, cura laboreque obtinemus: nihil enimpro  \pend
\phantomsection
\addcontentsline{toc}{subsection}{\textit{}}
\subsection*{\textit{}}
\section*{AD TIMOTH. CAPI VI. }
\marginpar{[ p.541 ]}\pstart ficimus, nisi eo largiente. Timeant semper diuites, ne domino irato diuinias suas possideant, ne non iuste obtinuerint, ut q̃ iuste utantur, iuste dispensent, solicite semper cogitent de graui ab illis ratione exigenda. Difficile est uere Christianum esse diuitem, nisi affatim erogando, quae deus abunde illis suppeditat, ut non nisi donorum dei sibi dispensator uideatur. Dum uero ad fruendum dicitur, improprie frui sumitur prouti. Non enim nisi propter nostram necessitatem et proximi, desiderandae uel possidendae sunt diuitiae, non propter se. Solo enim deo fruendum, quia solus est summum bonum, propter seipsum diligendum. Non ergo necessario obsequendum semper regulis sophistarum, qui hic haeresim Paulo impinge, rent dicenti, diuitiis fruendum esse. Libere et clare eloquamur quae uera sunt, ma gisque quid, quam quibus exacte uerbis explicemus curandum erit. Non uultut abunde fruamur, sed abunde praestat ea quibus modeste fruamur, hoc est parce, non ad luxum uel fastum.  \pend
\phantomsection
\addcontentsline{toc}{subsection}{\textit{Bene agere, diuites fieri in bonis operibus, facile communicare. }}
\subsection*{\textit{Bene agere, diuites fieri in bonis operibus, facile communicare. }}\pstart Docuit diuites primo humilitatem, secundo abrenunciationem, tertio spem, quarto fruitionem. Nunc quinto docet bene agere, non magnopere ea sectari quae deus non praecepit, ecclesias scilicet, monasteria, anniuersaria, praebendas, aras, imagines erigere: sed quae deus praecepit, Omni petenti tribuere, pauperes pasce re, potare, redimere, iuuare, consolari, mutuum praestare, nihil inde sperare. Non audire Missas, minime intelligendas et cantiones, sed mendicos clamantes. Non preculas adnumerare, sed precantes exaudire. Non testamentum ordinare dum moreris, sed sano sensu et corpore erogare superflua indigentibus, cum adhuc ui uis et tua possides, de eis domino redditurus rationem, ut dispensator bonorum domini dei tui fidelis et sedulus in commissis. Addit autem, Diuites fieri in bonis operibus.) Sexto monet, ut ditescant, non bonis corporalibus, sed spiritualibus. Non quidem ut de bonis suis praesumere incipiant, et sibi placeant, uel sensu abundent, sed ut bona opera a deo mandata, cum dei adiutorio, obediendo, perficiant, seque ut inutiles seruos credant, ac profiteantur ingenue. Quod facere debent, hoc faciant, non spe mercedis, sed instinctu charitatis, et fiduciae in deum. Nominat autem diuites eos, qui dediti sunt operibus misericordiae, qui soli diuitias habere censentur coram deo, cum scilicet legitime ac religiose administrant iuxta commissionem dei, cuius procuratores constituuntur. Vnde ad mentem et modum loquendi apostoli Pauli, Diuites fieri in operibus bonis, aliud non est, quam opera pietatis de diuitiarum abundantia affatim exercere, opera uere bona ac religiosa multiplicare, iuxta facultatem domini illis donatam abunde prae aliis. Ne contingat eos egere maxime, cum minime uelint, cum diuite epulone, cuigut ta aquae negata fuit, diuiti rebus olim, sed egeno gratia dei et bonis operibus. Bo na opera quae uere dicantur et sint, hactenus ignoratum fuit a multis: nunc disuaderi incipiunt non uere bona opera, sed quae falso putabantur bona: suadentur autem a doctis et piis hoc nostro tempore et docentur ea bona opera, quae uere bona sunt, quia a deo praecepta, laudata et muneranda. Diuites autem erimus non in hoc saeculo, sed in coelo. Illuc enim thesaurizandum, quia hic semper pauperes erimus, plura debentes, que praestare possimus, semper peccatores, semper orantes pro peccatis remittendis, de operibus bonis nihil praesumentes. De mise ricordia dei et fide nostra tantum gloriemur, non in nobis, qui nec fidem praestamus ex nobis nec opera, sed domino largiente utraque suscipimus. In eo autem quod sequitur, Facile communicare.) Adhuc hortari docentur diuites ab episcopo, ut non difficile adigantur ad dandum precantibus, non expectare precantem,  \pend
\section*{COMMENT. IN I. EPIST. }
\marginpar{[ p.542 ]}\pstart sed inquirere egentem, non tam suscipere quam cogere hospitem pauperem. Omnia communia putare, dei dona esse non propria, ad dei beneplacitum non grauatim dispensanda, cum sua non sint, sed eius qui magis indiget.  \pend
\phantomsection
\addcontentsline{toc}{subsection}{\textit{Thesaurizare sibi fundamentum bonum in futurum, ut apprehendant uitam aeternam. }}
\subsection*{\textit{Thesaurizare sibi fundamentum bonum in futurum, ut apprehendant uitam aeternam. }}\pstart Docet in coelum reponere diuitias tuto loco, et in aeternum usui nostro profuturas, alioqui et hic eas perdemus et alibi, perinde ac contigit diuiti in Euangello, repo nenti bona multa in multos annos. Quicquid enim non in deum retulerimus, pro pterque eundem egerimus, totum perdemus. Sic fiet omnibus bonis operibus nostris, nostri causa, non dei exactis. Sic quoque datae eleemosynae, non propter deum pure, perfecte, mercede carebunt optata. Thesaurus uero noster cumulabitur, si nullum ad cumulandum nobis oculum applicemus, sed deo potius reponamus omnia quam nobis: de suo praemiabimur, non de nostro: suis fruemur bonis, non nostris. Patitur nos deus improprie loqui, dum pie intelligamus et uelimus: alloqui qui omia debet alteri, sibiipsi proprie thesaurizare nil potest. Et apud deum nostrum omnes peccatores sumus, gratia quoque sola, non ex proprie dictis meritis, saluabimur. Sua in nobis opera deus coronat. Fundamentum bonum nominat.) Qui enim thesaurum temporarium et incertum bene collocant, et in sinum pauperum reponunt, imo Christo commendant cui donant, habebunt in futuro stabiles diuitias, iucundas, aeternas. Dant enim temporalia, et capiunt aeterna, uitam beatam impari compensione. Cum nihil homo dederit homini de suo, quicquid dederit: denique uoluntatem dandi non nisi a deo habuit, multo minus dona et opera. Veruntamen patitur se deus suis propriis bonis absolui, et reddi quae sibi debentur, donatque unde soluamus: unum quoque accipit pro mille, gratum a nobis suscipit, quod gratis concesserat et commodauerat. His uerbis, quasi adiectitiis, apostolus Paulus causam pauperum egit, utque admonerentur diuites beneficentiae, et pau peres insuper uerecundia liberentur, qui ut uellent non possunt erogare eleemo synas. Non omnibus praecipit de diuitiis expendendis, sed tantum diuitibus, qui citra dispendium proprii uictus et suorum, donare possunt de proprio. Huiusmodi collectarum consuetudo et oblationum in synaxibus reformanda est omnino, et non auferenda, cum sit omnium maxime necessaria pro egentibus. Sic tamen ut collecta hoc modo, ac reposita pro pauperibus, non aliter q̃ pro pauperibus pro fundantur, non coaugmententur in census, ut fit nostro saeculo in hospitalibus, quorum curatores continuo diuitias augere student: non tam pro pauperum usu et sustentatione, quam diuitum desperantium de uictu, uel otio inerti sese dedentium, non sine perfidia annonam huiusmodi ementium pariter atque uendentium.  \pend
\phantomsection
\addcontentsline{toc}{subsection}{\textit{O Timothee, depositum custodi. }}
\subsection*{\textit{O Timothee, depositum custodi. }}\pstart Deposuerat penes eum Paulus, imo spiritus sanctus, gratiam prophetiae, sanaeque doctrinae: nam ab adolescentia sacris literis operam impenderat, et donum uerbi uberius acceperat, pro ecclesiae profectu depositum. Hoc unum nunc iube tur sic custodire, ut non perdatur, non deficiat. Id autem est liberaliter distribuere, sine intermissione docere, creditum talentum non in terram defodere, sed ad usuram domino augmentandam exponere. Tunc autem auarissime custoditur, dum liberalissime profunditur. Depositum autem dicit, quia mera dei gratia, et sine meritis accepit, erogandum in tempore, familiae suae, illapsu et munere diuint spiritus, quasi aliunde donatum. Ostendit autem affectum dicens, O Timothee, excitat, hortatur et admonet, ceu omnino magnum aliquid necessario et iugiter  \pend
\section*{AD TIMOTH. CAP. VI. }
\marginpar{[ p.543 ]}\pstart cogitandum propositurus. Maxime necessarium officium episcopi est docere, orare, prospicere animabus saluandis. Non quidem episcoporum nostrorum temporum, sed uerorum, quorum nunc sunt paucissimi: deposita enim alia habent, pecunias nempe, et equos, et caetera.  \pend
\phantomsection
\addcontentsline{toc}{subsection}{\textit{Deuitans prophanas uocum nouitates. }}
\subsection*{\textit{Deuitans prophanas uocum nouitates. }}\pstart Vult discipulum suum Paulus non alia docere ecclesiam, que ̃ quae didicerat a magistro: non inducere uel confingere inconsueta uerba. Praesensit enim iam Apellis et Valentinianorum aliorumque figmenta, de Eonibus quatuor, octo, et tri ginta, singulas singulis nominibus nominando, quae iam apostolorum tempore a ludaeis et Graecis superstitiosioribus fingebantur. Ex Cabalistarum, ut nunc uocant, et philosophorum maxime Pythagoricorum dictis se recepisse et edoctos. iactabant. De quibus Irenaeus. Quibus nuper successerunt Parisienses theologastri, portentosis nominibus, quae nec Latini nec Graeci intelligunt, nisi de eis diu edoceantur, ad quaestum uberiorem. Praemonuit Apostolus, uitandas huiusmodi uocum nouitates, quas in dei uerbis sanctisque scripturis non recepimus. Id si seruatum fuisset, non tot differentiae doctrinarum et schismata audissemus, De ho musion, transubstantiatione, circumincessione, respectibus diuinis interioribus et exterioribus, essentialium et notionalium perfectionum, et reliqua huius generis loquutionibus, praeter scripturae authoritatem introductis. De quibus sicut et de dependentiis formalitatibus et quidditatibus, obmissis utilibus et necessariis tantopere digladiati sunt doctores illi neothericae ecclesiae, ut magis disputando Scholastice ueritatem perdiderint quam repererint, quin fidem et spiritum cum Christiana uera sanctimonia amiserint.  \pend
\phantomsection
\addcontentsline{toc}{subsection}{\textit{Et oppositiones falsi nominis scientiae. }}
\subsection*{\textit{Et oppositiones falsi nominis scientiae. }}\pstart Appellatur quidem ab illis superstitiosis doctoribus istarum rerum cognitig Scientia, sed indigno prorsus et falso uocabulo, cum sit re uera stultitia, error, opi nio, haeresis, superbiam docens et quaestum, minime uero pietatem. Oppositiones enim sunt apud illos, non doctrinae: subuersiones discentium, non ueritates. Videmus siquidem tales non aliud scripsisse quam oppositiones, argumenta, pro et con tra, pro opposito uel contra oppositum, ut loquuntur. Vide primum inter mendi cantes Albertum, quem Magnum uocant, in Magistrum Sententiarum, uel Summa eius, inuenies meras oppositiones, et praeter argumenta pro et contra, nihil, Ouem caeteri Parisienses omnes sequuti sunt, ueluti in re eximie illis pulchra uisa. Eos nimirum nominasse hic Paulus uidebitur, oppositionum falsi nominis sci£ entiae authores. Qui et semetipsos et scientiam suam caligine oppositionum inuoluunt, et sic improbant uicissim, ut syncera apud eosdem ueritas nulla relicta sit. fides quoque et uera pietas abolita fuerit, quam eorum nullus pro rei dignitate tractauerit, sed in profundum silentium sepulta per eosdem est uera fides, unde fide les nos nominari conueniebat, ut fere in totum expirarint et sancta opera ueraque ́ religio Christiana, donec explosa eorum schola, et falsi nominis scientia per dei gratiam, sacrae literae, Christi spiritus, apostolica fides et doctrina denuo nunc emergunt, ad lectionem antiquissimorum doctorum sanctorum, quos illi quoque Originales uocant. Quorum nobis deus doctrinas et libros Impressoriae artis be neficio communiores reddidit, et passim exposuit ruminandos. Interim Scholastici illi male feriati theologastri, cum suo sacrilego Aristotele stulti deprehendun tur, et uerae fidei perniciosi pharmacopolae. Vniuersitatum quoque studia, ut nunc sunt, inania deprehenduntur, et ecclesiasticae iuuentutis prostibula, necessario  \pend
\section*{COMMENT. IN I. EPIST. }
\marginpar{[ p.544 ]}\pstart in melius reformanda.  \pend
\phantomsection
\addcontentsline{toc}{subsection}{\textit{Quam quidam promittentes, circa fidem exciderunt. }}
\subsection*{\textit{Quam quidam promittentes, circa fidem exciderunt. }}\pstart Promittentes dicit, non exhibentes, eam quam promittunt sclentlam, non enim praestare possunt quod spondent. Contradictiones euadere non docent, sem per autem discipulus magistro reclamat, diuersum sentit, docet et scribit. Disceptare tantum norunt, omnia quaestionibus inuoluere, et super lana caprina digladiari. In uno tantum conuenientes, quod se ueram fidem nescisse, et opera ue rae pietatis non curasse fateri cogantur: sacra Biblia nunquam uidisse, legisse, didicisse, intellexisse, nec docere potuisse, cum tamen eandem canonicam scriptu, ram, et uerae fidei fundamentum negare non ausint, sed, ut ipsi loquuntur, adaequatum obiectum fidei Christianae scripserint omnes. Digna itaque sunt animaduersione puniti, ut dum sibi cisternas ueteres fodiunt, non ualentes continere aquas, salutarium et uiuarum aquarum fontem perdiderunt: dum humana sectantur, diuina amittunt: dum ratione animali suas doctrinas ut idolum colunt, fidem uerbi dei deseruerunt: dumque se doctores profitentur Aristotelis, amplissimis titu lis delectati, Christi domini et apostolorum discipuli esse desierunt, ac utriusque fructum et eruditionem perdiderunt. Omnem fidel iacturam et morum, cum incre mentis uanitatum ac superstitionum istis theologistis debemus. Exciderunt, inquit Apostolus, circa fidem, aberrarunt a scopo totius scripturae, quae docet fidem in deum, prouidentiam eius circa omnia, promissionesque diuinas in Christo. Ipsi uero errantes in foribus, ad opera dilapsi sunt: non quidem illa quae sacris literis commendantur, sed quae praeter uerbum dei et contra, auaris placuerunt: unde et traditionum humanarum ingluuies, abusuum abyssus, et mare superstitionum ac errorum in fide et moribus prodierunt.  \pend
\phantomsection
\addcontentsline{toc}{subsection}{\textit{Gratia tecum. Amen. }}
\subsection*{\textit{Gratia tecum. Amen. }}\pstart Haec manu Pauli adscripta dicuntur, cum alia per dictatum describi fecerit, iuxta eius morem. Optima autem est imprecatio. Quid enim molestiarum non tolerabimus, si dei assit gratia, in quam speramus? Si deus pro nobis, quis contra nos ? Facessat omnis humana gratia et potestas, dum nobis adest diuina.  \pend
\phantomsection
\addcontentsline{toc}{subsection}{\textit{FINIS. }}
\subsection*{\textit{FINIS. }}
\endnumbering
\end{pages}
\end{document}
        